% Généralités sur les moteurs à combustion interne — PCT 3ème — Cours signé OnBuch+
% Programme officiel MINESEC. Compiler avec : tectonic N18-generalites-moteurs-combustion-interne.tex
\newcommand{\DOCMATIERE}{PCT}
\newcommand{\DOCNIVEAU}{3ème}
\newcommand{\DOCMODULE}{Module 4 : Technologie et mécanique}
\newcommand{\DOCLECON}{N18}
\newcommand{\DOCTITRE}{Généralités sur les moteurs à combustion interne}
% OnBuch+ — préambule « cours signés » ENRICHI (compilé avec tectonic / XeLaTeX)
% Système visuel « Pop » d'OnBuch+ : fond crème (jamais blanc), bordures encre
% épaisses, ombres « dures » décalées, titres Archivo Black en capitales.
% Version MATHÉMATIQUES : graphes (pgfplots/tikz), boîtes pédagogiques
% (définition, propriété, méthode, attention, le savais-tu, exercice résolu,
% exercice, TP, bilan), page de garde et signature OnBuch+ ancrée en bas.
%
% Macros attendues dans main.tex AVANT \input{preamble} :
%   \DOCMATIERE  \DOCNIVEAU  \DOCMODULE  \DOCLECON (numéro)  \DOCTITRE
\documentclass[11pt]{article}

\usepackage{fontspec}
\usepackage[french]{babel}
\usepackage[a4paper,margin=2cm,top=3.4cm,bottom=2.8cm,headheight=1.4cm,footskip=1.3cm]{geometry}
\usepackage{amsmath,amssymb,amsfonts}
\usepackage{mathtools} % \xmapsto, \coloneqq... souvent utilisés par les modèles
\usepackage{cancel} % simplifications barrées (bilans de forces)
% \abs{z} (module, valeur absolue) et \norm{u} (norme), utilisés par les modèles
\providecommand{\abs}{}\let\abs\relax\DeclarePairedDelimiter{\abs}{\lvert}{\rvert}
\providecommand{\norm}{}\let\norm\relax\DeclarePairedDelimiter{\norm}{\lVert}{\rVert}
\usepackage[table]{xcolor} % [table] : fournit aussi \rowcolors (alternance de lignes)
\usepackage{pagecolor}
\usepackage{tikz}
\usetikzlibrary{arrows.meta,positioning,calc,shapes.geometric,shapes.misc,shapes.arrows,decorations.pathmorphing,decorations.pathreplacing,decorations.markings,patterns,shadows,mindmap,trees,fit,backgrounds,matrix,intersections,angles,quotes}
\usepackage{pgfplots}
\usepackage[version=4]{mhchem} % équations nucléaires \ce{^{235}_{92}U}
\usepackage[european,siunitx]{circuitikz} % schémas électriques
\pgfplotsset{compat=1.18}
\usepgfplotslibrary{fillbetween,groupplots} % aires entre courbes (calcul intégral)
\usepackage{enumitem}
\usepackage{fancyhdr}
% [most] charge toutes les bibliothèques tcolorbox (listings, minted, raster,
% documentation...) et gonfle inutilement la mémoire TeX sur un document long
% (~300 boîtes) : on ne charge que ce qui est réellement utilisé.
\usepackage[skins,breakable]{tcolorbox}
\usepackage{graphicx}
\usepackage{colortbl}
\usepackage{booktabs}
\usepackage{tabularx}
\usepackage{multirow}
\usepackage{diagbox} % case d'en-tête diagonale des tableaux à double entrée
% Colonnes extensibles centrée / gauche / droite (C, L, R), souvent utilisées par les modèles
\newcolumntype{C}{>{\centering\arraybackslash}X}
\newcolumntype{L}{>{\raggedright\arraybackslash}X}
\newcolumntype{R}{>{\raggedleft\arraybackslash}X}
\usepackage{array}
\usepackage{multicol}
\usepackage{siunitx}
% parse-numbers=false : accepte aussi \SI{2,5 \times 10^{-3}}{...}
\sisetup{locale=FR,output-decimal-marker={,},group-minimum-digits=4,parse-numbers=false}
\DeclareSIUnit\ohm{\ensuremath{\symup{\Omega}}} % U+2126 absent de la police
\DeclareSIUnit\division{div} % réglages d'oscilloscope (ms/div, V/div)
\DeclareSIUnit\tr{tr} % tours (vitesse de rotation, ex. \SI{1500}{\tr\per\minute})
\DeclareSIUnit\molar{M} % concentration molaire (ex. \SI{3}{\molar}, \SI{1}{\micro\molar})
\DeclareSIUnit\an{an} % année (le modèle confond parfois avec \year, un compteur TeX) — ex. \SI{5}{\kilo\meter\per\an}
\DeclareSIUnit\FCFA{FCFA} % franc CFA (ex. \SI{500}{\FCFA\per\kilo\gram})
\DeclareSIUnit\degreeBrix{\textdegree Brix} % teneur en sucre d'un sirop/jus (réfractométrie)
\DeclareSIUnit\month{mois} % le modèle utilise parfois \month, un compteur TeX, comme unité de durée
\DeclareSIUnit\microliter{\micro\liter} % le modèle écrit parfois \microliter en un seul token
\DeclareSIUnit\million{millions} % dénombrement (ex. \SI{15}{\million\per\mL} spermatozoïdes)
\usepackage{qrcode}
% \degree : commande fréquemment utilisée par les modèles pour un angle ou une
% température en dehors de \SI{}{} (non fournie par les paquets chargés).
\providecommand{\degree}{^{\circ}}
% \qed (fin de démonstration), utilisé par les modèles sans amsthm
\providecommand{\qed}{\hfill$\square$}
% \barre : double barre des valeurs interdites dans un tableau de variations (tkz-tab)
\providecommand{\barre}{\ensuremath{\|}}
% environnement proof (amsthm non chargé) utilisé par les modèles
\ifdefined\proof\else\newenvironment{proof}[1][Démonstration]{\par\noindent\textit{#1.}\ }{\qed\par}\fi
\usepackage{lastpage}
\usepackage{hyperref}
\hypersetup{hidelinks}

% ============================================================
% Palette « Pop » OnBuch+ — mêmes hex que la classe P (plus_ui.dart)
% ============================================================
\definecolor{poporange}{HTML}{F59321}
\definecolor{popdark}{HTML}{E07A0C}
\definecolor{popcream}{HTML}{FFF6E8}
\definecolor{popink}{HTML}{1C1714}
\definecolor{poppurple}{HTML}{7A5AE0}
\definecolor{popgreen}{HTML}{2E9E6A}
\definecolor{poppink}{HTML}{E0457B}
\definecolor{popblue}{HTML}{2F7DE1}
\definecolor{popgold}{HTML}{C98A12}
\definecolor{popmuted}{HTML}{8A7F76}
\definecolor{popline}{HTML}{EADFD2}
\definecolor{popgray}{HTML}{8A7F76}
\colorlet{popgrey}{popgray}
% Alias tolérants (le modèle nomme parfois les couleurs différemment)
\colorlet{popwhite}{white}
\colorlet{popblack}{popink}
\colorlet{popred}{poppink}
\colorlet{popyellow}{popgold}
% Teintes claires (fonds de tableaux/encadrés) — le modèle nomme parfois
% les couleurs avec un suffixe L (ex. popblueL) sans les avoir définies.
\colorlet{poporangeL}{poporange!20}
\colorlet{popdarkL}{popdark!20}
\colorlet{poppurpleL}{poppurple!20}
\colorlet{popgreenL}{popgreen!20}
\colorlet{poppinkL}{poppink!20}
\colorlet{popblueL}{popblue!20}
\colorlet{popgoldL}{popgold!20}
\colorlet{popredL}{popred!20}
\colorlet{popgrayL}{popgray!20}
% Teintes claires pour fonds de tableaux / zones
\colorlet{popblueL}{popblue!10!white}
\colorlet{poporangeL}{poporange!14!white}
\colorlet{popgreenL}{popgreen!12!white}
\colorlet{poppurpleL}{poppurple!12!white}
\colorlet{poppinkL}{poppink!12!white}
\colorlet{popgoldL}{popgold!14!white}

\pagecolor{popcream}

% ============================================================
% Typographie
% ============================================================
\defaultfontfeatures{Ligatures=TeX}
\setmainfont{PlusJakartaSans-Medium.ttf}[Path=../../../fonts/,BoldFont=PlusJakartaSans-Bold.ttf]
% Maths en sans-serif (Fira Math) avec chiffres et lettres droites de Plus
% Jakarta, pour que les formules se fondent dans le texte.
\usepackage[math-style=ISO,bold-style=ISO]{unicode-math}
\setmathfont{FiraMath-Regular.otf}[Path=../../../fonts/]
\setmathfont{PlusJakartaSans-Medium.ttf}[Path=../../../fonts/,range={up/num,up/{Latin,latin},\mathup}]
\usepackage{icomma} % virgule décimale sans espace en mode math
% Caractères mathématiques Unicode absents des polices de texte (Plus Jakarta,
% Archivo Black) : rendus par leur équivalent mathématique, en texte comme en
% formule — sinon ils s'affichent en carré vide (ex. « Arithmétique dans ℤ »).
\usepackage{newunicodechar}
\newunicodechar{ℝ}{\ensuremath{\mathbb{R}}}
\newunicodechar{ℤ}{\ensuremath{\mathbb{Z}}}
\newunicodechar{ℂ}{\ensuremath{\mathbb{C}}}
\newunicodechar{ℕ}{\ensuremath{\mathbb{N}}}
\newunicodechar{ℚ}{\ensuremath{\mathbb{Q}}}
\newunicodechar{𝔻}{\ensuremath{\mathbb{D}}}
\newunicodechar{α}{\ensuremath{\alpha}}
\newunicodechar{β}{\ensuremath{\beta}}
\newunicodechar{γ}{\ensuremath{\gamma}}
\newunicodechar{δ}{\ensuremath{\delta}}
\newunicodechar{ε}{\ensuremath{\varepsilon}}
\newunicodechar{θ}{\ensuremath{\theta}}
\newunicodechar{λ}{\ensuremath{\lambda}}
\newunicodechar{μ}{\ensuremath{\mu}}
\newunicodechar{σ}{\ensuremath{\sigma}}
\newunicodechar{τ}{\ensuremath{\tau}}
\newunicodechar{φ}{\ensuremath{\varphi}}
\newunicodechar{ψ}{\ensuremath{\psi}}
\newunicodechar{ω}{\ensuremath{\omega}}
\newunicodechar{Σ}{\ensuremath{\Sigma}}
% majuscules produites par \MakeUppercase dans les titres (α-aminés -> Α)
\newunicodechar{Α}{\ensuremath{\alpha}}
\newunicodechar{Β}{\ensuremath{\beta}}
\newunicodechar{Γ}{\ensuremath{\Gamma}}
\newunicodechar{Δ}{\ensuremath{\Delta}}
\newunicodechar{Ε}{\ensuremath{\varepsilon}}
\newunicodechar{Θ}{\ensuremath{\Theta}}
\newunicodechar{Λ}{\ensuremath{\Lambda}}
\newunicodechar{Μ}{\ensuremath{\mu}}
\newunicodechar{Τ}{\ensuremath{\tau}}
\newunicodechar{Φ}{\ensuremath{\Phi}}
\newunicodechar{Ψ}{\ensuremath{\Psi}}
\newunicodechar{Ω}{\ensuremath{\Omega}}
\newunicodechar{↦}{\ensuremath{\mapsto}}
\newunicodechar{∈}{\ensuremath{\in}}
\newunicodechar{∉}{\ensuremath{\notin}}
\newunicodechar{∧}{\ensuremath{\wedge}}
\newunicodechar{⇔}{\ensuremath{\Leftrightarrow}}
\newunicodechar{⟺}{\ensuremath{\Longleftrightarrow}}
\newunicodechar{⇒}{\ensuremath{\Rightarrow}}
\newunicodechar{⟹}{\ensuremath{\Longrightarrow}}
\newunicodechar{∘}{\ensuremath{\circ}}
\newunicodechar{∪}{\ensuremath{\cup}}
\newunicodechar{∩}{\ensuremath{\cap}}
\newunicodechar{≡}{\ensuremath{\equiv}}
\newunicodechar{⊂}{\ensuremath{\subset}}
\newunicodechar{⊥}{\ensuremath{\perp}}
\newunicodechar{∃}{\ensuremath{\exists}}
\newunicodechar{∀}{\ensuremath{\forall}}
\newunicodechar{∛}{\ensuremath{\sqrt[3]{\,}}}
\newunicodechar{∜}{\ensuremath{\sqrt[4]{\,}}}
\newunicodechar{⃗}{\ensuremath{{}^{\to}}}
\newunicodechar{ⁿ}{\ifmmode^{n}\else\textsuperscript{\ensuremath{n}}\fi}
\newunicodechar{ˣ}{\ifmmode^{x}\else\textsuperscript{\ensuremath{x}}\fi}
\newunicodechar{ᵇ}{\ifmmode^{b}\else\textsuperscript{\ensuremath{b}}\fi}
\newunicodechar{ʳ}{\ifmmode^{r}\else\textsuperscript{\ensuremath{r}}\fi}
\newunicodechar{ᵘ}{\ifmmode^{u}\else\textsuperscript{\ensuremath{u}}\fi}
\newunicodechar{ʸ}{\ifmmode^{y}\else\textsuperscript{\ensuremath{y}}\fi}
\newunicodechar{ᵃ}{\ifmmode^{a}\else\textsuperscript{\ensuremath{a}}\fi}
\newunicodechar{ᵉ}{\ifmmode^{e}\else\textsuperscript{\ensuremath{e}}\fi}
\newunicodechar{ᵏ}{\ifmmode^{k}\else\textsuperscript{\ensuremath{k}}\fi}
\newunicodechar{ᵖ}{\ifmmode^{p}\else\textsuperscript{\ensuremath{p}}\fi}
\newunicodechar{ᴺ}{\ifmmode^{N}\else\textsuperscript{\ensuremath{N}}\fi}
\newunicodechar{ᶜ}{\ifmmode^{c}\else\textsuperscript{\ensuremath{c}}\fi}
\newunicodechar{ʰ}{\ifmmode^{h}\else\textsuperscript{\ensuremath{h}}\fi}
\newunicodechar{ᵠ}{\ifmmode^{q}\else\textsuperscript{\ensuremath{q}}\fi}
\newunicodechar{ₙ}{\ifmmode_{n}\else\textsubscript{\ensuremath{n}}\fi}
\newunicodechar{ₐ}{\ifmmode_{a}\else\textsubscript{\ensuremath{a}}\fi}
\newunicodechar{ᵢ}{\ifmmode_{i}\else\textsubscript{\ensuremath{i}}\fi}
\newunicodechar{ᵣ}{\ifmmode_{r}\else\textsubscript{\ensuremath{r}}\fi}
\newunicodechar{ₖ}{\ifmmode_{k}\else\textsubscript{\ensuremath{k}}\fi}
\newunicodechar{ᵦ}{\ifmmode_{\beta}\else\textsubscript{\ensuremath{\beta}}\fi}
\newunicodechar{ₓ}{\ifmmode_{x}\else\textsubscript{\ensuremath{x}}\fi}
\newfontfamily\popdisplay{ArchivoBlack-Regular.ttf}[Path=../../../fonts/]
\newfontfamily\poplogo{SpaceGrotesk-Bold.ttf}[Path=../../../fonts/]
\newfontfamily\popsemi{PlusJakartaSans-SemiBold.ttf}[Path=../../../fonts/]
\newfontfamily\popxbold{PlusJakartaSans-ExtraBold.ttf}[Path=../../../fonts/]
\newfontfamily\popxboldls{PlusJakartaSans-ExtraBold.ttf}[Path=../../../fonts/,LetterSpace=3.0]

\setlist[enumerate]{leftmargin=1.7em,itemsep=5pt,topsep=4pt}
\setlist[itemize]{leftmargin=1.7em,itemsep=4pt,topsep=4pt,label={\color{poporange}\textbullet}}
\setlength{\parindent}{0pt}
\setlength{\parskip}{5pt}
\renewcommand{\arraystretch}{1.25}

% pgfplots : style Pop par défaut
\pgfplotsset{
  every axis/.append style={axis line style={popink,line width=1pt},tick style={popink},
    grid style={popline,line width=0.5pt},label style={font=\small},tick label style={font=\footnotesize}},
  popcurve/.style={poporange,line width=1.8pt,no markers,smooth},
}
% Même style disponible dans un \draw TikZ simple (hors pgfplots)
\tikzset{popcurve/.style={poporange,line width=1.8pt}}
\tikzset{popgrid/.style={popline,very thin}}
\colorlet{popcurve}{poporange} % le nom du style est parfois utilisé comme couleur

% ============================================================
% Pill / squiggle
% ============================================================
\newcommand{\poppill}[3]{%
  \tikz[baseline=-0.55ex]\node[draw=popink,line width=1.2pt,rounded corners=2.6pt,%
    fill=#2,inner xsep=6.5pt,inner ysep=3pt]{%
    \popxboldls\fontsize{7}{7}\selectfont\color{#3}\MakeUppercase{#1}};%
}
\newcommand{\popsquiggle}[1]{%
  \tikz[baseline=-0.4ex]{
    \draw[#1,line width=2.6pt,line cap=round]
      plot [smooth] coordinates {(0,0.09) (0.34,0.01) (0.68,0.21) (1.02,0.01) (1.36,0.10)};
  }%
}
% Mot-clé surligné (marqueur orange sous le texte)
\newcommand{\cle}[1]{{\popxbold\color{popdark}#1}}

% ============================================================
% En-tête / pied
% ============================================================
\pagestyle{fancy}
\fancyhf{}
\renewcommand{\headrulewidth}{0pt}
\renewcommand{\footrulewidth}{0pt}
\lhead{\raisebox{-0.15ex}{\poplogo\fontsize{13}{13}\selectfont\color{popink}OnBuch\color{poporange}+}}
\rhead{\poppill{\DOCMATIERE\ \textbullet\ \DOCNIVEAU\ \textbullet\ Leçon \DOCLECON}{white}{popink}}
\lfoot{\popxbold\fontsize{7.6}{7.6}\selectfont\color{popmuted}\MakeUppercase{Cours signé OnBuch+}\ \textbullet\ {\popsemi\DOCTITRE}}
\rfoot{\tikz[baseline=-0.6ex]\node[rounded corners=8.5pt,fill=popink,minimum height=17pt,minimum width=17pt,inner xsep=5pt,inner sep=0]{\popxbold\fontsize{8}{8}\selectfont\color{popcream}\thepage\,/\,\pageref*{LastPage}};}

% ============================================================
% Titres
% ============================================================
\newcommand{\chaptitre}[2]{%
  \begin{tcolorbox}[enhanced,colback=poporange,colframe=popink,boxrule=2.6pt,arc=11pt,
      drop shadow=popink,left=15pt,right=15pt,top=12pt,bottom=11pt,before skip=3pt,after skip=18pt]
    \poppill{#2}{popink}{popcream}\par\vspace{9pt}
    {\raggedright\hyphenpenalty=10000\exhyphenpenalty=10000\popdisplay\fontsize{22}{24}\selectfont\color{popcream}\MakeUppercase{#1}\par}
  \end{tcolorbox}
}
% Section numérotée : pastille orange avec le numéro + titre Archivo
\newcounter{popsec}
\newcommand{\coursec}[1]{%
  \stepcounter{popsec}\setcounter{popsub}{0}%
  \par\vspace{16pt}\noindent
  \begin{minipage}{\linewidth}
  \tikz[baseline=-0.7ex]\node[draw=popink,line width=1.6pt,fill=poporange,rounded corners=4pt,minimum size=22pt,inner sep=0]{\popdisplay\fontsize{12}{12}\selectfont\color{popcream}\thepopsec};%
  \hspace{8pt}{\popdisplay\fontsize{15}{17}\selectfont\color{popink}\MakeUppercase{#1}}\par
  \vspace{4pt}\noindent{\color{poporange}\rule{36pt}{3.6pt}}
  \end{minipage}\par\nopagebreak\vspace{8pt}
}
\newcounter{popsub}
\newcommand{\courssub}[1]{%
  \stepcounter{popsub}%
  \par\vspace{9pt}\noindent{\popxbold\fontsize{11.5}{13}\selectfont\color{popink}{\color{poporange}\thepopsec.\thepopsub}\hspace{6pt}#1}\par\nopagebreak\vspace{4pt}
}

% ============================================================
% Boîtes pédagogiques — même recette : fond blanc, bordure épaisse colorée,
% ombre dure encre, badge plein. Titre optionnel [#1].
% ============================================================
\tcbset{popbase/.style={breakable,enhanced,colback=white,boxrule=2.3pt,arc=9pt,
  shadow={3pt}{-3pt}{0pt}{popink},left=14pt,right=14pt,top=11pt,bottom=11pt,before skip=11pt,after skip=15pt,
  fontupper=\popsemi\fontsize{10.6}{14.5}\selectfont\color{popink},
  fontlower=\popsemi\fontsize{10.6}{14.5}\selectfont\color{popink}}}
% Coche dessinée (le glyphe ✓ n'existe pas dans les polices du document)
\newcommand{\popcheck}[1]{\tikz[baseline=-0.5ex]\draw[#1,line width=1.6pt,line cap=round,line join=round](-0.13,0.0)--(-0.04,-0.09)--(0.13,0.1);}
\providecommand{\coche}{\popcheck{popgreen}}
\providecommand{\resultat}[1]{\par\smallskip\noindent\fcolorbox{popgreen}{popgreenL}{\parbox{\dimexpr\linewidth-2\fboxsep-2\fboxrule\relax}{\textbf{Résultat.} #1}}\par\smallskip}
\newcommand{\popboxhead}[3]{\poppill{#1}{#2}{white}\ifx\relax#3\relax\else\hspace{6pt}{\popxbold\fontsize{10.6}{12}\selectfont\color{popink}#3}\fi\par\vspace{7pt}}

\newtcolorbox{aretenir}[1][]{popbase,colframe=popgreen,before upper={\popboxhead{À retenir}{popgreen}{#1}}}
\newtcolorbox{exemplebox}[1][]{popbase,colframe=poporange,before upper={\popboxhead{Exemple}{poporange}{#1}}}
\newtcolorbox{definition}[1][]{popbase,colframe=popblue,before upper={\popboxhead{Définition}{popblue}{#1}}}
\newtcolorbox{propriete}[1][]{popbase,colframe=poppurple,before upper={\popboxhead{Propriété}{poppurple}{#1}}}
\newtcolorbox{methode}[1][]{popbase,colframe=popgold,before upper={\popboxhead{Méthode}{popgold}{#1}}}
\newtcolorbox{attention}[1][]{popbase,colframe=poppink,before upper={\popboxhead{Attention}{poppink}{#1}}}
\newtcolorbox{experience}[1][]{popbase,colframe=popblue,colback=popblueL,before upper={\popboxhead{Expérience}{popblue}{#1}}}
% « Le savais-tu ? » : carte violette pleine (contexte, culture, Cameroun)
\newtcolorbox{savaistu}[1][]{popbase,colback=poppurple,colframe=popink,
  fontupper=\popsemi\fontsize{10.4}{14.2}\selectfont\color{white},
  before upper={\poppill{Le savais-tu\,?}{white}{poppurple}\ifx\relax#1\relax\else\hspace{6pt}{\popxbold\color{white}#1}\fi\par\vspace{7pt}}}
% Exercice résolu : énoncé en haut, \tcblower puis la solution sur fond crème
\newcounter{popexo}\newcounter{popexr}
\newtcolorbox{exoresolu}[1][]{popbase,colframe=poporange,colbacklower=popcream,
  segmentation style={popink,line width=1.2pt,dashed},
  before upper={\stepcounter{popexr}\popboxhead{Exercice résolu \thepopexr}{poporange}{#1}},
  before lower={\poppill{Solution}{popink}{popcream}\par\vspace{6pt}}}
% Exercice (non résolu en place) — la correction est dans la partie Corrigés
\newtcolorbox{exercice}[1][]{popbase,colframe=popink,
  before upper={\stepcounter{popexo}\popboxhead{Exercice \thepopexo}{popink}{#1}}}
% Corrigé
\newtcolorbox{corrige}[1][]{popbase,colframe=popgreen,colback=popgreenL,
  before upper={\popboxhead{Corrigé}{popgreen}{#1}}}
% Objectifs / prérequis / bilan
\newtcolorbox{objectifs}{popbase,colframe=popink,colback=poporangeL,
  before upper={\popboxhead{Objectifs de la leçon}{popink}{}}}
\newtcolorbox{prerequis}{popbase,colframe=popink,colback=popblueL,
  before upper={\popboxhead{Prérequis}{popblue}{}}}
\newtcolorbox{bilan}{popbase,colframe=popink,colback=white,boxrule=2.8pt,
  before upper={\popboxhead{Fiche bilan}{poporange}{L'essentiel à connaître pour l'examen}}}
% Figure encadrée avec légende
\newtcolorbox{popfigure}[1][]{enhanced,breakable=false,colback=white,colframe=popline,boxrule=1.4pt,arc=8pt,
  left=10pt,right=10pt,top=10pt,bottom=8pt,before skip=10pt,after skip=12pt,halign=center,
  lower separated=false,halign lower=center,
  fontlower=\popsemi\fontsize{9}{11.5}\selectfont\color{popmuted},
  #1}
\newcommand{\legende}[1]{\par\vspace{6pt}{\popsemi\fontsize{9}{11.5}\selectfont\color{popmuted}#1}}

% ============================================================
% Signature OnBuch+
% ============================================================
\newcommand{\poplogoinline}[1]{{\poplogo\fontsize{#1}{#1}\selectfont\color{popink}OnBuch\color{poporange}+}}
\newcommand{\signatureonbuch}{%
  \begin{tcolorbox}[enhanced,colback=popink,colframe=popink,boxrule=2.2pt,arc=10pt,
      drop shadow=poporange,left=14pt,right=14pt,top=10pt,bottom=10pt,before skip=0pt,after skip=0pt]
    \tikz[baseline=-0.7ex]\node[circle,fill=popgreen,draw=popcream,line width=1.3pt,minimum size=22pt,inner sep=0]{\popcheck{white}};%
    \hspace{9pt}\begin{minipage}[c]{0.62\linewidth}
      {\popxbold\fontsize{12}{13}\selectfont\color{popcream}Signé\ \textbullet\ {\poplogo OnBuch\color{poporange}+}}\par\vspace{2pt}
      {\raggedright\popsemi\fontsize{8}{10.5}\selectfont\color{popline}\MakeUppercase{Équipe pédagogique OnBuch+}\\\MakeUppercase{Conforme au programme officiel MINESEC}\par}
    \end{minipage}\hfill
    {\poplogo\fontsize{22}{22}\selectfont\color{popcream}OnBuch\color{poporange}+}
  \end{tcolorbox}%
}

% ============================================================
% Page de garde
% #1 titre · #2 matière · #3 niveau · #4 module · #5 n° leçon · #6 durée ·
% #7 description / objectifs (texte ou itemize)
% ============================================================
\newcommand{\pagegarde}[7]{%
  \thispagestyle{empty}
  \newgeometry{a4paper,margin=1.7cm,top=1.6cm,bottom=1.4cm}%
  \begin{tikzpicture}[remember picture,overlay]
    \fill[poporange!18] ([xshift=-3.2cm,yshift=-3.6cm]current page.north east) circle (4.6cm);
    \fill[poppurple!14] ([xshift=2.4cm,yshift=7.6cm]current page.south west) circle (3.8cm);
  \end{tikzpicture}%
  {\poplogo\fontsize{30}{30}\selectfont\color{popink}OnBuch\color{poporange}+}\hfill
  \raisebox{0.6ex}{\poppill{Cours signé \textbullet\ Premium}{popink}{popcream}}\par
  \vspace{1pt}\popsquiggle{poppurple}\par
  \vspace{8pt}
  \begin{tcolorbox}[enhanced,colback=poporange,colframe=popink,boxrule=2.8pt,arc=13pt,
      drop shadow=popink,left=18pt,right=18pt,top=12pt,bottom=13pt,after skip=10pt]
    \poppill{#2 \textbullet\ #3}{popink}{popcream}\hspace{5pt}\poppill{#4}{white}{popink}\par\vspace{8pt}
    {\popdisplay\fontsize{14}{14}\selectfont\color{popink}LEÇON #5}\par\vspace{4pt}
    {\raggedright\hyphenpenalty=10000\exhyphenpenalty=10000\popdisplay\fontsize{25}{27}\selectfont\color{popcream}\MakeUppercase{#1}\par}
    \vspace{8pt}
    {\popxbold\fontsize{9.5}{9.5}\selectfont\color{popink}\MakeUppercase{Durée conseillée : #6}}
  \end{tcolorbox}
  \begin{tcolorbox}[enhanced,colback=white,colframe=popink,boxrule=2.2pt,arc=10pt,
      drop shadow=popink,left=16pt,right=16pt,top=10pt,bottom=10pt,after skip=8pt]
    \poppill{Dans cette leçon}{poppurple}{white}\par\vspace{5pt}
    \popsemi\fontsize{9.3}{12.6}\selectfont\color{popink}#7
  \end{tcolorbox}
  \noindent\begin{minipage}[t]{0.487\linewidth}
    \begin{tcolorbox}[enhanced,colback=popgreen,colframe=popink,boxrule=2.2pt,arc=10pt,
        drop shadow=popink,left=11pt,right=11pt,top=10pt,bottom=10pt,height=3.1cm,valign=center]
      \tcbox[colback=white,colframe=popink,boxrule=1.3pt,arc=5pt,boxsep=0pt,nobeforeafter,
        left=3pt,right=3pt,top=3pt,bottom=3pt,valign=center]{\qrcode[height=1.9cm]{https://play.google.com/store/apps/details?id=com.luvvix.onbuch&hl=fr}}%
      \hspace{8pt}\begin{minipage}[c]{0.5\linewidth}
        {\popdisplay\fontsize{11}{12.5}\selectfont\color{white}\MakeUppercase{Télécharge OnBuch}}\par\vspace{4pt}
        {\raggedright\popsemi\fontsize{8.4}{11}\selectfont\color{white}Gratuit \textbullet\ Google Play\\Tuteur IA, annales, concours, résultats\par}
      \end{minipage}
    \end{tcolorbox}
  \end{minipage}\hfill
  \begin{minipage}[t]{0.487\linewidth}
    \begin{tcolorbox}[enhanced,colback=poppurple,colframe=popink,boxrule=2.2pt,arc=10pt,
        drop shadow=popink,left=11pt,right=11pt,top=10pt,bottom=10pt,height=3.1cm,valign=center]
      \tikz[baseline=-0.7ex]\node[circle,fill=white,draw=popink,line width=1.3pt,minimum size=1.6cm,inner sep=0]{\popdisplay\fontsize{24}{24}\selectfont\color{poppurple}+};%
      \hspace{8pt}\begin{minipage}[c]{0.6\linewidth}
        {\popdisplay\fontsize{11}{12.5}\selectfont\color{white}\MakeUppercase{Passe à OnBuch+}}\par\vspace{4pt}
        {\raggedright\popsemi\fontsize{8.4}{11}\selectfont\color{white}Tous les cours signés, Tuteur IA illimité, fiches PDF — onglet OnBuch+ de l'app\par}
      \end{minipage}
    \end{tcolorbox}
  \end{minipage}
  \vspace{8pt}
  \popcontenu
  \vfill
  \signatureonbuch
  \restoregeometry
  \newpage
}

% Rangée « Ce cours contient » (page de garde)
\newcommand{\poptile}[3]{% #1 couleur #2 gros texte #3 légende
  \begin{tcolorbox}[enhanced,colback=#1,colframe=popink,boxrule=2pt,arc=9pt,shadow={2.5pt}{-2.5pt}{0pt}{popink},
      left=6pt,right=6pt,top=9pt,bottom=9pt,halign=center,valign=center,height=1.75cm,width=\linewidth,nobeforeafter]
    {\popdisplay\fontsize{12.5}{13}\selectfont\color{white}\MakeUppercase{#2}}\par\vspace{3pt}
    {\centering\popsemi\fontsize{7.8}{9.5}\selectfont\color{white}#3\par}
  \end{tcolorbox}}
\newcommand{\popcontenu}{%
  \par\noindent{\popxboldls\fontsize{7.5}{7.5}\selectfont\color{popmuted}CE COURS SIGNÉ CONTIENT}\par\vspace{7pt}
  \noindent\begin{minipage}[t]{0.235\linewidth}\poptile{popblue}{Cours}{complet, illustré, pas à pas}\end{minipage}\hfill
  \begin{minipage}[t]{0.235\linewidth}\poptile{poporange}{Méthodes}{et exercices résolus}\end{minipage}\hfill
  \begin{minipage}[t]{0.235\linewidth}\poptile{poppink}{Exercices}{type examen + corrigés}\end{minipage}\hfill
  \begin{minipage}[t]{0.235\linewidth}\poptile{popgreen}{Fiche bilan}{l'essentiel à retenir}\end{minipage}\par}

% Sommaire
\newtcolorbox{sommaire}{enhanced,colback=white,colframe=popink,boxrule=2.2pt,arc=10pt,
  drop shadow=popink,left=18pt,right=18pt,top=13pt,bottom=13pt,before skip=6pt,after skip=20pt,
  before upper={\poppill{Sommaire}{poppurple}{white}\par\vspace{9pt}\popsemi\fontsize{10.6}{14.5}\selectfont\color{popink}}}

% Signature de fin de cours — poussée tout en bas de la dernière page
\newcommand{\findecours}{%
  \par\vfill\nopagebreak
  \begin{center}\popsquiggle{poporange}\end{center}\vspace{4pt}
  \signatureonbuch
}
\AtBeginDocument{\def\llbracket{\mathopen{[\mkern-3.5mu[}}\def\rrbracket{\mathclose{]\mkern-3.5mu]}}} % intervalles d'entiers (glyphe ⟦ absent de la police)
% Glyphes absents de Fira Math : redéfinis à partir de glyphes disponibles
\AtBeginDocument{%
  \def\setminus{\mathbin{\backslash}}\let\smallsetminus\setminus
  \def\vdots{\mathord{\vbox{\baselineskip3.2pt\lineskiplimit0pt\kern4pt\hbox{.}\hbox{.}\hbox{.}}}}%
  \def\ddots{\mathinner{\mkern1mu\raise6.5pt\hbox{.}\mkern2mu\raise3.6pt\hbox{.}\mkern2mu\raise0.7pt\hbox{.}\mkern1mu}}%
  \def\triangle{\mathord{\scalebox{0.9}{$\symup{\Delta}$}}}%
  \def\nabla{\mathord{\raisebox{\depth}{\scalebox{1}[-1]{$\symup{\Delta}$}}}}%
  \def\checkmark{\popcheck{popdark}}%
  \def\star{\mathbin{\ast}}\def\bigstar{\mathord{\ast}}%
  \def\diamond{\mathbin{\scalebox{0.6}{$\Diamond$}}}%
  \def\bigcup{\mathop{\mathchoice{\raisebox{-0.3em}{\scalebox{1.7}{$\cup$}}}{\scalebox{1.25}{$\cup$}}{\cup}{\cup}}\displaylimits}%
  \def\bigcap{\mathop{\mathchoice{\raisebox{-0.3em}{\scalebox{1.7}{$\cap$}}}{\scalebox{1.25}{$\cap$}}{\cap}{\cap}}\displaylimits}%
  \def\lesseqqgtr{\mathrel{\lessgtr}}\def\bigoplus{\mathop{\oplus}\displaylimits}%
}
\providecommand{\ssi}{si et seulement si} % abréviation fréquente
\AtBeginDocument{\providecommand{\wideparen}{\overparen}} % arc de cercle

%% FIGFIX-BEGIN v5 — correctifs globaux des figures (docs/onbuch-generation/tools/figfix)
\usepackage{adjustbox}\usepackage{etoolbox}\tcbuselibrary{raster}
% toute figure TikZ/pgfplots plus large que la ligne est réduite à la largeur disponible
\BeforeBeginEnvironment{tikzpicture}{\begin{adjustbox}{max width=\linewidth}}
\AfterEndEnvironment{tikzpicture}{\end{adjustbox}}
% légendes pgfplots lisibles par-dessus les courbes
\pgfplotsset{every axis/.append style={legend style={fill=white,fill opacity=0.92,text opacity=1,draw=black!15,inner sep=3pt}}}
% étiquettes d'arêtes (syntaxe edge["..."]) avec fond blanc
\tikzset{every edge quotes/.append style={fill=white,inner sep=1.2pt}}
%% FIGFIX-END
\begin{document}

\pagegarde{Généralités sur les moteurs à combustion interne}{PCT}{3ème}{Module 4 : Technologie et mécanique}{N18}{1 séance (1 h chacune)}{Cette leçon introduit les moteurs à combustion interne, machines qui convertissent l'énergie chimique d'un carburant en énergie mécanique. Elle présente leur définition, leur principe de fonctionnement et leurs différents classements (cycle, carburant, allumage, refroidissement), illustrés par les véhicules de la vie quotidienne camerounaise.
\vspace{4pt}
\begin{multicols}{2}\small
\begin{itemize}
\item À la fin de cette leçon, je sais définir un moteur à combustion interne et le distinguer d'un moteur à combustion externe.
\item Je sais identifier les organes essentiels d'un moteur à combustion interne (cylindre, piston, bielle, vilebrequin, soupapes, bougie ou injecteur) sur un schéma.
\item Je sais décrire les quatre temps du cycle d'un moteur à essence 4 temps.
\item Je sais classer les moteurs selon le cycle (2 temps / 4 temps), le carburant (essence / diesel / GPL), le mode d'allumage et le mode de refroidissement.
\item Je sais citer des applications concrètes de chaque type de moteur dans la vie courante (moto-taxi, groupe électrogène, camion, pirogue motorisée).
\item Je sais interpréter un schéma de moteur et y repérer la transformation d'énergie réalisée.
\end{itemize}\end{multicols}}

\begin{sommaire}
\begin{multicols}{2}
\begin{enumerate}
\item Définition et rôle d'un moteur à combustion interne
\item Constitution d'un moteur à combustion interne
\item Le cycle à 4 temps du moteur à essence
\item Le moteur à 2 temps et le moteur diesel
\item Classement des moteurs à combustion interne
\item Sécurité, entretien et impact environnemental
\item Activité d'intégration
\item Méthodes et exercices résolus
\item Exercices
\item Corrigés des exercices
\item Fiche bilan
\end{enumerate}
\end{multicols}
\end{sommaire}

\begin{objectifs}
\begin{itemize}
\item À la fin de cette leçon, je sais définir un moteur à combustion interne et le distinguer d'un moteur à combustion externe.
\item Je sais identifier les organes essentiels d'un moteur à combustion interne (cylindre, piston, bielle, vilebrequin, soupapes, bougie ou injecteur) sur un schéma.
\item Je sais décrire les quatre temps du cycle d'un moteur à essence 4 temps.
\item Je sais classer les moteurs selon le cycle (2 temps / 4 temps), le carburant (essence / diesel / GPL), le mode d'allumage et le mode de refroidissement.
\item Je sais citer des applications concrètes de chaque type de moteur dans la vie courante (moto-taxi, groupe électrogène, camion, pirogue motorisée).
\item Je sais interpréter un schéma de moteur et y repérer la transformation d'énergie réalisée.
\item Je sais appliquer les consignes de sécurité liées à l'usage et à l'entretien des moteurs (carburants inflammables, gaz d'échappement toxiques, pièces chaudes).
\end{itemize}
\end{objectifs}

\begin{prerequis}
\begin{itemize}
\item Les transformations d'énergie et la chaîne énergétique (leçon sur l'énergie, 3ème)
\item La combustion : combustible, comburant, énergie libérée (chimie, 3ème)
\item Notion de pression et de volume d'un gaz (physique, 4ème/3ème)
\item Les mouvements : translation, rotation, transformation de mouvement (technologie, 4ème)
\end{itemize}
\end{prerequis}

\newpage

\coursec{Définition et rôle d'un moteur à combustion interne}

\courssub{Situation d'accroche : le benskin qui refuse de démarrer}

Il est 6 h 30 à Douala, quartier Deïdo. Jean, conducteur de moto-taxi « benskin », tourne la clé de contact, appuie sur le démarreur\ldots{} rien. Le moteur toussote, puis se tait. Sa journée de travail est menacée. Il pousse sa moto jusqu'à l'atelier de Papa Etame, le mécanicien du quartier. Celui-ci démonte quelques pièces, renifle la bougie, et annonce : « Ton moteur est un moteur à essence à 4 temps. Ce n'est pas comme le moteur diesel du camion qui livre le marché, ni comme le petit moteur 2 temps de la mobylette de ton voisin. »

Jean est intrigué. Pourquoi parle-t-on de moteur « à \cle{combustion interne} » ? Qu'est-ce qui distingue un moteur à essence d'un moteur diesel, un moteur 2 temps d'un moteur 4 temps ? Et surtout : comment cette machine transforme-t-elle l'essence achetée à la station en mouvement qui fait avancer la moto ?

Ces machines sont partout au Cameroun : motos-taxis, voitures, camions de la côte, groupes électrogènes qui secourent les quartiers pendant les délestages, pirogues motorisées du Wouri, motopompes qui irriguent les champs de l'Extrême-Nord. Comprendre leur fonctionnement, c'est comprendre le cœur mécanique du pays. C'est l'objet de cette leçon, et nous commençons par la question fondamentale : \textbf{qu'est-ce qu'un moteur à combustion interne et quel est son rôle ?}

\courssub{Rappel : qu'est-ce qu'un moteur ?}

Avant de parler de combustion, revenons à une idée simple que tu connais déjà : l'\cle{énergie} ne se crée pas, elle se \emph{transforme}. Quand tu pédales sur un vélo, tu transformes l'énergie chimique de ton repas (attiéké, ndolé, eru\ldots) en énergie mécanique de mouvement. Quand ENEO alimente ton ventilateur, l'énergie électrique devient énergie mécanique (les pales tournent) et énergie thermique (le moteur chauffe un peu).

\begin{definition}[Moteur]
Un \cle{moteur} est une \textbf{machine qui transforme une forme d'énergie en énergie mécanique}, c'est-à-dire en mouvement (translation ou rotation).
\end{definition}

Le rôle d'un moteur est donc de \emph{convertir}, jamais de « fabriquer » de l'énergie. Cette précision est essentielle, et nous y reviendrons dans les pièges à éviter.

Quelques exemples de conversions réalisées par des moteurs :
\begin{itemize}
\item le moteur électrique du ventilateur : énergie électrique $\rightarrow$ énergie mécanique ;
\item le moteur d'une grue de chantier au port de Douala : énergie électrique $\rightarrow$ énergie mécanique ;
\item le moteur de la moto de Jean : énergie \emph{chimique} de l'essence $\rightarrow$ énergie mécanique.
\end{itemize}

C'est ce dernier type de moteur, celui qui « mange » du carburant, qui nous intéresse ici.

\courssub{Le fait expérimental fondateur : les gaz de combustion poussent}

\courssub{Expérience du canon à piston}

Toute la technologie des moteurs à combustion interne repose sur un constat expérimental simple et spectaculaire : \textbf{quand on brûle un carburant dans une enceinte fermée munie d'une paroi mobile, les gaz produits poussent cette paroi}.

\begin{experience}[Le canon à piston]
\textbf{Matériel :} un tube métallique robuste fermé à une extrémité (le « canon »), un piston (ou un bouchon) coulissant dans le tube, une petite quantité de carburant vaporisé (quelques gouttes d'alcool), un dispositif d'allumage (étincelle), un support fixe.

\textbf{Protocole :}
\begin{enumerate}
\item Introduire quelques gouttes d'alcool dans le tube et agiter pour le vaporiser : le mélange air + vapeur d'alcool est un mélange \cle{explosif}.
\item Placer le piston (ou le bouchon) à l'ouverture du tube.
\item Produire une étincelle à l'intérieur du tube.
\item Observer.
\end{enumerate}

\textbf{Observations :} on entend une détonation ; le piston est \textbf{violemment chassé} hors du tube ; le tube est chaud après l'expérience.

\textbf{Interprétation :} la combustion de l'alcool dans l'air dégage beaucoup de chaleur (réaction \emph{exothermique}). Les gaz produits (\ce{CO2} et vapeur d'eau), portés à haute température, se \emph{dilatent} brutalement : leur pression augmente fortement dans l'enceinte fermée. Cette pression s'exerce sur le piston, seule paroi mobile, et le \textbf{pousse} : l'énergie thermique s'est transformée en énergie mécanique.

\textbf{Conclusion :} brûler un carburant dans une enceinte fermée munie d'un piston mobile provoque le déplacement du piston. C'est le principe fondateur du moteur à combustion interne.
\end{experience}

\begin{attention}[Sécurité]
Cette expérience est une \textbf{démonstration de cours réalisée uniquement par le professeur}, avec de très petites quantités de carburant et un équipement adapté. Ne la reproduis jamais seul à la maison : une explosion non maîtrisée peut provoquer des brûlures graves et des incendies.
\end{attention}

On peut rapprocher cette expérience d'un phénomène de la vie courante : quand on gonfle un ballon et qu'on le lâche, l'air sous pression s'échappe et le ballon file dans tous les sens. Dans le moteur, ce ne sont pas des gaz qu'on a stockés sous pression, mais des gaz \emph{créés et chauffés sur place} par la combustion, qui poussent le piston.

\begin{propriete}[Principe du moteur à combustion interne]
La combustion d'un carburant \textbf{à l'intérieur} d'une enceinte fermée (la \cle{chambre de combustion}) libère des \textbf{gaz chauds sous forte pression}. Ces gaz exercent une force sur une paroi mobile (le \cle{piston}) et la \textbf{poussent directement} : l'énergie thermique de la combustion est ainsi convertie en énergie mécanique. Ce savoir est établi à partir de l'expérience du piston chassé par les gaz de combustion.
\end{propriete}

\courssub{Définition du moteur à combustion interne}

Nous pouvons maintenant énoncer la définition centrale de la leçon.

\begin{definition}[Moteur à combustion interne (MCI)]
Un \cle{moteur à combustion interne} est un moteur dans lequel la \textbf{combustion du carburant se produit à l'intérieur même du moteur}, dans une enceinte appelée \cle{chambre de combustion}. Les gaz chauds produits par cette combustion agissent \emph{directement} sur les organes mobiles du moteur (piston) pour produire le mouvement.
\end{definition}

Le mot « interne » est donc le cœur de la définition : le feu qui fournit l'énergie brûle \emph{dans} la machine, et non à côté. Le carburant utilisé (essence, gasoil, pétrole lampant, gaz) contient de l'\cle{énergie chimique} stockée dans ses molécules ; la combustion la libère.

\begin{aretenir}[L'essentiel de la définition]
Dans un moteur à combustion interne :
\begin{itemize}
\item le carburant brûle \textbf{à l'intérieur} du moteur (chambre de combustion) ;
\item les gaz chauds sous pression poussent \textbf{directement} le piston ;
\item le mouvement du piston est transformé en \textbf{rotation} par le vilebrequin.
\end{itemize}
\end{aretenir}

\courssub{Comparaison avec le moteur à combustion externe}

Pour bien comprendre ce que « interne » signifie, comparons avec son contraire historique : la \textbf{machine à vapeur}, qui a fait rouler les premiers trains et tourner les premières usines.

Dans une machine à vapeur, on brûle du bois ou du charbon \textbf{à l'extérieur} du moteur, sous une chaudière. La chaleur produite chauffe de l'eau qui se transforme en vapeur ; c'est cette vapeur, acheminée par des tuyaux, qui pousse ensuite le piston. La combustion et le moteur sont donc \emph{séparés} : on parle de \cle{moteur à combustion externe}.

\begin{exemplebox}[Deux philosophies de la machine]
\textbf{Machine à vapeur (combustion externe) :} le feu brûle sous la chaudière, à l'extérieur du cylindre. La chaleur traverse la paroi de la chaudière pour vaporiser l'eau, et la vapeur pousse le piston.

\textbf{Moteur de moto (combustion interne) :} l'essence brûle directement dans le cylindre, au-dessus du piston. Les gaz de combustion eux-mêmes poussent le piston, sans intermédiaire.
\end{exemplebox}

\begin{center}
\begin{tabularx}{\linewidth}{|l|X|X|}
\hline
\rowcolor{popblueL} \textbf{Critère} & \textbf{Combustion externe} & \textbf{Combustion interne} \\
\hline
Lieu de la combustion & À l'extérieur du moteur (foyer sous chaudière) & À l'intérieur du moteur (chambre de combustion) \\
\hline
Fluide qui pousse le piston & Vapeur d'eau chauffée par la chaleur extérieure & Gaz de combustion eux-mêmes \\
\hline
Exemple & Machine à vapeur (locomotive ancienne) & Moteur de moto, voiture, groupe électrogène \\
\hline
Encombrement, rendement & Volumineux, rendement faible & Compact, rendement meilleur \\
\hline
\end{tabularx}
\end{center}

\begin{popfigure}
\begin{tikzpicture}[scale=0.92, every node/.style={font=\small}]
% ===== Combustion externe (gauche) =====
\begin{scope}
\node[font=\small\bfseries, popblue] at (2.6,4.6) {Combustion externe (machine à vapeur)};
% chaudière
\draw[fill=popblueL, thick, popblue] (0.4,1.6) rectangle (2.6,3.0);
\node at (1.5,2.3) {Chaudière (eau)};
% feu sous la chaudière
\draw[fill=poporangeL, thick, poporange] (0.6,0.5) rectangle (2.4,1.4);
\node[popdark] at (1.5,0.95) {Foyer (feu)};
\draw[->, thick, poporange] (1.5,1.45) -- (1.5,1.6);
% tuyau vapeur
\draw[thick, popblue] (2.6,2.6) -- (3.6,2.6) -- (3.6,2.3);
\node[popblue, font=\scriptsize] at (3.3,2.85) {vapeur};
% cylindre + piston
\draw[thick, popdark] (3.6,1.7) rectangle (5.2,2.9);
\draw[fill=popgreenL, thick, popgreen] (4.0,1.75) rectangle (4.35,2.85);
\draw[thick, popdark] (4.35,2.3) -- (5.6,2.3);
\node[font=\scriptsize] at (4.17,1.5) {piston};
\node[font=\scriptsize, popdark] at (5.0,3.15) {cylindre};
\end{scope}
% ===== Combustion interne (droite) =====
\begin{scope}[xshift=7.2cm]
\node[font=\small\bfseries, poporange] at (2.6,4.6) {Combustion interne (moteur de moto)};
% cylindre
\draw[thick, popdark] (1.4,1.7) rectangle (3.4,3.9);
% chambre de combustion avec explosion (remplacement starburst par tracé manuel)
\draw[fill=poporangeL, poporange, thick] (2.4,3.35) circle (0.42);
\foreach \ang in {0,30,60,90,120,150,180,210,240,270,300,330} {
  \draw[poporange, thick] (2.4,3.35) -- ++(\ang:0.7);
}
\node[font=\scriptsize, popdark] at (2.4,4.25) {chambre de combustion};
\node[font=\scriptsize, poporange] at (2.4,3.35) {feu};
% piston
\draw[fill=popgreenL, thick, popgreen] (1.5,1.8) rectangle (3.3,2.6);
\node[font=\scriptsize] at (2.4,2.2) {piston};
% force sur piston
\draw[->, very thick, poppink] (2.4,2.95) -- (2.4,2.62);
% bielle + vilebrequin
\draw[thick, popdark] (2.4,1.8) -- (2.75,0.85);
\draw[fill=popblueL, thick, popblue] (2.4,0.55) circle (0.35);
\fill[popdark] (2.4,0.55) circle (0.05);
\node[font=\scriptsize, popblue] at (3.35,0.55) {vilebrequin};
% flèche rotation
\draw[->, thick, popblue] (3.0,0.15) arc (-40:60:0.35);
\end{scope}
\end{tikzpicture}
\legende{Comparaison des deux types de moteurs thermiques. À gauche (combustion externe) : le feu brûle hors du moteur et chauffe l'eau d'une chaudière ; la vapeur produite pousse le piston. À droite (combustion interne) : le carburant brûle dans la chambre de combustion, à l'intérieur du cylindre ; les gaz chauds poussent directement le piston, dont le mouvement fait tourner le vilebrequin.}
\end{popfigure}

L'avantage décisif de la combustion interne est la \textbf{compacité} : pas de chaudière, pas de réserve d'eau, pas de tuyauterie. C'est pourquoi ces moteurs équipent les véhicules légers et les machines portables, alors que la machine à vapeur a disparu de nos routes.

\courssub{La chaîne énergétique du moteur à combustion interne}

Suivons maintenant le chemin complet de l'énergie, de la station-service jusqu'aux roues de la moto. C'est ce qu'on appelle la \cle{chaîne énergétique}.

\begin{enumerate}
\item \textbf{Énergie chimique} : l'essence est un carburant ; ses molécules (hydrocarbures) stockent de l'énergie chimique, libérable par combustion avec le dioxygène de l'air.
\item \textbf{Énergie thermique} : dans la chambre de combustion, l'étincelle de la bougie déclenche la combustion. Cette réaction exothermique produit des gaz très chauds (\ce{CO2}, vapeur d'eau) dont la pression augmente brutalement.
\item \textbf{Énergie mécanique de translation} : les gaz sous pression poussent le piston, qui descend dans le cylindre.
\item \textbf{Énergie mécanique de rotation} : la bielle et le \cle{vilebrequin} transforment le mouvement de va-et-vient du piston en mouvement de rotation.
\item \textbf{Mouvement du véhicule} : la rotation est transmise (embrayage, boîte de vitesses, chaîne ou cardan) jusqu'aux roues.
\end{enumerate}

\begin{popfigure}
\begin{tikzpicture}[
box/.style={draw=popblue, thick, fill=popblueL, rounded corners=3pt, minimum height=0.95cm, align=center, font=\small},
ebox/.style={draw=poporange, thick, fill=poporangeL, rounded corners=3pt, minimum height=0.95cm, align=center, font=\small},
fleche/.style={->, very thick, popdark},
every node/.style={font=\small}]
\node[ebox, align=center] (carb) at (0,0){Carburant\\ (essence)};
\node[box, align=center] (comb) at (3.1,0){Combustion\\ (chambre de\\ combustion)};
\node[box, align=center] (gaz) at (6.2,0){Gaz chauds\\ sous pression};
\node[box, align=center] (pist) at (9.3,0){Piston\\ (va-et-vient)};
\node[box, align=center] (vil) at (12.4,0){Vilebrequin\\ (rotation)};
\node[ebox, align=center] (roue) at (15.3,0){Roues\\ (mouvement)};
\draw[fleche] (carb) -- (comb);
\draw[fleche] (comb) -- (gaz);
\draw[fleche] (gaz) -- (pist);
\draw[fleche] (pist) -- (vil);
\draw[fleche] (vil) -- (roue);
% formes d'énergie en dessous
\node[font=\scriptsize\itshape, poppurple, align=center] at (1.55,-1.15) {énergie\\ chimique};
\node[font=\scriptsize\itshape, poppurple, align=center] at (4.65,-1.15) {énergie\\ thermique};
\node[font=\scriptsize\itshape, poppurple, align=center] at (10.8,-1.15) {énergie mécanique};
% Accolade supprimée (librairie decorations.pathreplacing non dispo)
\end{tikzpicture}
\legende{Chaîne énergétique d'un moteur à combustion interne : l'énergie chimique du carburant devient énergie thermique par la combustion, puis énergie mécanique grâce à la poussée des gaz sur le piston et à la rotation du vilebrequin, transmise enfin aux roues.}
\end{popfigure}

\begin{aretenir}[Chaîne énergétique à connaître par cœur]
\[
\text{Énergie chimique} \;\xrightarrow{\text{combustion}}\; \text{Énergie thermique} \;\xrightarrow{\text{poussée des gaz}}\; \text{Énergie mécanique}
\]
Le piston fournit un mouvement de \emph{translation} (va-et-vient) ; le vilebrequin le convertit en \emph{rotation}.
\end{aretenir}

\begin{attention}[Erreur fréquente : le moteur ne « fabrique » pas d'énergie]
Beaucoup d'élèves écrivent au BEPC : « le moteur produit de l'énergie » ou « le moteur crée du mouvement ». C'est faux et cela coûte des points. Un moteur \textbf{transforme} (ou \emph{convertit}) une énergie déjà existante : ici, l'énergie chimique du carburant. Sans essence dans le réservoir, le moteur le plus perfectionné ne produit rien du tout.

Autre confusion à éviter : ne pas confondre le \textbf{moteur} (qui transforme une énergie en énergie mécanique) et le \textbf{générateur électrique} comme l'alternateur (qui transforme l'énergie mécanique en énergie électrique). Dans un groupe électrogène, les deux sont associés : le moteur à essence fait tourner l'alternateur, qui fournit le courant. Le moteur consomme de l'énergie chimique ; l'alternateur consomme de l'énergie mécanique.
\end{attention}

\courssub{Exemples d'usage au Cameroun}

Les moteurs à combustion interne font littéralement bouger le pays. En voici les applications les plus visibles :

\begin{itemize}
\item \textbf{Transports :} motos-taxis « benskins » (moteurs à essence), voitures et taxis, camions et bus de voyage (souvent diesel) ;
\item \textbf{Production d'électricité :} les \textbf{groupes électrogènes} des quartiers, des boutiques et des hôpitaux, précieux pendant les délestages ;
\item \textbf{Transport fluvial et pêche :} les \textbf{pirogues motorisées} du Wouri, du fleuve Sanaga et des côtes de Kribi, équipées de moteurs hors-bord ;
\item \textbf{Agriculture :} les \textbf{motopompes} qui puisent l'eau pour irriguer les champs, les motoculteurs, les décortiqueuses ;
\item \textbf{Artisanat :} les moulins à maïs et à manioc des villages, entraînés par de petits moteurs diesel ou essence.
\end{itemize}

\begin{exemplebox}[Le groupe électrogène du quartier]
Pendant une coupure ENEO, le menuisier du quartier démarre son groupe électrogène. Que se passe-t-il du point de vue énergétique ?
\begin{itemize}
\item le \textbf{moteur à essence} du groupe transforme l'énergie chimique de l'essence en énergie mécanique (rotation de l'arbre) ;
\item l'\textbf{alternateur}, entraîné par le moteur, transforme cette énergie mécanique en énergie électrique ;
\item les machines de l'atelier (scie, perceuse) transforment à leur tour l'énergie électrique en énergie mécanique.
\end{itemize}
Chaîne complète : énergie chimique $\rightarrow$ énergie mécanique $\rightarrow$ énergie électrique $\rightarrow$ énergie mécanique. Le groupe électrogène illustre à merveille le rôle du moteur à combustion interne : c'est lui la « source de mouvement » de toute l'installation.
\end{exemplebox}

\begin{savaistu}[Les pionniers du moteur à combustion interne]
Le premier moteur à combustion interne réellement fonctionnel a été construit en \textbf{1860} par l'ingénieur belge \textbf{Étienne Lenoir} : il brûlait du gaz d'éclairage et ressemblait encore à une machine à vapeur. Son rendement était faible, mais il prouvait que l'idée fonctionnait. En \textbf{1876}, l'Allemand \textbf{Nikolaus Otto} construisit le premier moteur à \textbf{quatre temps} (admission, compression, combustion-détente, échappement), beaucoup plus efficace : ce cycle, étudié dans la suite de la leçon, équipe encore aujourd'hui la quasi-totalité des voitures et des motos du Cameroun. Plus tard, \textbf{Rudolf Diesel} (1893) inventa le moteur qui porte son nom, sans bougie, où le carburant s'enflamme par la seule chaleur de la compression. En un peu plus de trente ans, ces inventions ont changé le monde : elles ont rendu possibles l'automobile, la moto, le camion\ldots{} et le benskin de Jean !
\end{savaistu}

\courssub{Exercice résolu}

\begin{exoresolu}[Identifier et compléter une chaîne énergétique]
\textbf{Énoncé.} Une motopompe à essence irrigue un champ de tomates près de Maroua. Le moteur entraîne une pompe qui monte l'eau du forage vers le champ.
\begin{enumerate}
\item De quel type de moteur s'agit-il ? Justifier.
\item Compléter la chaîne énergétique : énergie \ldots\ du carburant $\rightarrow$ énergie \ldots\ (combustion) $\rightarrow$ énergie \ldots\ (mouvement).
\item Un élève affirme : « le moteur de la motopompe fabrique l'énergie qui fait monter l'eau ». Corriger cette phrase.
\end{enumerate}
\tcblower
\textbf{Solution détaillée.}
\begin{enumerate}
\item Il s'agit d'un \textbf{moteur à combustion interne}, car la combustion de l'essence se produit \textbf{à l'intérieur} du moteur, dans la chambre de combustion, et les gaz chauds poussent directement le piston.
\item Énergie \textbf{chimique} du carburant $\rightarrow$ énergie \textbf{thermique} (combustion) $\rightarrow$ énergie \textbf{mécanique} (mouvement du piston, puis rotation du vilebrequin qui entraîne la pompe).
\item Phrase corrigée : « le moteur de la motopompe \textbf{transforme} l'énergie chimique de l'essence en énergie mécanique, qui permet à la pompe de faire monter l'eau ». Un moteur ne fabrique jamais d'énergie : il la convertit. Sans carburant, aucune énergie mécanique ne serait disponible.
\end{enumerate}
\end{exoresolu}

\begin{exercice}[Pour t'entraîner]
\begin{enumerate}
\item Définir un moteur à combustion interne en une phrase complète.
\item Pourquoi la machine à vapeur est-elle appelée moteur à combustion \emph{externe} ?
\item Dans l'expérience du canon à piston, citer deux observations qui prouvent qu'une transformation d'énergie a eu lieu.
\item Classer les appareils suivants selon qu'ils contiennent ou non un moteur à combustion interne : ventilateur électrique ; pirogue à moteur hors-bord ; groupe électrogène à gasoil ; vélo ; camion-citerne.
\item Reproduire et légender le schéma de la chaîne énergétique d'une moto, de l'essence aux roues.
\end{enumerate}
\end{exercice}

\begin{corrige}[Exercice « Pour t'entraîner »]
\begin{enumerate}
\item Un moteur à combustion interne est un moteur dans lequel la combustion du carburant se produit à l'intérieur même du moteur, dans la chambre de combustion, et dont les gaz chauds sous pression poussent directement le piston pour produire le mouvement.
\item Parce que le combustible (bois, charbon) brûle à l'extérieur du moteur, sous une chaudière ; la chaleur produite est transmise à l'eau, et c'est la vapeur d'eau qui pousse le piston. La combustion et le moteur sont séparés.
\item On observe une détonation et le piston violemment chassé du tube (apparition d'un mouvement : énergie mécanique), ainsi que l'échauffement du tube (énergie thermique). L'énergie chimique de l'alcool s'est transformée en énergie thermique puis mécanique.
\item Avec MCI : pirogue à moteur hors-bord, groupe électrogène à gasoil, camion-citerne. Sans MCI : ventilateur électrique (moteur électrique), vélo (énergie musculaire).
\item Carburant (essence) $\rightarrow$ combustion dans la chambre de combustion $\rightarrow$ gaz chauds sous pression $\rightarrow$ piston (va-et-vient) $\rightarrow$ vilebrequin (rotation) $\rightarrow$ roues. Formes d'énergie : chimique $\rightarrow$ thermique $\rightarrow$ mécanique.
\end{enumerate}
\end{corrige}

\begin{aretenir}[Bilan de la section]
\begin{itemize}
\item Un \textbf{moteur} transforme une forme d'énergie en \textbf{énergie mécanique} ; il ne fabrique jamais d'énergie.
\item Dans un \textbf{moteur à combustion interne}, le carburant brûle \emph{à l'intérieur} du moteur (chambre de combustion) ; les gaz chauds sous pression poussent directement le piston.
\item À l'inverse, dans un moteur à \textbf{combustion externe} (machine à vapeur), le feu brûle hors du moteur et chauffe un fluide intermédiaire (la vapeur d'eau).
\item Chaîne énergétique : énergie \textbf{chimique} $\rightarrow$ énergie \textbf{thermique} $\rightarrow$ énergie \textbf{mécanique} (translation du piston, puis rotation du vilebrequin).
\item Applications au Cameroun : motos, voitures, camions, groupes électrogènes, pirogues motorisées, motopompes.
\end{itemize}
Dans la prochaine section, nous ouvrirons le capot : nous découvrirons la \textbf{constitution} d'un moteur à combustion interne, pièce par pièce.
\end{aretenir}

\coursec{Constitution d'un moteur à combustion interne}

Dans la section précédente, nous avons défini le moteur à combustion interne : c'est une machine qui transforme l'énergie chimique d'un carburant (essence, gasoil) en énergie mécanique, grâce à une combustion qui se produit \emph{à l'intérieur} même de la machine. Nous avons vu que cette énergie fait tourner les roues d'une moto-taxi ou l'arbre d'une pompe à eau. Mais \emph{comment} cette transformation est-elle réalisée concrètement ? Quelles pièces entrent en jeu ? Ouvrons maintenant le capot et observons l'intérieur du moteur, pièce par pièce, comme le ferait un mécanicien de quartier à Douala ou à Yaoundé.

\courssub{Les organes fixes du moteur}

\textbf{Observation préalable.} Quand un mécanicien démonte le moteur d'une moto, il sépare d'abord les pièces qui \emph{ne bougent pas} pendant le fonctionnement : elles sont vissées au bloc moteur. Ce sont les \cle{organes fixes}. Puis il extrait les pièces qui se déplacent : ce sont les \cle{organes mobiles}. Cette distinction simple va nous guider.

\begin{definition}[Les organes fixes]
Les organes fixes d'un moteur à combustion interne sont les pièces qui restent immobiles pendant le fonctionnement. Ce sont principalement :
\begin{itemize}
\item le \cle{cylindre} : tube creux, aux parois très lisses, dans lequel coulisse le piston ; il est souvent réalisé en fonte ou en aluminium traité ;
\item la \cle{culasse} : pièce qui ferme le cylindre par le haut ; elle porte les soupapes, la bougie (moteur essence) ou l'injecteur (moteur diesel) ;
\item le \cle{carter} (ou \emph{bas-moteur}) : carter inférieur qui ferme le bas du moteur ; il contient l'huile de graissage et loge le vilebrequin.
\end{itemize}
\end{definition}

Le cylindre est le « cœur » du moteur : c'est en son sein que tout se passe. Ses parois doivent être parfaitement polies et résister à de très hautes températures, car la combustion y produit des gaz portés à plus de \SI{2000}{\degreeCelsius}. La culasse, boulonnée sur le cylindre, joue le rôle de « couvercle » : sans elle, les gaz brûlés s'échapperaient sans pousser le piston. Le carter, enfin, est le « carter d'huile » : il protège les organes mobiles et assure leur lubrification.

\begin{attention}[Ne pas confondre culasse et carter]
La culasse n'est pas le carter ! La \textbf{culasse} ferme le cylindre \emph{par le haut} et porte les soupapes et la bougie (ou l'injecteur). Le \textbf{carter} ferme le moteur \emph{par le bas} et contient l'huile. Au BEPC, cette inversion est une erreur fréquente.
\end{attention}

\courssub{Les organes mobiles : piston, bielle, vilebrequin}

\begin{definition}[Les organes mobiles]
Les organes mobiles sont les pièces qui se déplacent pendant le fonctionnement du moteur :
\begin{itemize}
\item le \cle{piston} : pièce cylindrique qui coulisse (monte et descend) dans le cylindre ; il reçoit la poussée des gaz brûlés. Il est muni de \emph{segments} (anneaux élastiques) qui assurent l'étanchéité entre le piston et la paroi du cylindre ;
\item la \cle{bielle} : tige articulée qui relie le piston au vilebrequin ;
\item le \cle{vilebrequin} : arbre coudé (une sorte de manivelle) qui transforme le mouvement du piston en mouvement de rotation ; c'est lui qui entraîne, à travers la boîte de vitesses, les roues du véhicule.
\end{itemize}
\end{definition}

Le piston est comme un « bouchon coulissant » : il est ajusté au cylindre au dixième de millimètre près, grâce à ses segments. Quand les gaz brûlent au-dessus de lui, ils le poussent violemment vers le bas : c'est là que naît la force motrice.

\begin{attention}[Erreur fréquente : inverser bielle et vilebrequin]
Beaucoup d'élèves confondent ces deux pièces. Retiens bien : la \textbf{bielle} est la tige qui \emph{relie le piston au vilebrequin} ; le \textbf{vilebrequin} est l'arbre coudé qui \emph{tourne}. Astuce de mémorisation : « \textbf{bi}elle » commence comme « \textbf{bi}pied » qui relie deux points, et le « vilebrequin » est celui qui « \textbf{vibre} en tournant ». Sur un schéma, la bielle est la tige droite, le vilebrequin est la manivelle coudée en bas.
\end{attention}

\courssub{Le système bielle-manivelle : transformer la translation en rotation}

\textbf{Situation concrète.} Regarde un cycliste : ses jambes montent et descendent (mouvement rectiligne alternatif), pourtant la roue du vélo tourne en continu. Le secret, ce sont les pédales et la manivelle du pédalier. Le moteur utilise exactement le même principe !

\begin{experience}[Observer la transformation de mouvement]
\textbf{Matériel :} un vieux moteur de moto démonté (ou une maquette), ou simplement une manivelle de moulin à café.
\textbf{Protocole :} fais tourner lentement le vilebrequin à la main et observe le piston. Puis pousse le piston vers le bas et observe le vilebrequin.
\textbf{Observations :} quand le vilebrequin tourne, le piston monte et descend alternativement. Réciproquement, quand on pousse le piston, le vilebrequin tourne.
\textbf{Interprétation :} l'ensemble bielle + vilebrequin constitue un système \cle{bielle-manivelle} qui transforme un mouvement rectiligne alternatif (va-et-vient du piston) en mouvement de rotation continue (du vilebrequin), et inversement.
\end{experience}

\begin{propriete}[Rôle du système bielle-manivelle]
Le système bielle-manivelle transforme le mouvement \textbf{rectiligne alternatif} du piston en mouvement de \textbf{rotation continue} du vilebrequin (temps moteur), et assure la transformation inverse (rotation vers translation) lors des temps d'admission, compression et échappement. C'est ce mouvement de rotation qui est transmis aux roues du véhicule.
\end{propriete}

\begin{popfigure}
\begin{center}
\begin{tikzpicture}[scale=0.9,>=Stealth]
% Cylindre
\draw[thick,popdark] (-1.2,0) -- (-1.2,4.6);
\draw[thick,popdark] (1.2,0) -- (1.2,4.6);
% Culasse
\draw[thick,popdark,fill=popblueL] (-1.6,4.6) rectangle (1.6,5.2);
\node[popdark] at (0,5.55) {\small \textbf{Culasse}};
% Chambre de combustion
\draw[fill=popgoldL] (-1.2,3.9) rectangle (1.2,4.6);
\node[popdark] at (0,4.25) {\small chambre de combustion};
% Soupapes
\draw[thick,popgreen] (-0.7,4.6) -- (-0.7,5.2);
\draw[thick,popgreen,fill=popgreen] (-0.95,4.6) -- (-0.45,4.6) -- (-0.7,4.35) -- cycle;
\draw[thick,popgreen] (0.7,4.6) -- (0.7,5.2);
\draw[thick,popgreen,fill=popgreen] (0.95,4.6) -- (0.45,4.6) -- (0.7,4.35) -- cycle;
\node[popgreen,anchor=west] at (1.7,5.0) {\small soupapes};
\draw[->,popgreen] (1.65,5.0) -- (0.85,4.7);
% Bougie
\draw[thick,poporange] (0,5.2) -- (0,5.9);
\draw[thick,poporange] (-0.12,4.6) rectangle (0.12,5.2);
\draw[poporange] (-0.08,4.6) -- (0,4.45) -- (0.08,4.6);
\node[poporange,anchor=east] at (-1.7,5.5) {\small bougie};
\draw[->,poporange] (-1.65,5.5) -- (-0.15,5.4);
% Piston
\draw[thick,popdark,fill=popblue!40] (-1.15,2.9) rectangle (1.15,3.9);
\draw[popdark] (-1.15,3.7) -- (1.15,3.7);
\draw[popdark] (-1.15,3.55) -- (1.15,3.55);
\node[popdark,anchor=west] at (1.7,3.4) {\small \textbf{piston}};
\draw[->,popdark] (1.65,3.4) -- (1.2,3.4);
% Axe de piston
\fill[popdark] (0,3.4) circle (0.09);
% Bielle
\draw[very thick,poppink] (0,3.4) -- (0.55,1.1);
\node[poppink,anchor=west] at (1.7,2.2) {\small \textbf{bielle}};
\draw[->,poppink] (1.65,2.2) -- (0.35,2.3);
% Vilebrequin
\draw[very thick,poppurple] (0.55,1.1) -- (0,0.55);
\draw[very thick,poppurple] (0,0.55) -- (-0.9,0.55);
\fill[poppurple] (0,0.55) circle (0.12);
\fill[poppink] (0.55,1.1) circle (0.09);
\node[poppurple,anchor=west] at (1.7,0.55) {\small \textbf{vilebrequin}};
\draw[->,poppurple] (1.65,0.55) -- (0.15,0.55);
% Rotation
\draw[->,poppurple,thick] (0.9,0.2) arc (-40:80:0.45);
\node[poppurple] at (1.55,-0.15) {\small rotation};
% Carter
\draw[thick,popdark,dashed] (-1.6,0) rectangle (1.6,-0.4);
\node[popdark] at (0,-0.2) {\small carter (huile)};
% PMH / PMB
\draw[<->,popblue,thick] (-2.1,3.9) -- (-2.1,2.9);
\node[popblue,rotate=90] at (-2.4,3.4) {\small course};
\draw[dashed,popblue] (-1.2,3.9) -- (-2.3,3.9);
\draw[dashed,popblue] (-1.2,2.9) -- (-2.3,2.9);
\node[popblue,anchor=east] at (-2.35,3.9) {\small PMH};
\node[popblue,anchor=east] at (-2.35,2.9) {\small PMB};
\end{tikzpicture}
\end{center}
\legende{Schéma en coupe d'un cylindre de moteur à essence. Le piston coulisse dans le cylindre ; la bielle le relie au vilebrequin. La culasse porte les soupapes et la bougie. PMH : point mort haut ; PMB : point mort bas.}
\end{popfigure}

\begin{popfigure}
\begin{center}
\begin{tikzpicture}[scale=1.0,>=Stealth]
% Manivelle
\draw[thick,popmuted] (0,0) circle (1.5);
\fill[popdark] (0,0) circle (0.08);
\node[popdark,below left] at (-0.1,-0.1) {\small vilebrequin};
\fill[poporange] (60:1.5) circle (0.1);
\draw[very thick,poporange] (0,0) -- (60:1.5);
\node[poporange] at (60:1.9) {\small maneton};
% Bielle
\draw[very thick,poppink] (60:1.5) -- ++(0,2.2) coordinate(pied);
\node[poppink,anchor=west] at ($(60:1.5)+(0.15,1.1)$) {\small bielle};
% Piston
\draw[thick,popdark,fill=popblue!40] ($(pied)+(-0.7,0)$) rectangle ($(pied)+(0.7,0.8)$);
\node[popdark,anchor=west] at ($(pied)+(0.8,0.4)$) {\small piston};
% Cylindre
\draw[thick,popdark] ($(pied)+(-0.85,-0.3)$) -- ($(pied)+(-0.85,1.6)$);
\draw[thick,popdark] ($(pied)+(0.85,-0.3)$) -- ($(pied)+(0.85,1.6)$);
% Flèches
\draw[<->,very thick,popblue] ($(pied)+(-1.4,0)$) -- ($(pied)+(-1.4,1.2)$);
\node[popblue,anchor=east,align=right] at ($(pied)+(-1.5,0.6)$) {\small translation\\ \small alternée};
\draw[->,very thick,poppurple] (-1.9,-0.4) arc (200:340:1.1);
\node[poppurple] at (0,-1.9) {\small rotation continue};
% Équivalence
\node[popdark,align=center] at (5.6,1.4) {\small \textbf{translation du piston} $\Longleftrightarrow$ \textbf{rotation du vilebrequin}};
\end{tikzpicture}
\end{center}
\legende{Le système bielle-manivelle : le va-et-vient du piston (flèche bleue) est transformé en rotation continue du vilebrequin (flèche violette), comme les jambes du cycliste font tourner le pédalier.}
\end{popfigure}

\courssub{Les organes d'admission, d'échappement, d'allumage et d'injection}

Pour que la combustion ait lieu, il faut faire entrer le mélange air-carburant, l'enflammer, puis évacuer les gaz brûlés. D'où trois familles d'organes supplémentaires.

\begin{definition}[Admission, échappement, allumage, injection]
\begin{itemize}
\item Les organes d'\cle{admission} laissent entrer le mélange air-carburant (essence) ou l'air seul (diesel) dans le cylindre ; les organes d'\cle{échappement} laissent sortir les gaz brûlés. Dans un moteur à \textbf{4 temps}, ce sont des \cle{soupapes} (clapets commandés par un arbre à cames) ; dans un moteur à \textbf{2 temps}, ce sont de simples ouvertures appelées \cle{lumières}, découvertes et fermées par le piston lui-même.
\item L'organe d'\cle{allumage} d'un moteur à essence est la \cle{bougie} : elle produit une étincelle électrique qui enflamme le mélange compressé.
\item Dans un moteur diesel, il n'y a pas de bougie : c'est l'\cle{injecteur} qui pulvérise le gasoil finement dans l'air fortement comprimé et très chaud ; le mélange s'enflamme alors spontanément (auto-allumage par compression).
\end{itemize}
\end{definition}

\begin{methode}[Identifier les organes d'un moteur sur un schéma en coupe]
Pour reconnaître les organes sur un schéma, suis le \textbf{trajet des gaz} :
\begin{enumerate}
\item repère l'\textbf{admission} : par où le mélange air-carburant (ou l'air) entre (soupape ou lumière d'admission) ;
\item suis le gaz jusqu'à la \textbf{chambre de combustion} : au-dessus du piston, sous la culasse ; c'est là que se trouvent la bougie (essence) ou l'injecteur (diesel) ;
\item repère l'\textbf{échappement} : par où les gaz brûlés sortent (soupape ou lumière d'échappement) ;
\item identifie enfin les organes mobiles dans l'ordre, de haut en bas : \textbf{piston} $\rightarrow$ \textbf{bielle} $\rightarrow$ \textbf{vilebrequin}.
\end{enumerate}
En suivant ce trajet logique (admission $\rightarrow$ combustion $\rightarrow$ échappement), tu ne peux plus te tromper !
\end{methode}

\courssub{La chambre de combustion et les repères du piston}

\begin{definition}[Chambre de combustion]
La \cle{chambre de combustion} est l'espace compris entre le piston lorsqu'il est au \textbf{point mort haut (PMH)} et la culasse. C'est dans cet espace réduit que le mélange air-carburant est comprimé puis enflammé.
\end{definition}

Pour décrire le mouvement du piston, les mécaniciens utilisent des repères précis :

\begin{definition}[PMH, PMB, course, cylindrée]
\begin{itemize}
\item le \cle{point mort haut (PMH)} : position la plus haute atteinte par le piston dans le cylindre (volume minimal) ;
\item le \cle{point mort bas (PMB)} : position la plus basse atteinte par le piston (volume maximal) ;
\item la \cle{course} : distance parcourue par le piston entre le PMH et le PMB ;
\item la \cle{cylindrée} (ou \emph{cylindrée unitaire}) : volume balayé par le piston entre le PMB et le PMH. Elle s'exprime en centimètres cubes (\si{cm^3}) ou en litres (\si{L}).
\end{itemize}
\end{definition}

\begin{aretenir}[Cylindrée totale]
La cylindrée totale est le volume total balayé par l'ensemble des pistons. Pour un moteur à plusieurs cylindres, on multiplie la cylindrée unitaire par le nombre de cylindres :
\[ \text{Cylindrée totale} = \text{cylindrée unitaire} \times \text{nombre de cylindres}. \]
Unités usuelles : \si{cm^3} (motos, petites cylindrées) ou \si{L} (voitures).
\end{aretenir}

\begin{exemplebox}[La moto « 125 » et la voiture « 2,0 L »]
Quand on dit qu'une moto-taxi est une « 125 », cela signifie que sa \textbf{cylindrée} vaut \SI{125}{cm^3} : le piston balaie un volume de \SI{125}{cm^3} entre le PMB et le PMH. De même, une voiture « 2,0 L » possède une cylindrée totale de \SI{2000}{cm^3}, répartie en général sur 4 cylindres de \SI{500}{cm^3} chacun.
\end{exemplebox}

\begin{exoresolu}[Calcul d'une cylindrée totale]
Énoncé : le moteur d'une voiture possède 4 cylindres identiques. Dans chaque cylindre, le piston balaie un volume de \SI{400}{cm^3} entre le PMB et le PMH. Calcule la cylindrée totale de ce moteur, en \si{cm^3} puis en litres.
\tcblower
Solution détaillée :
\begin{align*}
\text{Cylindrée totale} &= \text{cylindrée unitaire} \times \text{nombre de cylindres} \\
&= 400 \times 4 \\
&= \SI{1600}{cm^3}
\end{align*}
Or $\SI{1}{L} = \SI{1000}{cm^3}$, donc :
\[ \SI{1600}{cm^3} = \frac{1600}{1000} = \SI{1,6}{L}. \]
Conclusion : la cylindrée totale du moteur est de \SI{1600}{cm^3}, soit \SI{1,6}{L}. On dira que c'est une voiture « 1,6 L ».
\end{exoresolu}

\begin{savaistu}[Un vilebrequin infatigable]
Le vilebrequin d'un moteur de voiture peut tourner à plus de \num{6000} tours par minute, soit 100 tours par seconde ! À chaque tour, la bielle pousse et tire le piston. Cela signifie que le piston monte et descend des dizaines de fois par seconde, tout en supportant des températures et des pressions énormes. C'est pourquoi l'huile contenue dans le carter est indispensable : sans graissage, ces pièces en métal s'useraient et se gripperaient en quelques minutes. Les moteurs des motos de course dépassent même \num{15000} tours par minute !
\end{savaistu}

\begin{exercice}[Pour t'entraîner]
\begin{enumerate}
\item Nomme les trois organes fixes d'un moteur à combustion interne et donne le rôle de chacun.
\item Sur un schéma en coupe, un élève affirme : « la bielle est l'arbre coudé qui tourne ». Corrige son erreur.
\item Une moto possède un moteur de cylindrée \SI{150}{cm^3}. Que signifie cette valeur ?
\item Un moteur de voiture a 4 cylindres de \SI{375}{cm^3} chacun. Calcule sa cylindrée totale en litres.
\end{enumerate}
\end{exercice}

\begin{corrige}[Exercice « Pour t'entraîner »]
\begin{enumerate}
\item Le \textbf{cylindre} (guide le piston), la \textbf{culasse} (ferme le cylindre par le haut et porte soupapes et bougie/injecteur) et le \textbf{carter} (ferme le bas du moteur et contient l'huile de graissage).
\item C'est faux : l'arbre coudé qui tourne est le \textbf{vilebrequin}. La \textbf{bielle} est la tige articulée qui relie le piston au vilebrequin.
\item Le piston balaie un volume de \SI{150}{cm^3} entre le PMB et le PMH (cylindrée unitaire).
\item Cylindrée totale $= 375 \times 4 = \SI{1500}{cm^3} = \SI{1,5}{L}$.
\end{enumerate}
\end{corrige}

Nous connaissons désormais toutes les pièces du moteur et leurs rôles. Reste à comprendre comment elles travaillent ensemble dans le temps : c'est l'objet de la section suivante, consacrée au cycle à quatre temps du moteur à essence.

\coursec{Le cycle à 4 temps du moteur à essence}

Dans la section précédente, tu as découvert la constitution d'un moteur à combustion interne : le cylindre, le piston, la bielle, le vilebrequin, les soupapes, la bougie et le volant moteur. Tu sais donc \emph{de quoi} le moteur est fait. Reste la grande question : \textbf{comment toutes ces pièces travaillent-elles ensemble pour faire avancer la moto-taxi qui t'emmène au collège ?} La réponse tient en une phrase : le moteur répète sans cesse une suite de quatre opérations bien précises, appelée \cle{cycle à 4 temps}. C'est ce cycle que nous allons étudier pas à pas.

\courssub{Principe général du cycle à 4 temps}

\begin{definition}[Cycle à 4 temps]
Dans un moteur à essence à 4 temps, un \cle{cycle} complet comporte \textbf{quatre courses du piston} (deux descentes et deux montées), appelées \cle{temps} : l'\cle{admission}, la \cle{compression}, la \cle{combustion-détente} et l'\cle{échappement}. Comme une course du piston correspond à un demi-tour de vilebrequin, un cycle complet correspond à \textbf{deux tours de vilebrequin}.
\end{definition}

Pour bien comprendre, imagine une seringue dont tu bouches l'extrémité avec le doigt : quand tu tires le piston, l'air entre ; quand tu pousses, l'air est comprimé. Le moteur fait la même chose, mais de façon automatique et très rapide, avec un mélange d'air et d'essence qui peut brûler. Chaque « temps » porte un nom qui décrit exactement ce qui se passe dans le cylindre.

\begin{aretenir}[L'essentiel sur le cycle]
\begin{itemize}
\item 1 cycle = \textbf{4 courses du piston} = \textbf{2 tours de vilebrequin}.
\item Un seul temps sur quatre produit du travail : c'est la \textbf{combustion-détente}, appelée \cle{temps moteur}.
\item Les trois autres temps sont \emph{résistants} : ils consomment de l'énergie, fournie par l'inertie du \cle{volant moteur}.
\end{itemize}
\end{aretenir}

\courssub{Étude détaillée des quatre temps}

\textbf{1\textsuperscript{er} temps : l'admission.}\par
La \textbf{soupape d'admission s'ouvre}, la soupape d'échappement reste fermée. Le piston \textbf{descend} du point mort haut (PMH) vers le point mort bas (PMB). En descendant, il crée une dépression dans le cylindre, exactement comme quand tu tires le piston d'une seringue : le \textbf{mélange air-essence} préparé par le carburateur (ou l'injection) est \emph{aspiré} dans le cylindre.

\textbf{2\textsuperscript{e} temps : la compression.}\par
Les \textbf{deux soupapes sont fermées}. Le piston \textbf{remonte} et comprime le mélange air-essence dans un volume de plus en plus petit (la chambre de combustion). Conséquence physique importante : la \textbf{pression} et la \textbf{température} du mélange augmentent fortement. Un gaz comprimé chauffe : c'est pourquoi le corps d'une pompe à vélo devient chaud quand on gonfle un pneu. Ce mélange chaud et comprimé est prêt à brûler très vite.

\textbf{3\textsuperscript{e} temps : la combustion-détente (temps moteur).}\par
Les deux soupapes sont toujours fermées. Au moment où le piston arrive en haut, la \textbf{bougie produit une étincelle} électrique entre ses électrodes. Le mélange air-essence s'enflamme brutalement : c'est une combustion très rapide, assimilée à une petite explosion. Les gaz brûlés, très chauds et sous très forte pression, \textbf{poussent violemment le piston vers le bas}. Cette poussée, transmise par la bielle au vilebrequin, fait tourner l'arbre moteur.

\begin{definition}[Temps moteur]
Le \cle{temps moteur} est le seul temps du cycle pendant lequel les gaz poussent le piston et \textbf{produisent du travail} mécanique. Dans le moteur à essence, c'est le temps de \textbf{combustion-détente} (3\textsuperscript{e} temps). Les trois autres temps (admission, compression, échappement) ne produisent aucun travail : ils sont réalisés grâce à l'énergie emmagasinée par le volant moteur.
\end{definition}

\textbf{4\textsuperscript{e} temps : l'échappement.}\par
La \textbf{soupape d'échappement s'ouvre}, la soupape d'admission reste fermée. Le piston \textbf{remonte} et \textbf{chasse les gaz brûlés} hors du cylindre, vers le pot d'échappement. Le cylindre est ainsi nettoyé, et un nouveau cycle peut recommencer par une admission.

\begin{popfigure}
\begin{tikzpicture}[scale=0.85, every node/.style={font=\small}]
% ===== Temps 1 : Admission =====
\begin{scope}[shift={(0,0)}]
\draw[thick, popdark] (0,0) rectangle (2.2,3);
\draw[fill=popblueL, draw=popdark] (0.15,0.6) rectangle (2.05,1.3);
\draw[thick, popdark] (1.1,0.6) -- (1.1,-0.5);
\draw[thick, popdark] (0.7,-0.5) -- (1.5,-0.5);
% soupape admission ouverte
\draw[thick, popgreen] (0.55,3) -- (0.55,2.45);
\draw[thick, popgreen] (0.3,2.45) -- (0.8,2.45);
% soupape echappement fermée
\draw[thick, popink] (1.65,3) -- (1.65,3.0);
\draw[thick, popink] (1.4,3) -- (1.9,3);
% flèche mélange
\draw[-{Stealth}, very thick, popgreen] (0.55,3.8) -- (0.55,2.7);
% flèche piston descend
\draw[-{Stealth}, very thick, popblue] (2.5,1.6) -- (2.5,0.7);
\node[popdark, font=\bfseries] at (1.1,-1) {1. Admission};
\node[popgreen, font=\scriptsize] at (1.1,4.1) {mélange air-essence};
\end{scope}
% ===== Temps 2 : Compression =====
\begin{scope}[shift={(4.6,0)}]
\draw[thick, popdark] (0,0) rectangle (2.2,3);
\draw[fill=popblueL, draw=popdark] (0.15,1.9) rectangle (2.05,2.6);
\draw[thick, popdark] (1.1,1.9) -- (1.1,-0.5);
\draw[thick, popdark] (0.7,-0.5) -- (1.5,-0.5);
\draw[thick, popink] (0.55,3) -- (0.55,3);
\draw[thick, popink] (0.3,3) -- (0.8,3);
\draw[thick, popink] (1.4,3) -- (1.9,3);
% petites flèches de compression
\draw[-{Stealth}, thick, poporange] (0.6,3.35) -- (0.6,2.95);
\draw[-{Stealth}, thick, poporange] (1.1,3.35) -- (1.1,2.95);
\draw[-{Stealth}, thick, poporange] (1.6,3.35) -- (1.6,2.95);
\draw[-{Stealth}, very thick, popblue] (2.5,1.0) -- (2.5,2.0);
\node[popdark, font=\bfseries] at (1.1,-1) {2. Compression};
\node[popmuted, font=\scriptsize] at (1.1,4.1) {soupapes fermées};
\end{scope}
% ===== Temps 3 : Combustion-détente =====
\begin{scope}[shift={(9.2,0)}]
\draw[thick, popdark] (0,0) rectangle (2.2,3);
\draw[fill=popblueL, draw=popdark] (0.15,1.9) rectangle (2.05,2.6);
\draw[thick, popdark] (1.1,1.9) -- (1.1,-0.5);
\draw[thick, popdark] (0.7,-0.5) -- (1.5,-0.5);
\draw[thick, popink] (0.3,3) -- (0.8,3);
\draw[thick, popink] (1.4,3) -- (1.9,3);
% bougie + étincelle
\draw[thick, poppurple] (1.1,3.5) -- (1.1,3.05);
\draw[thick, popgold] (1.0,3.05) -- (1.2,2.85);
\draw[thick, popgold] (1.2,3.05) -- (1.0,2.85);
% explosion
\node[star, star points=7, fill=poporangeL, draw=poporange, minimum size=0.7cm, inner sep=0pt] at (1.1,2.85) {};
\draw[-{Stealth}, very thick, poporange] (2.5,2.2) -- (2.5,1.1);
\node[popdark, font=\bfseries] at (1.1,-1) {3. Combustion-détente};
\node[popgold, font=\scriptsize] at (1.1,4.1) {étincelle de la bougie};
\end{scope}
% ===== Temps 4 : Échappement =====
\begin{scope}[shift={(13.8,0)}]
\draw[thick, popdark] (0,0) rectangle (2.2,3);
\draw[fill=popblueL, draw=popdark] (0.15,1.9) rectangle (2.05,2.6);
\draw[thick, popdark] (1.1,1.9) -- (1.1,-0.5);
\draw[thick, popdark] (0.7,-0.5) -- (1.5,-0.5);
\draw[thick, popink] (0.3,3) -- (0.8,3);
% soupape échappement ouverte
\draw[thick, poppink] (1.65,3) -- (1.65,2.45);
\draw[thick, poppink] (1.4,2.45) -- (1.9,2.45);
\draw[-{Stealth}, very thick, poppink] (1.65,2.7) -- (1.65,3.8);
\draw[-{Stealth}, very thick, popblue] (2.5,1.0) -- (2.5,2.0);
\node[popdark, font=\bfseries] at (1.1,-1) {4. Échappement};
\node[poppink, font=\scriptsize] at (1.1,4.1) {gaz brûlés};
\end{scope}
\end{tikzpicture}
\legende{Les quatre temps du moteur à essence. En bleu clair : le piston ; les flèches bleues indiquent son sens de déplacement. En vert : la soupape d'admission ouverte ; en rose : la soupape d'échappement ouverte ; en noir : soupape fermée. Au 3\textsuperscript{e} temps, la bougie (violet) produit l'étincelle (jaune) qui enflamme le mélange (étoile orange) : c'est le seul temps moteur.}
\end{popfigure}

\courssub{Observation expérimentale : un seul temps sur quatre est moteur}

\begin{experience}[Observer le cycle sur une maquette de moteur]
\textbf{Matériel :} maquette de moteur à 4 temps à carter transparent (ou moteur réel décapé et tourné à la main), manivelle.
\begin{enumerate}
\item On tourne lentement la manivelle et on observe le piston, les soupapes et (sur la maquette) le moment de l'étincelle simulée par une petite lampe.
\item On note, pour chaque demi-tour de vilebrequin : le sens de déplacement du piston et l'état des soupapes.
\end{enumerate}
\textbf{Observations :} sur quatre courses successives du piston, une seule correspond à une poussée des gaz vers le bas (combustion-détente). Pendant les trois autres courses, il faut \emph{fournir} un effort sur la manivelle : on le sent nettement à la compression, qui est dure à tourner.
\textbf{Interprétation :} un cycle comporte 4 courses, soit 2 tours de vilebrequin, et \textbf{une seule course sur quatre est motrice}. Les trois autres sont \emph{entretenues} grâce à l'énergie accumulée par le volant moteur.
\textbf{Conclusion :} le moteur à 4 temps ne produit du travail que pendant un quart du cycle ; il a donc besoin d'un organe qui stocke l'énergie et la restitue entre deux temps moteurs.
\end{experience}

\begin{definition}[Rôle du volant moteur]
Le \cle{volant moteur} est un disque lourd fixé sur le vilebrequin. Pendant le temps moteur, il \textbf{emmagasine} une partie de l'énergie sous forme d'énergie de rotation (inertie). Pendant les trois temps résistants, il \textbf{restitue} cette énergie pour entraîner le piston. Il \textbf{régularise la rotation} du moteur entre deux temps moteurs : sans lui, le moteur tournerait par à-coups et s'arrêterait.
\end{definition}

\begin{methode}[Reconnaître chaque temps sur un schéma]
Pour identifier un temps du cycle sur un schéma, pose-toi deux questions dans l'ordre :
\begin{enumerate}
\item \textbf{Quel est le sens de déplacement du piston ?} (regarde la flèche ou compare deux positions successives) : descente ou montée.
\item \textbf{Quel est l'état des soupapes ?} (ouverte ou fermée, admission ou échappement).
\end{enumerate}
Puis conclus avec le tableau :
\begin{center}
\begin{tabularx}{\linewidth}{|l|X|X|X|}
\hline
\rowcolor{popblueL} \textbf{Temps} & \textbf{Piston} & \textbf{Soupapes} & \textbf{Indice} \\
\hline
Admission & descend & admission ouverte & mélange qui entre \\
\hline
Compression & monte & toutes fermées & pas d'étincelle \\
\hline
Combustion-détente & descend & toutes fermées & étincelle à la bougie \\
\hline
Échappement & monte & échappement ouverte & gaz brûlés qui sortent \\
\hline
\end{tabularx}
\end{center}
Astuce : si le piston descend, c'est admission ou détente ; l'étincelle (ou la flamme) désigne alors la détente. Si le piston monte, c'est compression ou échappement ; la soupape ouverte désigne l'échappement.
\end{methode}

\begin{attention}[Erreurs fréquentes]
\begin{itemize}
\item \textbf{La bougie ne fonctionne pas en continu !} Elle ne produit l'étincelle qu'\textbf{une fois par cycle}, au début du temps de combustion-détente. Dire « la bougie brûle l'essence en permanence » est une erreur classique au BEPC.
\item Ne confonds pas \emph{temps} (une course du piston) et \emph{tour} du vilebrequin : 4 temps = 2 tours, pas 4 tours.
\item La combustion n'est pas une explosion de la bougie : la bougie ne fait que \emph{déclencher} la combustion du mélange air-essence.
\end{itemize}
\end{attention}

\courssub{Application numérique : combien de temps moteurs par minute ?}

\begin{exoresolu}[Temps moteurs d'une moto à \SI{3000}{tr/min}]
\textbf{Énoncé.} Le moteur à 4 temps d'une moto-taxi (un seul cylindre) tourne à \SI{3000}{tr/min} (tours de vilebrequin par minute). 1) Combien de cycles complets le cylindre réalise-t-il par minute ? 2) Combien de temps moteurs par minute cela représente-t-il ? 3) Que devient ce nombre pour une voiture dont le moteur possède 4 cylindres tournant à la même vitesse ?
\tcblower
\textbf{Solution détaillée.}
\begin{enumerate}
\item Un cycle complet correspond à \textbf{2 tours de vilebrequin}. Le nombre de cycles par minute est donc :
\[ N_{\text{cycles}} = \frac{3000}{2} = 1500 \text{ cycles par minute.} \]
\item Chaque cycle comporte \textbf{un seul temps moteur}. Le cylindre réalise donc :
\[ N_{\text{moteurs}} = 1500 \text{ temps moteurs par minute.} \]
\item Avec 4 cylindres, chaque cylindre fournit 1500 temps moteurs par minute, donc au total :
\[ 4 \times 1500 = 6000 \text{ temps moteurs par minute.} \]
\end{enumerate}
\textbf{Conclusion :} à \SI{3000}{tr/min}, le moteur de la moto réalise \textbf{1500 temps moteurs par minute et par cylindre}, soit 25 par seconde ! On comprend pourquoi le volant moteur est indispensable : les « poussées » sont séparées les unes des autres. Avec 4 cylindres décalés, les temps moteurs se succèdent plus régulièrement : c'est pourquoi le moteur d'une voiture tourne plus « doux » que celui d'une moto.
\end{exoresolu}

\courssub{Complément (hors programme strict) : le diagramme pression-volume}

Pour les ingénieurs, on résume le cycle en traçant la \textbf{pression} des gaz dans le cylindre en fonction du \textbf{volume} de la chambre. Ce diagramme, appelé \emph{cycle théorique de Beau de Rochas -- Otto}, n'est pas exigible au BEPC, mais il t'aide à visualiser d'un coup d'œil les quatre temps.

\begin{popfigure}
\begin{tikzpicture}
\begin{axis}[
width=0.8\linewidth, height=6.5cm,
xlabel={Volume $V$ (arbitraire)},
ylabel={Pression $p$ (arbitraire)},
xmin=0.4, xmax=3.4, ymin=0, ymax=10,
xtick=\empty, ytick=\empty,
axis lines=left,
every axis x label/.style={at={(ticklabel* cs:1)}, anchor=west},
every axis y label/.style={at={(ticklabel* cs:1)}, anchor=south},
]
% admission (V augmente, p faible)
\addplot[popblue, very thick, domain=0.6:3, samples=2] {0.8};
% compression (V diminue, p monte) : courbe de 3 -> 0.6
\addplot[popgreen, very thick, domain=0.6:3, samples=60] {0.8*(3/x)^1.3};
% combustion (V constant, p monte)
\draw[poporange, very thick] (axis cs:0.6,4.5) -- (axis cs:0.6,9);
% détente (V augmente, p descend)
\addplot[popink, very thick, domain=0.6:3, samples=60] {9*(0.6/x)^1.3};
% échappement (retour)
\draw[poppink, very thick, dashed] (axis cs:3,1.6) -- (axis cs:3,0.8);
% annotations
\node[popblue, font=\scriptsize] at (axis cs:1.9,0.35) {1. Admission};
\node[popgreen, font=\scriptsize, rotate=-32] at (axis cs:1.75,2.6) {2. Compression};
\node[poporange, font=\scriptsize, anchor=west] at (axis cs:0.68,6.8) {3. Combustion};
\node[popink, font=\scriptsize, rotate=-30] at (axis cs:1.9,4.6) {3. Détente (temps moteur)};
\node[poppink, font=\scriptsize, anchor=west] at (axis cs:3.02,1.2) {4. Échapp.};
% flèches de sens
\draw[-{Stealth}, popblue] (axis cs:2.2,0.8) -- (axis cs:2.6,0.8);
\draw[-{Stealth}, popgreen] (axis cs:1.3,2.15) -- (axis cs:1.0,2.9);
\draw[-{Stealth}, popink] (axis cs:1.5,3.1) -- (axis cs:1.9,2.35);
\end{axis}
\end{tikzpicture}
\legende{Diagramme pression--volume simplifié du cycle à 4 temps (complément hors programme). L'admission se fait à basse pression (bleu), la compression fait monter la pression (vert), la combustion se produit à volume presque constant (orange), la détente repousse le piston (noir), puis les gaz brûlés sont chassés (rose). L'aire de la boucle représente le travail fourni par le moteur.}
\end{popfigure}

\begin{savaistu}[D'où vient le cycle à 4 temps ?]
Le cycle à 4 temps est aussi appelé \cle{cycle Beau de Rochas -- Otto}. En 1862, l'ingénieur français \textbf{Alphonse Beau de Rochas} dépose un brevet décrivant les quatre temps et les conditions d'un bon rendement. En 1876, l'ingénieur allemand \textbf{Nikolaus Otto} construit le premier moteur pratique fonctionnant sur ce principe : c'est la naissance du moteur à essence moderne. Près de 150 ans plus tard, le même cycle fait encore tourner les moto-taxis de Douala, les groupes électrogènes qui dépannent pendant les délestages et les pirogues à moteur hors-bord du fleuve Wouri !
\end{savaistu}

\begin{exercice}[À toi de jouer]
Sur le schéma d'un cylindre de moteur à essence, on observe : le piston qui \textbf{monte}, la soupape d'admission \textbf{fermée} et la soupape d'échappement \textbf{ouverte}.
\begin{enumerate}
\item De quel temps du cycle s'agit-il ? Justifie ta réponse.
\item Ce temps est-il moteur ou résistant ? Quel organe permet de le réaliser ?
\item Le vilebrequin tourne à \SI{2400}{tr/min}. Combien de temps moteurs le cylindre réalise-t-il en une minute ?
\end{enumerate}
\end{exercice}

\begin{corrige}[Exercice 1]
\begin{enumerate}
\item Il s'agit de l'\textbf{échappement} (4\textsuperscript{e} temps) : le piston monte et chasse les gaz brûlés par la soupape d'échappement ouverte. (Si c'était la compression, les deux soupapes seraient fermées.)
\item C'est un temps \textbf{résistant} : il ne produit pas de travail. Il est réalisé grâce à l'énergie emmagasinée par le \textbf{volant moteur} pendant le temps moteur.
\item Un cycle = 2 tours de vilebrequin et contient 1 temps moteur. Nombre de temps moteurs : $2400 \div 2 = 1200$ temps moteurs par minute.
\end{enumerate}
\end{corrige}

\begin{aretenir}[Bilan de la section]
\begin{itemize}
\item Le cycle à 4 temps = \textbf{admission, compression, combustion-détente, échappement} = 4 courses du piston = \textbf{2 tours de vilebrequin}.
\item Seule la \textbf{combustion-détente} est un temps moteur : les gaz brûlés poussent le piston et produisent le travail.
\item La bougie ne déclenche l'étincelle \textbf{qu'au début du temps moteur}, une fois par cycle.
\item Le \textbf{volant moteur} régularise la rotation en restituant, pendant les trois temps résistants, l'énergie emmagasinée pendant le temps moteur.
\item À $N$ tours par minute, un cylindre réalise $N/2$ temps moteurs par minute.
\end{itemize}
\end{aretenir}

\coursec{Le moteur à 2 temps et le moteur diesel}

Dans la section précédente, nous avons étudié le \cle{cycle à 4 temps} du moteur à essence : admission, compression, combustion-détente et échappement, réalisés en 4 courses du piston, soit 2 tours de vilebrequin. Ce moteur équipe la plupart des voitures. Mais sur les routes et les pistes du Cameroun, tu entends aussi le bruit aigu des motos-taxis et le ronronnement grave des camions et des groupes électrogènes. Ces machines n'utilisent pas toutes le même type de moteur. Dans cette section, nous découvrons deux autres familles de moteurs à combustion interne : le \cle{moteur à 2 temps}, simple et nerveux, et le \cle{moteur diesel}, puissant et économe.

\courssub{Le moteur à essence à 2 temps}

\textbf{Situation de départ.}\par
Observe une tronçonneuse en action ou une moto légère qui démarre : le moteur est petit, léger, et pourtant il crie fort et tourne vite. Le mécanicien verse dans le réservoir un mélange d'essence et d'huile, jamais de l'essence seule. Pourquoi ? Parce que ces engins utilisent un moteur à \cle{2 temps}, dont le fonctionnement diffère du 4 temps.

\begin{definition}[Moteur à 2 temps]
Un \cle{moteur à 2 temps} est un moteur à combustion interne dans lequel le cycle complet (admission, compression, combustion-détente, échappement) s'effectue en \textbf{2 courses du piston}, c'est-à-dire en \textbf{1 seul tour de vilebrequin}. Il ne possède \textbf{pas de soupapes} : les gaz entrent et sortent par des ouvertures percées dans la paroi du cylindre, appelées \cle{lumières}, que le piston découvre et obstrue en coulissant.
\end{definition}

\textbf{Fonctionnement.}\par
Le cycle se déroule ainsi :
\begin{enumerate}
\item \textbf{Première course (montée du piston)} : le piston monte et comprime le mélange air-essence déjà présent au-dessus de lui. En même temps, sa montée crée une dépression dans le carter (sous le piston) qui aspire un nouveau mélange frais par la \cle{lumière d'admission}. Au point mort haut, la bougie produit l'étincelle : \textbf{combustion}.
\item \textbf{Deuxième course (descente du piston)} : les gaz brûlés poussent le piston vers le bas (temps moteur). En descendant, le piston découvre d'abord la \cle{lumière d'échappement} (les gaz brûlés s'échappent), puis la \cle{lumière de transfert} : le mélange frais, mis en pression dans le carter, remonte dans le cylindre et chasse les gaz brûlés restants (balayage). Le cycle recommence.
\end{enumerate}

\begin{popfigure}
\begin{center}
\begin{tikzpicture}[scale=1.0, every node/.style={font=\small}]
% Cylindre
\draw[thick, popdark] (-1.6,0) rectangle (1.6,4.2);
% Piston
\fill[popblue!40] (-1.5,1.5) rectangle (1.5,2.3);
\draw[thick, popdark] (-1.5,1.5) rectangle (1.5,2.3);
\node at (0,1.9) {Piston};
% Bielle et vilebrequin
\draw[thick, popdark] (0,1.5) -- (0.5,0.2);
\draw[thick, popdark, fill=popmuted] (0.5,0.2) circle (0.12);
\draw[thick, popdark] (0.5,0.2) -- (0,-0.6);
\draw[thick, popdark, fill=popmuted] (0,-0.6) circle (0.1);
\node[below] at (0,-0.75) {Vilebrequin};
% Carter
\draw[thick, popdark] (-1.6,0) -- (-1.6,-1.1) -- (1.6,-1.1) -- (1.6,0);
\node at (0,-0.85) {\footnotesize Carter (mélange frais)};
% Bougie
\draw[thick, poporange, fill=poporangeL] (-0.25,4.2) rectangle (0.25,4.7);
\node[right, poporange] at (0.3,4.45) {Bougie};
% Lumières
\fill[popgreen!50] (-1.6,2.6) rectangle (-1.9,3.2);
\draw[thick, popdark] (-1.6,2.6) -- (-1.9,2.6) (-1.6,3.2) -- (-1.9,3.2) (-1.9,2.6) -- (-1.9,3.2);
\node[left, popgreen!60!black, align=right] at (-1.95,2.9) {Lumière\\ d'échappement};
\fill[popgold!50] (1.6,2.6) rectangle (1.9,3.2);
\draw[thick, popdark] (1.6,2.6) -- (1.9,2.6) (1.6,3.2) -- (1.9,3.2) (1.9,2.6) -- (1.9,3.2);
\node[right, popgold!60!black, align=left] at (1.95,2.9) {Lumière\\ de transfert};
\fill[popblue!30] (1.6,0.3) rectangle (1.9,0.9);
\draw[thick, popdark] (1.6,0.3) -- (1.9,0.3) (1.6,0.9) -- (1.9,0.9) (1.9,0.3) -- (1.9,0.9);
\node[right, popblue, align=left] at (1.95,0.6) {Lumière\\ d'admission};
% Flèche gaz brûlés
\draw[-{Stealth}, thick, popgreen!60!black] (-1.2,3.6) -- (-2.1,3.6);
% Flèche mélange
\draw[-{Stealth}, thick, popgold!60!black] (1.2,0.6) .. controls (1.4,1.6) .. (1.2,2.9);
\end{tikzpicture}
\end{center}
\legende{Schéma simplifié d'un moteur à 2 temps : pas de soupapes, mais trois lumières (admission, transfert, échappement) découvertes par le piston. Le mélange frais passe par le carter.}
\end{popfigure}

\textbf{Pourquoi un mélange essence-huile ?}\par
Dans un moteur 4 temps, l'huile est stockée dans un carter séparé et lubrifie les pièces en mouvement. Dans le 2 temps, le carter sert au passage du mélange frais : on ne peut donc pas y mettre d'huile. La solution : on mélange directement l'huile à l'essence (environ 2 à 4 \% d'huile). L'huile brûle avec l'essence, ce qui explique la fumée bleutée caractéristique des motos 2 temps.

\begin{propriete}[Avantages et inconvénients du moteur 2 temps]
\begin{itemize}
\item \textbf{Avantages} : moteur \textbf{simple} (pas de soupapes ni de distribution), \textbf{léger}, \textbf{puissant pour sa taille} (une combustion à chaque tour de vilebrequin, contre une tous les deux tours pour le 4 temps).
\item \textbf{Inconvénients} : il \textbf{consomme davantage} (une partie du mélange frais s'échappe avec les gaz brûlés lors du balayage) et il \textbf{pollue plus} (huile brûlée, mauvaise évacuation des gaz).
\end{itemize}
\end{propriete}

\begin{exemplebox}[Applications du moteur 2 temps]
On trouve les moteurs 2 temps partout où il faut un moteur léger et compact : les \textbf{motos légères} et motos-taxis, les \textbf{tronçonneuses} des bûcherons, les \textbf{débroussailleuses} des jardiniers, et les \textbf{moteurs hors-bord} qui propulsent les pirogues sur le Wouri ou le fleuve Sanaga.
\end{exemplebox}

\courssub{Le moteur diesel : l'allumage par compression}

\textbf{Situation de départ.}\par
Ouvre le capot d'un camion ou d'un groupe électrogène : tu ne trouveras \textbf{aucune bougie}. Pourtant, le gasoil s'enflamme bien dans les cylindres. Comment le mélange peut-il s'allumer sans étincelle ? Toute la réponse tient dans un phénomène que tu as déjà ressenti sans le savoir.

\begin{experience}[La compression d'un gaz l'échauffe]
\textbf{Matériel} : une pompe à vélo (ou une seringue bouchée).\\
\textbf{Protocole} :
\begin{enumerate}
\item Pompe vigoureusement un pneu de vélo pendant une minute.
\item Touche le corps de la pompe : il est chaud.
\item Variante : bouche l'extrémité d'une seringue avec le doigt et enfonce brutalement le piston. Le piston résiste et revient un peu : l'air comprimé s'est échauffé.
\end{enumerate}
\textbf{Observations} : le corps de la pompe chauffe nettement ; dans la seringue, l'air comprimé se comporte comme un ressort chaud.\\
\textbf{Interprétation} : comprimer rapidement un gaz, c'est concentrer son énergie dans un volume plus petit. L'agitation des molécules augmente : la \textbf{température du gaz s'élève}.\\
\textbf{Conclusion} : la compression rapide d'un gaz élève sa température.
\end{experience}

\begin{propriete}[Compression et température d'un gaz]
Lorsqu'on comprime rapidement un gaz, sa \textbf{température augmente} d'autant plus que la compression est forte. Dans un moteur diesel, l'air est comprimé environ 2 fois plus fort que dans un moteur à essence : sa température dépasse alors \SI{500}{\celsius}, ce qui suffit à enflammer le gasoil sans aucune étincelle.
\end{propriete}

\begin{definition}[Allumage par compression]
Dans un \cle{moteur diesel}, on dit que l'allumage se fait \textbf{par compression} : l'air seul est d'abord fortement comprimé, ce qui le chauffe à très haute température ; le gasoil est ensuite \textbf{injecté} sous forme de fines gouttelettes dans cet air brûlant et \textbf{s'enflamme spontanément}, sans bougie ni étincelle.
\end{definition}

\textbf{Fonctionnement du cycle diesel (4 temps).}\par
Le moteur diesel fonctionne lui aussi en 4 temps, mais avec deux différences majeures :
\begin{enumerate}
\item \textbf{Admission} : on aspire de l'\textbf{air seul} (pas de mélange).
\item \textbf{Compression} : l'air est comprimé très fortement (rapport volumétrique de 16 à 22, contre 8 à 10 pour l'essence) : il atteint une température très élevée.
\item \textbf{Injection et combustion-détente} : l'\cle{injecteur} pulvérise le gasoil dans l'air brûlant ; le carburant s'enflamme spontanément ; les gaz poussent le piston.
\item \textbf{Échappement} : les gaz brûlés sont rejetés.
\end{enumerate}

\begin{popfigure}
\begin{center}
\begin{tikzpicture}[scale=0.92, every node/.style={font=\small}]
% ===== Moteur essence =====
\begin{scope}
\draw[thick, popdark] (-1.4,0) rectangle (1.4,3.6);
\fill[popblue!30] (-1.3,1.2) rectangle (1.3,2.0);
\draw[thick, popdark] (-1.3,1.2) rectangle (1.3,2.0);
\draw[thick, popdark] (0,1.2) -- (0,0.2);
\draw[thick, popdark, fill=popmuted] (0,0.2) circle (0.1);
% Bougie
\draw[thick, poporange, fill=poporangeL] (-0.22,3.6) rectangle (0.22,4.1);
\draw[poporange, thick] (0,3.6) -- (0,3.35);
\node[poporange] at (0,4.4) {\textbf{Bougie}};
% Étincelle
\draw[popgold, thick] (0,3.3) -- (-0.15,3.05) (0,3.3) -- (0.15,3.05) (0,3.3) -- (0,2.95);
% Mélange
\node[popblue] at (0,2.9) {\footnotesize Mélange air + essence};
\node at (0,-0.5) {\textbf{Moteur à essence}};
\node[align=center] at (0,-1.2) {\footnotesize L'étincelle de la bougie\\ \footnotesize enflamme le mélange};
\end{scope}
% ===== Moteur diesel =====
\begin{scope}[xshift=8cm]
\draw[thick, popdark] (-1.4,0) rectangle (1.4,3.6);
\fill[popgreen!30] (-1.3,1.2) rectangle (1.3,2.0);
\draw[thick, popdark] (-1.3,1.2) rectangle (1.3,2.0);
\draw[thick, popdark] (0,1.2) -- (0,0.2);
\draw[thick, popdark, fill=popmuted] (0,0.2) circle (0.1);
% Injecteur
\draw[thick, poppink, fill=poppinkL] (-0.22,3.6) rectangle (0.22,4.1);
\draw[poppink, thick] (0,3.6) -- (0,3.35);
\node[poppink] at (0,4.4) {\textbf{Injecteur}};
% Pulvérisation
\draw[poppink, thick] (0,3.3) -- (-0.4,2.9) (0,3.3) -- (0,2.85) (0,3.3) -- (0.4,2.9);
% Air chaud
\node[popgreen!50!black] at (0,2.55) {\footnotesize Air seul très comprimé};
\node[popgreen!50!black] at (0,2.25) {\footnotesize (très chaud)};
\node at (0,-0.5) {\textbf{Moteur diesel}};
\node[align=center] at (0,-1.2) {\footnotesize Le gasoil injecté s'enflamme\\ \footnotesize au contact de l'air brûlant};
\end{scope}
\end{tikzpicture}
\end{center}
\legende{Comparaison : dans le moteur à essence, une bougie produit l'étincelle qui enflamme le mélange ; dans le moteur diesel, un injecteur pulvérise le gasoil dans l'air fortement comprimé et chauffé, qui s'enflamme spontanément.}
\end{popfigure}

\begin{propriete}[Avantages et inconvénients du moteur diesel]
\begin{itemize}
\item \textbf{Avantages} : \textbf{consommation réduite} (le gasoil est mieux exploité), \textbf{couple élevé} (grande force à bas régime, idéal pour tirer de lourdes charges), \textbf{robustesse} et longue durée de vie.
\item \textbf{Inconvénients} : moteur plus \textbf{lourd}, plus \textbf{bruyant}, plus coûteux à l'achat, et ses gaz d'échappement contiennent des particules fines nocives.
\end{itemize}
\end{propriete}

\begin{exemplebox}[Applications du moteur diesel]
Les \textbf{camions} de la SABC qui livrent les boissons, les \textbf{bus de voyage} qui relient Douala à Yaoundé, les \textbf{tracteurs} des plantations, les \textbf{groupes électrogènes} qui prennent le relais d'ENEO pendant les délestages, et les \textbf{pirogues de pêche} motorisées utilisent des moteurs diesel : il leur faut de la force, de l'endurance et une consommation maîtrisée.
\end{exemplebox}

\begin{attention}[Erreur fréquente : se tromper de carburant]
Mettre du \textbf{gasoil dans un moteur à essence}, ou de l'\textbf{essence dans un moteur diesel}, risque d'\textbf{endommager gravement le moteur} : le moteur à essence n'a pas de compression suffisante pour enflammer le gasoil (calage, encrassage), et l'essence détruit la lubrification des pompes à injection d'un diesel (grippage, casse). À la station-service, \textbf{vérifie toujours} le carburant adapté avant de faire le plein : regarde l'étiquette sur la trappe du réservoir ou demande au pompiste.
\end{attention}

\begin{savaistu}[Rudolf Diesel et l'huile d'arachide]
L'ingénieur allemand \textbf{Rudolf Diesel} a déposé le brevet de son moteur en \textbf{1892}. Fait remarquable : il voulait le faire fonctionner à l'\textbf{huile végétale}, notamment l'huile d'arachide, pour permettre aux agriculteurs de produire eux-mêmes leur carburant ! Plus d'un siècle plus tard, cette idée revient au goût du jour avec les \textbf{biocarburants} fabriqués à partir d'huile de palme ou de soja. Le moteur diesel porte d'ailleurs le nom de son inventeur.
\end{savaistu}

\begin{exoresolu}[Choisir le bon moteur]
Un pêcheur de Limbé veut motoriser sa pirogue légère pour de courtes sorties en mer. Un transporteur veut acheter un camion pour livrer des sacs de ciment entre Douala et Bafoussam. Pour chacun, indique le type de moteur le plus adapté (2 temps, 4 temps essence ou diesel) et justifie ton choix par deux arguments.
\tcblower
\textbf{Pour le pêcheur} : le \textbf{moteur 2 temps} (hors-bord) est le plus adapté.
\begin{itemize}
\item Il est \textbf{léger et compact} : facile à installer et à retirer sur une pirogue.
\item Il est \textbf{puissant pour sa taille} et \textbf{simple} : entretien facile, démarrage rapide pour de courtes sorties.
\end{itemize}
\textbf{Pour le transporteur} : le \textbf{moteur diesel} est le plus adapté.
\begin{itemize}
\item Il a un \textbf{couple élevé} : il peut tirer de lourdes charges de ciment, même en côte.
\item Il \textbf{consomme moins} et il est \textbf{robuste} : sur de longues distances, l'économie de carburant et la durée de vie sont décisives.
\end{itemize}
\textbf{Conclusion} : le choix d'un moteur dépend de l'usage : légèreté et simplicité pour le 2 temps, force et économie pour le diesel.
\end{exoresolu}

\begin{methode}[Reconnaître le type de moteur d'un engin]
\begin{enumerate}
\item \textbf{Écoute et observe} : un bruit aigu et une fumée bleutée indiquent un 2 temps ; un ronronnement grave et sourd indique un diesel.
\item \textbf{Cherche la bougie} : présence d'une bougie = moteur à essence ; présence d'injecteurs = moteur diesel.
\item \textbf{Regarde le carburant} : mélange essence-huile = 2 temps ; essence seule = 4 temps essence ; gasoil = diesel.
\item \textbf{Conclus} en reliant le type de moteur à l'usage de l'engin (léger ou lourd).
\end{enumerate}
\end{methode}

\begin{aretenir}[L'essentiel de la section]
\begin{itemize}
\item Le \textbf{moteur 2 temps} réalise son cycle en 2 courses du piston (1 tour de vilebrequin), sans soupapes mais avec des \textbf{lumières} ; il fonctionne au \textbf{mélange essence-huile}. Simple, léger et puissant, il consomme et pollue davantage. Applications : motos légères, tronçonneuses, hors-bord, débroussailleuses.
\item \textbf{Comprimer un gaz l'échauffe} : c'est le principe du moteur diesel.
\item Dans le \textbf{moteur diesel}, l'air seul est fortement comprimé et chauffé, puis le gasoil injecté \textbf{s'enflamme spontanément} : c'est l'\textbf{allumage par compression}. Pas de bougie.
\item Le diesel est \textbf{économe, puissant (couple élevé) et robuste}, mais lourd et bruyant. Applications : camions, bus, tracteurs, groupes électrogènes.
\item \textbf{Ne jamais confondre les carburants} : essence et gasoil ne sont pas interchangeables.
\end{itemize}
\end{aretenir}

\coursec{Classement des moteurs à combustion interne}

Dans la section précédente, nous avons découvert le moteur à 2 temps et le moteur diesel. Tu as donc déjà croisé plusieurs « familles » de moteurs : celui de la mobylette du voisin, celui de la moto-taxi 125 cm³, celui du camion qui monte la côte de Mbanga en grondant, celui du groupe électrogène qui ronronne pendant la coupure d'ENEO. Tous ces moteurs sont des moteurs à combustion interne, mais ils ne se ressemblent pas. Comment s'y retrouver ? Les techniciens et les mécaniciens les classent selon des \cle{critères} précis, comme on classe les élèves d'un collège par classe, puis par série. C'est l'objet de cette section : apprendre à ranger chaque moteur dans la bonne case, puis à lire la fiche technique d'un véhicule pour identifier son moteur.

\courssub{Pourquoi classer les moteurs ?}

\begin{exemplebox}[Une situation de la vie courante]
Monsieur Etoundi veut acheter une moto d'occasion à Bafoussam. Le vendeur lui dit : « C'est une 125, quatre temps, essence, refroidie à air. » Ces quatre informations ne sont pas du bavardage : elles décrivent le \cle{classement} du moteur. Elles permettent de savoir quel carburant acheter, quelle huile utiliser, quel mécanicien consulter et quelle consommation prévoir. Classer un moteur, c'est donc établir sa \emph{carte d'identité technique}.
\end{exemplebox}

\begin{definition}[Classement des moteurs à combustion interne]
Classer les moteurs à combustion interne, c'est les regrouper en catégories selon des critères techniques précis. Les quatre critères principaux au programme sont :
\begin{enumerate}
\item le \cle{cycle} de fonctionnement (2 temps ou 4 temps) ;
\item le \cle{carburant} utilisé (essence, gasoil, GPL, biocarburant) ;
\item le mode d'\cle{allumage} (commandé par bougie ou par compression) ;
\item le mode de \cle{refroidissement} (à air ou à eau).
\end{enumerate}
On ajoute, à titre de complément, la disposition des cylindres.
\end{definition}

\courssub{Critère 1 : le cycle de fonctionnement}

Tu connais déjà ce critère : c'est le nombre de temps (courses du piston) nécessaires pour réaliser un cycle complet.

\begin{itemize}
\item \textbf{Moteur à 4 temps} : le cycle complet (admission, compression, combustion-détente, échappement) s'effectue en \textbf{4 courses du piston}, soit 2 tours de vilebrequin. C'est le cas de la plupart des motos-taxis 125 cm³, des voitures et des groupes électrogènes.
\item \textbf{Moteur à 2 temps} : le cycle complet s'effectue en \textbf{2 courses du piston}, soit 1 seul tour de vilebrequin. C'est le cas des mobylettes, des tronçonneuses et des moteurs hors-bord de pirogues.
\end{itemize}

\begin{aretenir}[L'essentiel sur le cycle]
Le critère « cycle » répond à la question : \emph{en combien de courses du piston le moteur réalise-t-il son cycle complet ?} Réponse possible : 2 temps ou 4 temps. Ce critère ne dit rien sur le nombre de cylindres.
\end{aretenir}

\courssub{Critère 2 : le carburant utilisé}

Le carburant est le « aliment » du moteur : c'est lui qui, en brûlant, fournit l'énergie. On distingue :

\begin{itemize}
\item \textbf{L'essence} (supercarburant) : carburant léger et très volatil, utilisé par les moteurs à allumage commandé (motos, voitures particulières, groupes électrogènes). À la station Total ou Ola Energy, c'est le pistolet « super ».
\item \textbf{Le gasoil} (ou diesel) : carburant plus lourd et moins volatil, utilisé par les moteurs à allumage par compression (camions, bus, pick-up, gros groupes électrogènes). Il est produit notamment par la SONARA à Limbé.
\item \textbf{Le GPL} (gaz de pétrole liquéfié) : le même gaz que celui des bonbonnes de cuisine (butane, propane). Certaines voitures et certains groupes électrogènes sont adaptés pour fonctionner au GPL, moins polluant et moins cher.
\item \textbf{Les biocarburants} : carburants fabriqués à partir de végétaux (huile de palme, éthanol de canne à sucre). Ils peuvent être mélangés au gasoil ou à l'essence. C'est une piste d'avenir pour un pays agricole comme le Cameroun.
\end{itemize}

\begin{attention}[Ne jamais se tromper de carburant]
Mettre de l'essence dans un moteur diesel (ou l'inverse) provoque de graves dommages : le moteur diesel n'a pas de bougie et compte sur la compression pour enflammer le gasoil ; l'essence détruit la pompe d'injection. Avant de faire le plein, vérifie toujours la mention sur le bouchon du réservoir.
\end{attention}

\courssub{Critère 3 : le mode d'allumage}

Pour que le mélange air-carburant brûle, il faut déclencher la combustion. Deux solutions existent :

\begin{itemize}
\item \textbf{Allumage commandé (par bougie)} : une \cle{bougie} produit une étincelle électrique au bon moment, qui enflamme le mélange air-essence. C'est le cas des moteurs à \textbf{essence}.
\item \textbf{Allumage par compression (auto-allumage)} : l'air seul est comprimé très fortement, ce qui l'échauffe à plus de \SI{500}{\celsius} ; le gasoil injecté alors s'enflamme \emph{tout seul}, sans bougie. C'est le cas des moteurs \textbf{diesel}.
\end{itemize}

\begin{experience}[Observer la différence d'allumage]
\textbf{Matériel} : une moto-taxi à l'arrêt, moteur ouvert par le mécanicien du quartier (observation encadrée).
\textbf{Protocole} : repérer, au sommet du cylindre, une petite pièce blanche (porcelaine) coiffée d'un fil noir : c'est la bougie. Comparer avec la photo du cylindre d'un moteur diesel de camion : à la place de la bougie, on trouve un \emph{injecteur}.
\textbf{Observation} : le moteur essence possède une bougie ; le moteur diesel possède un injecteur et pas de bougie d'allumage.
\textbf{Interprétation} : l'essence s'enflamme grâce à l'étincelle de la bougie (allumage commandé) ; le gasoil s'enflamme grâce à la chaleur de l'air fortement comprimé (allumage par compression).
\textbf{Conclusion} : le mode d'allumage est directement lié au carburant : essence $\rightarrow$ bougie ; gasoil $\rightarrow$ compression.
\end{experience}

\courssub{Critère 4 : le mode de refroidissement}

La combustion dégage énormément de chaleur. Sans refroidissement, les pièces métalliques du moteur se dilateraient, gripperaient et se détruiraient. Deux grandes solutions :

\begin{itemize}
\item \textbf{Refroidissement à air} : le cylindre est entouré d'\cle{ailettes} métalliques qui augmentent la surface de contact avec l'air ; l'air en mouvement emporte la chaleur. Solution simple et légère, utilisée sur les \textbf{motos}, mobylettes et petits groupes électrogènes. Observe le moteur d'une moto-taxi : tu verras ces fines lamelles autour du cylindre.
\item \textbf{Refroidissement à eau (à liquide)} : un liquide (eau + liquide de refroidissement) circule dans des canaux autour des cylindres, puis se refroidit dans un \cle{radiateur} balayé par l'air. Solution plus efficace mais plus lourde, utilisée sur les \textbf{voitures}, camions et gros moteurs.
\end{itemize}

\begin{popfigure}
\begin{tikzpicture}[scale=0.9]
% Cylindre refroidi à air
\draw[fill=popblueL,thick] (0,0) rectangle (1.6,2.2);
\foreach \y in {0.3,0.7,1.1,1.5,1.9}{
  \draw[thick,fill=popblue!40] (-0.5,\y) rectangle (0,\y+0.12);
  \draw[thick,fill=popblue!40] (1.6,\y) rectangle (2.1,\y+0.12);}
\draw[->,poporange,thick] (-1.2,1.1) -- (-0.6,1.1);
\draw[->,poporange,thick] (2.2,1.1) -- (2.8,1.1);
\node[popdark] at (0.8,-0.5) {\small \textbf{Refroidissement à air (moto)}};
\node[popmuted] at (0.8,2.6) {\small ailettes};
% Cylindre refroidi à eau
\begin{scope}[xshift=6.5cm]
\draw[fill=popgreenL,thick] (0,0) rectangle (1.6,2.2);
\draw[thick,popblue] (-0.4,0) -- (-0.4,2.2) -- (2.0,2.2) -- (2.0,0);
\draw[->,popblue,thick] (-0.4,-0.4) -- (-0.4,0.6);
\draw[->,popblue,thick] (2.0,1.6) -- (2.0,2.6);
\draw[fill=popblue!25,thick] (3.0,0.4) rectangle (3.5,2.0);
\foreach \x in {3.08,3.18,3.28,3.38}{\draw[popblue] (\x,0.4) -- (\x,2.0);}
\draw[->,popblue] (2.0,2.4) .. controls (2.6,2.8) .. (3.25,2.15);
\draw[->,popblue] (3.25,0.25) .. controls (2.4,-0.3) .. (-0.4,-0.25);
\node[popdark] at (1.6,-0.5) {\small \textbf{Refroidissement à eau (voiture)}};
\node[popmuted] at (3.25,2.4) {\small radiateur};
\end{scope}
\end{tikzpicture}
\legende{Les deux modes de refroidissement. À gauche : les ailettes (bleu) du cylindre cèdent la chaleur à l'air (flèches orange). À droite : le liquide de refroidissement (circuit bleu) transporte la chaleur du cylindre vers le radiateur.}
\end{popfigure}

\courssub{Critère 5 : la disposition des cylindres (complément)}

Ce critère n'est pas exigible au BEPC, mais il est utile pour lire une fiche technique. Quand un moteur possède plusieurs cylindres, on peut les disposer :

\begin{itemize}
\item \textbf{en ligne} : les cylindres sont alignés côte à côte (exemple : 4 cylindres en ligne d'une Toyota Corolla) ;
\item \textbf{en V} : les cylindres forment deux rangées inclinées qui se rejoignent au vilebrequin, comme les branches d'un V (moteurs V6, V8 des grosses voitures et 4$\times$4) ;
\item \textbf{à plat} (dits « boxer ») : les cylindres sont horizontaux, face à face (certaines motos BMW et voitures Subaru).
\end{itemize}

\begin{attention}[L'erreur classique : 4 temps $\neq$ 4 cylindres]
Beaucoup d'élèves confondent « 4 temps » et « 4 cylindres ». C'est faux !
\begin{itemize}
\item Le nombre de \textbf{temps} décrit le \emph{cycle de fonctionnement} (2 ou 4 courses du piston).
\item Le nombre de \textbf{cylindres} décrit la \emph{construction} du moteur.
\end{itemize}
Une moto-taxi 125 cm³ est un moteur \textbf{4 temps} qui n'a qu'\textbf{un seul cylindre}. Une voiture peut être 4 temps avec 3, 4, 6 ou 8 cylindres. Au BEPC, cette confusion coûte des points : relis-toi !
\end{attention}

\courssub{Arbre de classification et tableau récapitulatif}

\begin{popfigure}
\begin{tikzpicture}[
  grow=right,
  level 1/.style={sibling distance=3.4cm, level distance=3.1cm},
  level 2/.style={sibling distance=1.6cm, level distance=3.4cm},
  every node/.style={font=\small},
  edge from parent/.style={draw=popmuted,thick}]
\node[fill=popdark,text=white,rounded corners=2pt,inner sep=4pt] {Moteurs à combustion interne}
  child { node[fill=popgoldL,rounded corners=2pt,inner sep=3pt] {Refroidissement}
    child { node[fill=popcream,rounded corners=2pt,inner sep=2pt] {à eau : voiture} }
    child { node[fill=popcream,rounded corners=2pt,inner sep=2pt] {à air : moto} } }
  child { node[fill=poppurpleL,rounded corners=2pt,inner sep=3pt] {Allumage}
    child { node[fill=popcream,rounded corners=2pt,inner sep=2pt] {par compression : diesel} }
    child { node[fill=popcream,rounded corners=2pt,inner sep=2pt] {par bougie : essence} } }
  child { node[fill=popgreenL,rounded corners=2pt,inner sep=3pt] {Carburant}
    child { node[fill=popcream,rounded corners=2pt,inner sep=2pt] {GPL, biocarburants} }
    child { node[fill=popcream,rounded corners=2pt,inner sep=2pt] {gasoil, essence} } }
  child { node[fill=popblueL,rounded corners=2pt,inner sep=3pt] {Cycle}
    child { node[fill=popcream,rounded corners=2pt,inner sep=2pt] {4 temps : moto 125} }
    child { node[fill=popcream,rounded corners=2pt,inner sep=2pt] {2 temps : tronçonneuse} } };
\end{tikzpicture}
\legende{Arbre de classification des moteurs à combustion interne selon les quatre critères au programme : le cycle, le carburant, l'allumage et le refroidissement.}
\end{popfigure}

\begin{center}
\small
\begin{tabularx}{\linewidth}{|l|X|X|X|}
\hline
\rowcolor{popblueL} \textbf{Critère} & \textbf{Essence 2 temps} & \textbf{Essence 4 temps} & \textbf{Diesel 4 temps} \\
\hline
Carburant & essence + huile mélangée & essence & gasoil \\
\hline
Allumage & bougie (étincelle) & bougie (étincelle) & compression (sans bougie) \\
\hline
Refroidissement & à air & à air (moto) ou à eau (voiture) & à eau \\
\hline
Avantages & léger, simple, puissant pour sa taille & économique, propre, durable & couple élevé, très robuste \\
\hline
Exemples d'usage & tronçonneuse, mobylette, hors-bord & moto-taxi 125, voiture, groupe électrogène & camion, bus, pick-up, tracteur \\
\hline
\end{tabularx}
\end{center}

\courssub{Méthode : classer un moteur à partir d'une fiche technique}

\begin{methode}[Classer un moteur en 4 questions]
Face à la fiche technique d'un véhicule, pose-toi successivement quatre questions, dans l'ordre :
\begin{enumerate}
\item \textbf{Quel carburant ?} Cherche « essence », « diesel/gasoil » ou « GPL » sur la fiche.
\item \textbf{Quel cycle ?} Cherche « 2 temps » ou « 4 temps » (souvent noté 2T ou 4T).
\item \textbf{Quel allumage ?} Essence $\rightarrow$ allumage commandé par bougie ; diesel $\rightarrow$ allumage par compression.
\item \textbf{Quel refroidissement ?} Cherche « refroidi à air » (ailettes) ou « refroidi à eau/liquide » (radiateur).
\end{enumerate}
Tu peux ensuite compléter par le nombre de cylindres et la cylindrée (en cm³), qui précisent la taille du moteur sans changer sa catégorie.
\end{methode}

\begin{exemplebox}[La moto-taxi 125 cm³ passée au crible]
Appliquons la méthode à la moto-taxi qui t'amène au collège :
\begin{enumerate}
\item Carburant : \textbf{essence} (elle fait le plein de super à la station) ;
\item Cycle : \textbf{4 temps} (mention 4T sur le carter) ;
\item Allumage : \textbf{commandé par bougie} (conséquence du carburant essence) ;
\item Refroidissement : \textbf{à air} (ailettes visibles autour du cylindre unique).
\end{enumerate}
Conclusion : c'est un moteur à combustion interne, \emph{essence, 4 temps, à allumage commandé, refroidi à air, monocylindre} (un seul cylindre).
\end{exemplebox}

\begin{exoresolu}[Lecture de la fiche technique d'un pick-up]
Énoncé. Voici l'extrait de la fiche technique d'un pick-up utilisé pour le transport entre Douala et Bafoussam : « Moteur : 4 cylindres en ligne ; cylindrée : \num{2494} cm³ ; carburant : diesel ; refroidissement : liquide ; cycle : 4 temps. »
\begin{enumerate}
\item Classer ce moteur selon les quatre critères du programme.
\item Le conducteur affirme : « 4 temps, donc 4 cylindres. » A-t-il raison ? Justifier.
\end{enumerate}
\tcblower
Solution détaillée.
\begin{enumerate}
\item \textbf{Carburant} : gasoil (diesel). \textbf{Cycle} : 4 temps. \textbf{Allumage} : par compression, car tout moteur diesel fonctionne sans bougie, par auto-allumage du gasoil dans l'air surchauffé. \textbf{Refroidissement} : à eau (liquide + radiateur). Le moteur est donc : \emph{diesel, 4 temps, à allumage par compression, refroidi à eau}, avec 4 cylindres en ligne.
\item Le conducteur a tort. « 4 temps » décrit le \emph{cycle} (4 courses du piston) et non le nombre de cylindres. Ici, les deux nombres valent 4 par pure coïncidence : une moto 4 temps n'a souvent qu'un seul cylindre, et il existe des moteurs 4 temps à 6 ou 8 cylindres.
\end{enumerate}
\end{exoresolu}

\begin{savaistu}[Des moteurs géants]
Les moteurs à combustion interne ne se trouvent pas que sur la route ! Les avions de ligne utilisent des dérivés de moteurs à combustion (turbopropulseurs), et les navires sont poussés par d'immenses moteurs diesel. Le plus gros moteur diesel du monde (le Wärtsilä RT-flex96C, qui équipe des porte-conteneurs) mesure environ \SI{13,5}{m} de haut : chacun de ses 14 cylindres est plus grand qu'une salle de classe de ton collège, et un seul piston pèse plus de 5 tonnes ! À l'autre extrémité, le moteur d'une tronçonneuse tient dans une main. Même principe, même famille : la combustion interne.
\end{savaistu}

\begin{exercice}[À toi de classer]
Pour chaque engin, classe le moteur selon les quatre critères (carburant, cycle, allumage, refroidissement) :
\begin{enumerate}
\item Une tronçonneuse (mélange essence-huile, 2 temps, ailettes) ;
\item Un bus de voyage de 70 places (gasoil, 4 temps, radiateur) ;
\item Un groupe électrogène de quartier (essence, 4 temps, ailettes).
\end{enumerate}
\end{exercice}

\begin{corrige}[Exercice : À toi de classer]
\begin{enumerate}
\item Tronçonneuse : essence, 2 temps, allumage commandé par bougie, refroidissement à air.
\item Bus : diesel (gasoil), 4 temps, allumage par compression, refroidissement à eau.
\item Groupe électrogène : essence, 4 temps, allumage commandé par bougie, refroidissement à air.
\end{enumerate}
\end{corrige}

\begin{aretenir}[Ce qu'il faut retenir pour le BEPC]
\begin{itemize}
\item Les moteurs à combustion interne se classent selon 4 critères : le \cle{cycle} (2T/4T), le \cle{carburant} (essence, gasoil, GPL, biocarburant), l'\cle{allumage} (bougie ou compression) et le \cle{refroidissement} (air ou eau).
\item Essence $\rightarrow$ bougie ; diesel $\rightarrow$ compression. Ce lien est automatique.
\item « 4 temps » ne signifie jamais « 4 cylindres » : le temps concerne le cycle, le cylindre concerne la construction.
\item Pour classer un moteur réel, lis sa fiche technique et réponds aux 4 questions de la méthode.
\end{itemize}
\end{aretenir}

\coursec{Sécurité, entretien et impact environnemental}

Dans la section précédente, nous avons classé les moteurs à combustion interne selon leur carburant, leur cycle et leur allumage. Ces machines sont partout autour de nous : moto-taxi, groupe électrogène pendant les délestages, pompe à eau au champ, voiture familiale. Mais un moteur, c'est aussi un réservoir de carburant inflammable, des pièces qui tournent vite, des surfaces brûlantes et des gaz d'échappement toxiques. Un moteur mal utilisé ou mal entretenu peut blesser, tuer, et polluer gravement notre environnement. Cette dernière section est peut-être la plus importante de la leçon : elle peut sauver des vies, à commencer par la tienne.

\courssub{Les risques liés aux carburants}

\begin{definition}[Carburant inflammable]
Un \cle{carburant} (essence, gasoil) est un combustible liquide qui peut s'enflammer au contact d'une flamme ou d'une étincelle. L'\cle{essence} est particulièrement dangereuse : ses \cle{vapeurs} s'enflamment même à distance du liquide, dès la température ambiante.
\end{definition}

Observe une station-service : il est interdit d'y fumer, d'y téléphoner près des pompes ou de remplir un bidon posé sur un plastique. Pourquoi ? L'essence dégage en permanence des vapeurs invisibles qui se mélangent à l'air. Ce mélange air-vapeurs d'essence est explosif : une simple étincelle (briquet, étincelle électrique, frottement) suffit à l'enflammer. Le gasoil s'enflamme moins facilement, mais il brûle aussi une fois allumé.

\begin{attention}[Règles d'or avec les carburants]
\begin{itemize}
\item Stocker l'essence et le gasoil dans des \textbf{bidons fermés}, à l'\textbf{ombre}, loin de toute flamme, de toute source de chaleur et hors de portée des enfants.
\item \textbf{Ne jamais remplir le réservoir d'un moteur chaud ou en marche} : une goutte d'essence sur le moteur brûlant peut s'enflammer instantanément. On arrête le moteur et on attend qu'il refroidisse avant le ravitaillement.
\item Ne jamais fumer ni approcher une flamme (bougie, lampe tempête, cuisinière) pendant une manipulation de carburant.
\item En cas de tache de carburant, essuyer immédiatement et aérer.
\end{itemize}
\end{attention}

\begin{savaistu}[Le danger de la vente d'essence au bord des routes]
Dans plusieurs villes du Cameroun, on vend de l'essence en bouteilles au bord des routes. Des accidents graves se produisent quand ce stockage se fait dans des boutiques où l'on cuisine au feu de bois ou au gaz juste à côté. Les vapeurs d'essence, plus lourdes que l'air, s'accumulent au niveau du sol et peuvent s'enflammer à plusieurs mètres du bidon. C'est pourquoi la réglementation impose des conditions strictes de stockage.
\end{savaistu}

\courssub{Les risques liés aux gaz d'échappement : le monoxyde de carbone}

\begin{experience}[Observation : que contiennent les gaz d'échappement ?]
\textbf{Observation 1 :} place une feuille blanche quelques secondes derrière le pot d'échappement d'une vieille moto en marche (à l'extérieur !). La feuille se couvre de particules noires : ce sont des \cle{suies}.\par
\textbf{Observation 2 :} on ne voit rien, on ne sent rien de particulier... et pourtant les gaz d'échappement contiennent un poison mortel : le \cle{monoxyde de carbone} \ce{CO}, produit par une combustion incomplète du carburant.\par
\textbf{Interprétation :} les gaz d'échappement contiennent du dioxyde de carbone \ce{CO2}, de la vapeur d'eau, des suies, des oxydes d'azote et du monoxyde de carbone \ce{CO}. Le \ce{CO} est \textbf{invisible, inodore et sans saveur} : aucun de nos sens ne peut le détecter.
\end{experience}

\begin{definition}[Monoxyde de carbone]
Le \cle{monoxyde de carbone}, de formule \ce{CO}, est un gaz toxique produit par les combustions incomplètes. Respiré, il prend la place du dioxygène \ce{O2} dans le sang et empêche l'oxygénation des organes : c'est l'\cle{intoxication au CO}, qui peut provoquer maux de tête, nausées, perte de conscience, puis la mort.
\end{definition}

\begin{attention}[Danger mortel : jamais de moteur dans une pièce fermée]
Le monoxyde de carbone \ce{CO} \textbf{tue sans odeur et sans couleur}. On s'endort sans s'en rendre compte. Par conséquent :
\begin{itemize}
\item un \textbf{groupe électrogène} fonctionne \textbf{toujours à l'extérieur}, dans un endroit ventilé, loin des portes et des fenêtres (les gaz peuvent rentrer dans la maison !) ;
\item on ne fait jamais tourner une moto ou une voiture dans un garage fermé ;
\item les premiers signes d'intoxication (maux de tête, fatigue, nausées) doivent alerter immédiatement : sortir à l'air libre et appeler de l'aide.
\end{itemize}
\end{attention}

\begin{exemplebox}[Une tragédie évitable]
Chaque année au Cameroun, pendant les longues coupures d'électricité, des familles installent le groupe électrogène dans la cour fermée, sous le porche, voire dans un couloir de la maison « à cause des voleurs ». Des intoxications au \ce{CO} s'ensuivent, parfois mortelles pendant le sommeil. La solution n'est pas de rentrer le groupe, mais de le \textbf{sécuriser dehors} : cadenas, chaîne, petit abri ventilé. La sécurité électrique ne doit jamais faire oublier la sécurité respiratoire.
\end{exemplebox}

\begin{popfigure}
\begin{center}
\begin{tikzpicture}[scale=0.9]
% Pictogramme inflammable
\begin{scope}[shift={(0,0)}]
\draw[fill=poppinkL,draw=poppink,very thick] (0,1.5) -- (1.5,0) -- (0,-1.5) -- (-1.5,0) -- cycle;
\draw[fill=poporange,draw=popdark] (0,0.7) .. controls (0.35,0.2) and (0.45,-0.1) .. (0.3,-0.45) .. controls (0.15,-0.15) and (-0.15,-0.15) .. (-0.3,-0.45) .. controls (-0.45,-0.1) and (-0.35,0.2) .. (0,0.7);
\draw[fill=poppink!70!poporange] (0,0.15) .. controls (0.15,-0.05) and (0.15,-0.2) .. (0,-0.4) .. controls (-0.15,-0.2) and (-0.15,-0.05) .. (0,0.15);
\node[below,font=\small\bfseries] at (0,-1.7) {Inflammable};
\end{scope}
% Pictogramme gaz toxique
\begin{scope}[shift={(4.5,0)}]
\draw[fill=poppurpleL,draw=poppurple,very thick] (0,1.5) -- (1.5,0) -- (0,-1.5) -- (-1.5,0) -- cycle;
\draw[fill=white,draw=popdark,thick] (0,0.55) circle (0.42);
\fill[popdark] (-0.16,0.62) circle (0.09);
\fill[popdark] (0.16,0.62) circle (0.09);
\draw[popdark,thick] (-0.12,0.38) -- (0.12,0.38);
\draw[popdark,thick] (-0.35,-0.35) -- (0.35,-0.75);
\draw[popdark,thick] (0.35,-0.35) -- (-0.35,-0.75);
\fill[popdark] (-0.42,-0.32) circle (0.07);
\fill[popdark] (0.42,-0.32) circle (0.07);
\fill[popdark] (-0.42,-0.78) circle (0.07);
\fill[popdark] (0.42,-0.78) circle (0.07);
\node[below,font=\small\bfseries] at (0,-1.7) {Gaz toxique (CO)};
\end{scope}
% Pictogramme surface chaude
\begin{scope}[shift={(9,0)}]
\draw[fill=popgoldL,draw=popgold,very thick] (0,1.5) -- (1.5,0) -- (0,-1.5) -- (-1.5,0) -- cycle;
\draw[popdark,very thick] (-0.5,-0.75) -- (0.5,-0.75);
\draw[popdark,thick] (-0.3,-0.5) .. controls (-0.45,-0.2) and (-0.15,0) .. (-0.3,0.3) .. controls (-0.4,0.5) and (-0.25,0.6) .. (-0.3,0.75);
\draw[popdark,thick] (0.05,-0.5) .. controls (-0.1,-0.2) and (0.2,0) .. (0.05,0.3) .. controls (-0.05,0.5) and (0.1,0.6) .. (0.05,0.75);
\draw[popdark,thick] (0.4,-0.5) .. controls (0.25,-0.2) and (0.55,0) .. (0.4,0.3) .. controls (0.3,0.5) and (0.45,0.6) .. (0.4,0.75);
\node[below,font=\small\bfseries] at (0,-1.7) {Surface chaude};
\end{scope}
\end{tikzpicture}
\end{center}
\legende{Pictogrammes de sécurité à connaître : produit inflammable (flamme), gaz toxique (tête de mort), surface chaude (chaleur qui monte d'une surface).}
\end{popfigure}

\courssub{Les risques mécaniques et thermiques}

Un moteur en marche, c'est aussi de la mécanique en mouvement et de la chaleur :

\begin{itemize}
\item \textbf{Pièces en rotation} : courroies, poulies, ventilateur, chaîne de moto. Un doigt, un pagne, une cravate ou des cheveux longs happés par une courroie en rotation provoquent des blessures graves. On ne met \textbf{jamais les mains} près d'une pièce qui tourne.
\item \textbf{Surfaces brûlantes} : le pot d'échappement dépasse largement \SI{100}{\celsius} ; beaucoup d'enfants portent des cicatrices de brûlures au mollet après être descendus d'une moto du côté du pot. Le moteur et le radiateur restent brûlants plusieurs minutes après l'arrêt.
\item \textbf{Radiateur sous pression} : ne jamais ouvrir le bouchon du radiateur à chaud ! L'eau surchauffée peut jaillir en vapeur brûlante. On attend le refroidissement complet.
\end{itemize}

\begin{methode}[Manipuler un moteur en atelier : les bons gestes]
\begin{enumerate}
\item \textbf{Arrêter} le moteur et retirer la clé de contact (ou débrancher la bougie sur un petit moteur).
\item \textbf{Laisser refroidir} le moteur, le pot d'échappement et le radiateur.
\item Vérifier qu'\textbf{aucune flamme} ni cigarette ne se trouve à proximité (vapeurs d'essence).
\item Porter des \textbf{gants} et des \textbf{lunettes de protection} ; attacher les cheveux longs, retirer les vêtements flottants.
\item Travailler dans un endroit \textbf{ventilé} ; si le moteur doit tourner, ce sera dehors.
\item Après la manipulation, \textbf{se laver les mains} (huile, essence, graisse sont nocives pour la peau).
\end{enumerate}
\end{methode}

\courssub{L'entretien de base d'un moteur}

Un moteur bien entretenu dure plus longtemps, consomme moins et pollue moins. Voici les gestes d'entretien que tout utilisateur doit connaître.

\begin{center}
\begin{tabularx}{\linewidth}{|l|X|X|}\hline
\rowcolor{popblueL} \textbf{Opération} & \textbf{Rôle} & \textbf{Fréquence indicative} \\ \hline
Vidange d'huile & Remplacer l'huile usée qui lubrifie les pièces en mouvement et évite leur usure & Tous les \num{1000} à \num{3000} km (moto) ou selon le constructeur \\ \hline
Filtre à air & Retenir la poussière de l'air aspiré ; un filtre encrassé étouffe le moteur et augmente la consommation & Nettoyage régulier, surtout en saison sèche et sur pistes poussiéreuses \\ \hline
Bougie & Produire l'étincelle d'allumage (moteur à essence) ; une bougie encrassée provoque des ratés et des fumées & Vérification et nettoyage périodiques, remplacement si usée \\ \hline
Niveau de carburant et de liquide de refroidissement & Éviter la panne sèche et la surchauffe du moteur & Avant chaque trajet important \\ \hline
\end{tabularx}
\end{center}

\begin{exoresolu}[Choisir la bonne conduite à tenir]
Énoncé : Ngono possède une moto. Depuis quelques semaines, le moteur « broute », démarre mal et rejette des fumées noires. Un ami lui conseille de continuer à rouler ainsi « tant que ça roule ». Quelles vérifications simples Ngono doit-il faire, et pourquoi ?
\tcblower
Solution détaillée : les symptômes (démarrage difficile, ratés, fumées noires) indiquent une \textbf{mauvaise combustion}. Les vérifications de base sont :
\begin{enumerate}
\item \textbf{Nettoyer ou remplacer le filtre à air} : s'il est encrassé de poussière, l'air arrive mal et le mélange devient trop riche en essence, d'où les fumées noires.
\item \textbf{Vérifier et nettoyer la bougie} : une bougie encrassée produit une étincelle faible, d'où les ratés d'allumage.
\item \textbf{Contrôler l'huile} : une huile trop vieille use le moteur.
\end{enumerate}
Continuer à rouler sans entretien augmente la consommation, pollue davantage et finit par détruire le moteur : le conseil de l'ami est donc mauvais. Un entretien simple et régulier coûte bien moins cher qu'une réparation complète.
\end{exoresolu}

\courssub{L'impact environnemental des moteurs}

\begin{definition}[Pollution liée aux moteurs à combustion interne]
En brûlant le carburant, les moteurs rejettent dans l'air des \cle{polluants} : monoxyde de carbone \ce{CO} (toxique), dioxyde de carbone \ce{CO2} (gaz à effet de serre qui contribue au réchauffement climatique), suies et fumées noires, oxydes d'azote. Ils produisent aussi une \cle{pollution sonore} (bruit), surtout quand le pot d'échappement est défectueux ou modifié.
\end{definition}

Le dioxyde de carbone \ce{CO2} n'est pas toxique directement, mais sa quantité croissante dans l'atmosphère renforce l'effet de serre et dérègle le climat. Les fumées noires des vieux moteurs diesel et des motos mal réglées irritent les bronches, surtout chez les enfants et les personnes âgées. Le bruit permanent des moteurs dans les grandes villes fatigue et rend malade.

\begin{savaistu}[Un moteur mal réglé, c'est du gaspillage]
Un moteur mal réglé (filtre à air encrassé, bougie usée, mauvais réglage du carburateur) peut consommer jusqu'à \textbf{25 \% de carburant en plus} pour parcourir la même distance, et il pollue beaucoup davantage. Entretenir son moteur, c'est donc à la fois économiser de l'argent et protéger l'air que respire tout le quartier.
\end{savaistu}

\courssub{Les gestes éco-citoyens}

Chacun peut réduire les dangers et la pollution liés aux moteurs par des gestes simples :

\begin{itemize}
\item \textbf{Couper le moteur} lors d'un arrêt prolongé (devant une école, dans un embouteillage immobile) : un moteur qui tourne à l'arrêt consomme et pollue pour rien.
\item \textbf{Entretenir régulièrement} son moteur : vidange, filtre à air, bougie.
\item \textbf{Éviter les fuites d'huile} : l'huile usée ne se jette jamais au sol ni dans les caniveaux ; elle s'infiltre dans le sol, pollue les nappes d'eau souterraines et les puits. On la récupère dans un bidon et on la remet à un garage ou à un point de collecte.
\item \textbf{Installer le groupe électrogène dehors}, sur un sol stable, à distance des fenêtres, et le faire réviser régulièrement.
\item Privilégier, quand c'est possible, la \textbf{marche, le vélo ou le transport en commun} pour les courts trajets.
\end{itemize}

\begin{popfigure}
\begin{center}
\begin{tikzpicture}[scale=1]
% Sol
\draw[fill=popgreenL] (-0.5,0) rectangle (11.5,-0.25);
\draw[popgreen] (-0.5,0) -- (11.5,0);
% Maison
\draw[fill=popcream,draw=popdark,thick] (7,0) rectangle (11,3);
\draw[fill=poporange,draw=popdark,thick] (6.7,3) -- (9,4.4) -- (11.3,3) -- cycle;
% Fenêtre
\draw[fill=popblueL,draw=popdark,thick] (7.5,1.4) rectangle (8.5,2.4);
\draw[popdark] (8,1.4) -- (8,2.4);
\draw[popdark] (7.5,1.9) -- (8.5,1.9);
\node[font=\small] at (9,3.5) {\textbf{Maison}};
% Porte
\draw[fill=popgoldL,draw=popdark,thick] (9.5,0) rectangle (10.4,1.7);
% Groupe électrogène dehors
\draw[fill=popblue,draw=popdark,thick] (1,0) rectangle (3,1.3);
\draw[fill=popblueL,draw=popdark] (1.2,1.3) rectangle (2.8,1.7);
\fill[popdark] (1.4,0) circle (0.18);
\fill[popdark] (2.6,0) circle (0.18);
\draw[popdark,thick] (3,0.9) -- (3.5,0.9);
\node[font=\small,align=center] at (2,-0.6) {\textbf{Groupe électrogène}\\ \textbf{à l'extérieur}};
% Fumée qui s'éloigne de la maison
\draw[popmuted,thick,decorate,decoration={snake,amplitude=1.5mm,segment length=4mm}] (3.5,0.9) -- (4.6,1.6);
\draw[popmuted,thick,decorate,decoration={snake,amplitude=1.5mm,segment length=4mm}] (4.6,1.6) -- (3.9,2.6);
\draw[popmuted,thick,decorate,decoration={snake,amplitude=1.5mm,segment length=4mm}] (3.9,2.6) -- (2.6,3.3);
\node[font=\small,popmuted] at (2.2,3.7) {gaz d'échappement évacués au loin};
% Distance sécurité
\draw[<->,popgreen!60!popdark,very thick] (3.5,0.35) -- (7.3,0.35);
\node[font=\small,popgreen!60!popdark] at (5.4,0.65) {distance de sécurité};
% Croix sur fenêtre proche
\draw[poppink,very thick] (7.4,1.3) -- (8.6,2.5);
\draw[poppink,very thick] (8.6,1.3) -- (7.4,2.5);
\node[font=\small,poppink,align=center] at (9.2,2.9) {loin des fenêtres\\ et des portes};
% Ventilation flèches
\draw[->,popblue,thick] (0.2,2.2) -- (1.2,2.2);
\draw[->,popblue,thick] (0.2,1.7) -- (1.2,1.7);
\node[font=\small,popblue] at (0.9,2.6) {air libre};
\end{tikzpicture}
\end{center}
\legende{Bonnes pratiques d'installation d'un groupe électrogène : à l'extérieur, sur un sol stable, dans un endroit ventilé, à distance des portes et des fenêtres pour que les gaz d'échappement (dont le \ce{CO}) ne pénètrent pas dans la maison.}
\end{popfigure}

\begin{exercice}[À toi de jouer]
\begin{enumerate}
\item Pourquoi ne faut-il jamais remplir le réservoir d'une moto dont le moteur vient de tourner ?
\item Cite deux propriétés du monoxyde de carbone qui le rendent particulièrement dangereux.
\item Le père de Marc fait tourner le groupe électrogène dans le couloir de la maison pendant une coupure de courant nocturne. Explique-lui le danger et propose-lui une solution.
\item Cite trois gestes d'entretien qui réduisent à la fois la consommation et la pollution d'une moto.
\end{enumerate}
\end{exercice}

\begin{corrige}[Exercice]
\begin{enumerate}
\item Le moteur et le pot d'échappement sont encore brûlants : une goutte d'essence ou ses vapeurs peuvent s'enflammer au contact des surfaces chaudes. On arrête le moteur et on le laisse refroidir avant le ravitaillement.
\item Le \ce{CO} est \textbf{invisible} (sans couleur) et \textbf{inodore} (sans odeur) : on ne peut pas le détecter avec nos sens, et il tue pendant le sommeil.
\item Le couloir est un espace fermé : le \ce{CO} s'y accumule et peut intoxiquer mortellement toute la famille pendant la nuit. Solution : installer le groupe \textbf{dehors}, dans un endroit ventilé, loin des fenêtres, et le protéger du vol par une chaîne et un cadenas.
\item Vidanger l'huile régulièrement, nettoyer le filtre à air, vérifier et nettoyer la bougie (et faire régler le moteur par un mécanicien).
\end{enumerate}
\end{corrige}

\begin{aretenir}[L'essentiel de la section]
\begin{itemize}
\item L'essence et le gasoil sont \textbf{inflammables} : stockage fermé, à l'ombre, loin des flammes ; jamais de ravitaillement moteur chaud ou en marche.
\item Les gaz d'échappement contiennent le \textbf{monoxyde de carbone \ce{CO}}, gaz toxique invisible et inodore : groupe électrogène et moteur \textbf{toujours à l'extérieur}, jamais dans une pièce fermée.
\item Attention aux \textbf{pièces en rotation} (courroies) et aux \textbf{surfaces brûlantes} (pot d'échappement, radiateur sous pression).
\item L'entretien de base (vidange, filtre à air, bougie, niveaux) fait durer le moteur, économise le carburant et réduit la pollution.
\item Les moteurs polluent l'air (\ce{CO}, \ce{CO2}, fumées noires) et par le bruit : adoptons les gestes éco-citoyens — couper le moteur à l'arrêt prolongé, entretenir, ne jamais jeter l'huile usée au sol.
\end{itemize}
\end{aretenir}

\newpage

\coursec{Activité d'intégration}

\coursec{Généralités sur les moteurs à combustion interne}

\courssub{Objectifs de la leçon}

À l'issue de cette leçon, tu seras capable de :
\begin{itemize}
\item identifier et nommer les \cle{organes principaux} d'un moteur à combustion interne sur un schéma en coupe ;
\item déterminer, pour une machine réelle, son \cle{carburant}, son \cle{cycle} (2 temps ou 4 temps), son \cle{mode d'allumage} (par bougie ou par compression) et son \cle{mode de refroidissement} (à air ou à eau) ;
\item \cle{classer} des moteurs à combustion interne selon plusieurs critères ;
\item rédiger un \cle{conseil de sécurité} adapté à chaque utilisateur ;
\item argumenter un choix technologique (moteur pour un moulin à maïs) en comparant avantages et inconvénients.
\end{itemize}

\courssub{Situation déclenchante : Au garage de Monsieur Etoundi}

Monsieur Etoundi tient un petit garage de quartier à Bafoussam. Ce samedi matin, trois clients se présentent en même temps, et son apprenti, en classe de 3ème au collège voisin, doit l'aider à préparer les fiches de diagnostic. Voici les trois machines déposées au garage.

\textbf{Machine 1 — la moto-taxi de Bouba.} Bouba est conducteur de moto-taxi. Sa moto est équipée d'un moteur à essence \cle{monocylindre} de cylindrée \SI{125}{cm^3}, à \cle{4 temps}, avec \cle{allumage par bougie} et \cle{refroidissement à air} (ailettes sur le cylindre). Consommation moyenne : \SI{2,5}{L} d'essence aux \SI{100}{km}.

\textbf{Machine 2 — le camion de Mama Ngo.} Mama Ngo transporte du sable entre Foumban et Bafoussam. Son camion possède un moteur \cle{diesel} à \cle{6 cylindres}, à \cle{4 temps}, avec \cle{allumage par compression} (sans bougie) et \cle{refroidissement à eau} (radiateur). Consommation moyenne : \SI{18}{L} de gasoil aux \SI{100}{km}.

\textbf{Machine 3 — le moteur hors-bord de Pa Pierre.} Pa Pierre est pêcheur sur le fleuve Noun. Son moteur hors-bord fonctionne à l'essence, c'est un moteur à \cle{2 temps} qui exige un \cle{mélange essence-huile} (4 \% d'huile, soit \SI{40}{cL} d'huile pour \SI{10}{L} d'essence), avec allumage par bougie et refroidissement par l'eau du fleuve.

L'apprenti dispose du schéma en coupe d'un moteur à essence 4 temps (figure ci-dessous) et du tableau de caractéristiques à compléter.

\begin{popfigure}
\begin{center}
\begin{tikzpicture}[scale=0.9,>=Stealth]
% Cylindre
\draw[thick,fill=popblueL] (0,0) rectangle (4,4.5);
\draw[thick,fill=white] (1,0.5) rectangle (3,4.5);
% Piston (position mi-course pour voir la chambre)
\draw[thick,fill=poporangeL] (1,1.2) rectangle (3,2.1);
\node at (2,1.65) {\small piston};
% Bielle et vilebrequin
\draw[thick] (2,1.2) -- (2.6,-0.4);
\draw[thick,fill=popgoldL] (2.6,-0.4) circle (0.18);
\draw[thick] (2,-0.9) circle (0.7);
\fill (2,-0.9) circle (0.06);
\node[below] at (2,-1.7) {\small vilebrequin};
% Chambre de combustion
\draw[thick,fill=poppinkL] (1,3.4) rectangle (3,4.5);
\node at (2,3.95) {\small chambre de combustion};
% Bougie
\draw[thick,fill=popgreenL] (1.8,4.5) rectangle (2.2,5.1);
\draw[thick] (2,5.1) -- (2,5.5);
\node[right] at (2.3,5.2) {\small bougie};
% Soupapes
\draw[thick,fill=popmuted!40] (0.7,4.5) rectangle (1.3,4.9);
\draw[thick] (1,4.9) -- (1,5.6);
\draw[thick,fill=popmuted!40] (2.7,4.5) rectangle (3.3,4.9);
\draw[thick] (3,4.9) -- (3,5.6);
\node[left] at (0.6,5.2) {\small soupape d'admission};
\node[right] at (3.4,5.2) {\small soupape d'échappement};
% Flèches entrées/sorties
\draw[->,thick,popblue] (-1.2,4.7) -- (0.6,4.7) node[midway,above]{\small air + essence};
\draw[->,thick,popdark] (3.4,4.7) -- (5.2,4.7) node[midway,above]{\small gaz brûlés};
% Cotes utiles
\draw[|<->|,thin,popmuted] (4.2,0.5) -- (4.2,4.5) node[midway,right] {\small course};
\draw[|<->|,thin,popmuted] (1,0.2) -- (3,0.2) node[midway,below] {\small alésage};
\end{tikzpicture}
\end{center}
\legende{Schéma en coupe simplifié d'un cylindre de moteur à essence 4 temps : piston, bielle, vilebrequin, chambre de combustion, bougie, soupapes d'admission et d'échappement.}
\end{popfigure}

\courssub{Consignes de l'activité}

\begin{enumerate}
\item Sur le schéma en coupe, nomme chaque organe repéré et donne son rôle en une phrase.
\item Pour chacune des trois machines du garage, indique : le carburant utilisé, le cycle (2 temps ou 4 temps), le mode d'allumage et le mode de refroidissement. Justifie chaque réponse par une information du texte.
\item Recopie et complète le tableau de classement suivant :
\begin{center}
\begin{tabularx}{\linewidth}{|l|X|X|X|}
\hline
\rowcolor{popblueL} \textbf{Critère} & \textbf{Moto-taxi} & \textbf{Camion} & \textbf{Hors-bord} \\
\hline
Carburant & & & \\
\hline
Cycle (2 ou 4 temps) & & & \\
\hline
Mode d'allumage & & & \\
\hline
Mode de refroidissement & & & \\
\hline
Nombre de cylindres & & & \\
\hline
\end{tabularx}
\end{center}
\item Calcule la quantité d'huile que Pa Pierre doit ajouter à \SI{25}{L} d'essence pour respecter le mélange à 4 \%.
\item Rédige un conseil de sécurité (2 lignes maximum) pour chaque utilisateur : Bouba (moto), Mama Ngo (camion), Pa Pierre (hors-bord).
\item \textbf{Débat.} Le village de Pa Pierre, non électrifié, veut acheter un moteur pour entraîner un moulin à maïs qui fonctionnera 6 heures par jour. Faut-il choisir un petit moteur essence 4 temps, un moteur diesel ou un moteur 2 temps ? Argumente en comparant au moins deux avantages et deux inconvénients de chaque solution.
\end{enumerate}

\courssub{Ressources à mobiliser : Rappels de cours}

\begin{definition}[Moteur à combustion interne]
Un \cle{moteur à combustion interne} est une machine thermique qui transforme l'\cle{énergie chimique} d'un carburant (essence, gasoil) en \cle{énergie mécanique} (rotation), la combustion ayant lieu \emph{à l'intérieur} du moteur, dans la \cle{chambre de combustion}, contrairement au moteur à combustion externe (machine à vapeur).
\end{definition}

\begin{propriete}[Organes principaux d'un moteur alternatif]
\begin{itemize}
\item \cle{Cylindre} : tube métallique (souvent en fonte ou aluminium) dans lequel coulisse le piston.
\item \cle{Piston} : pièce mobile étanche qui comprime le mélange et reçoit la poussée des gaz brûlés ; il assure la liaison avec la bielle.
\item \cle{Bielle} : pièce de liaison qui transmet la force du piston au vilebrequin ; elle transforme le mouvement alternatif en rotation.
\item \cle{Vilebrequin} (ou arbre à cames) : arbre muni de manetons excentrés ; il reçoit le mouvement alternatif de la bielle et délivre un mouvement de rotation continu.
\item \cle{Chambre de combustion} : volume variable situé au-dessus du piston (culasse + tête de piston) où se produit la combustion.
\item \cle{Soupapes} : clapets commandés par l'arbre à cames (via poussoirs/culbuteurs) : la \cle{soupape d'admission} laisse entrer le mélange frais (ou l'air), la \cle{soupape d'échappement} laisse sortir les gaz brûlés.
\item \cle{Bougie d'allumage} (moteur à essence) : produit l'étincelle haute tension qui enflamme le mélange compressé.
\item \cle{Injecteur} (moteur diesel) : pulvérise le gasoil dans l'air comprimé très chaud.
\end{itemize}
\end{propriete}

\begin{methode}[Le cycle à 4 temps (cycle Beau de Rochas / Otto) -- Démonstration exigible au BEPC]
Le cycle se déroule en \textbf{2 tours de vilebrequin} (720°) et \textbf{4 courses du piston} :
\begin{enumerate}
\item \textbf{Admission (1er temps, descente)} : Soupape d'admission \textbf{ouverte}, échappement fermée. Le piston descend, créant une dépression qui aspire le mélange air-essence (ou l'air seul pour le diesel).
\item \textbf{Compression (2ème temps, montée)} : Les deux soupapes \textbf{fermées}. Le piston monte, comprimant le mélange (taux de compression \SIrange{8}{12}{ pour l'essence}, \SIrange{14}{22}{ pour le diesel}). Température et pression augmentent fortement.
\item \textbf{Combustion -- Détente (3ème temps, descente = \cle{temps moteur})} : Les deux soupapes \textbf{fermées}. 
\begin{itemize}
\item Essence : la bougie déclenche l'inflammation \emph{isochore} (quasi instantanée).
\item Diesel : l'injecteur pulvérise le carburant qui s'enflamme par \emph{auto-inflammation} (compression).
\end{itemize}
La pression monte brutalement (\SIrange{30}{50}{bar}), le piston est violemment repoussé vers le bas : c'est le seul temps qui \emph{produit} de l'énergie.
\item \textbf{Échappement (4ème temps, montée)} : Soupape d'échappement \textbf{ouverte}, admission fermée. Le piston monte et chasse les gaz brûlés vers le pot d'échappement.
\end{enumerate}
\emph{Astuce mnémotechnique :} \textbf{A}dmision -- \textbf{C}ompression -- \textbf{E}xplosion (Détente) -- \textbf{E}chappement (\textbf{A C E E}).
\end{methode}

\begin{propriete}[Moteur à 2 temps et Moteur Diesel]
\begin{itemize}
\item \textbf{Moteur à 2 temps} : Le cycle (admission-compression + combustion-échappement) se réalise en \textbf{1 tour de vilebrequin} (360°). Pas de soupapes mais des \cle{luminaires} (orifices) découverts par le piston. Nécessite souvent un \cle{mélange essence-huile} (lubrification par perte totale). Plus léger, plus simple, mais plus polluant et moins rendement.
\item \textbf{Moteur Diesel (allumage par compression)} : Admet \textbf{uniquement de l'air} à l'admission. Taux de compression très élevé ($\sim$ 20) $\rightarrow$ température $\sim$ \SI{600}{°C} en fin de compression. Injection du gasoil $\rightarrow$ auto-inflammation. Pas de bougie. Rendement meilleur, couple plus élevé, plus robuste, mais plus lourd, plus bruyant, plus cher à l'achat.
\end{itemize}
\end{propriete}

\begin{savaistu}[Histoire et contexte camerounais]
\begin{itemize}
\item \textbf{Nicolaus Otto} (1876) brevette le 4 temps ; \textbf{Rudolf Diesel} (1892) brevette l'allumage par compression.
\item Au Cameroun, la \cle{SONARA} (Société Nationale de Raffinage) de Limbé transforme le brut en essence super, gasoil, kérosène. Les moto-taxis (\emph{benskin}) utilisent massivement des moteurs 4 temps essence (Yamaha, Honda, TVS) pour leur fiabilité et leur faible consommation d'huile.
\item Les groupes électrogènes ENEO ou de secours (hôpitaux, mairies) sont quasi exclusivement \textbf{diesel} (robustesse, autonomie, gasoil moins cher).
\item Les moteurs hors-bord 2 temps (Yamaha, Mercury, Suzuki) équipent la pêche artisanale sur le Wouri, la Sanaga, le Noun, le Lac Tchad : légers, simples, refroidis par l'eau du fleuve.
\end{itemize}
\end{savaistu}

\begin{aretenir}[Synthèse des critères de classement]
\begin{itemize}
\item \textbf{Carburant} : Essence (mélange air/essence) ou Gasoil (air seul + injection) ou Mélange essence-huile (2 temps).
\item \textbf{Cycle} : 4 temps (2 tours vilebrequin, soupapes) ou 2 temps (1 tour, luminaires).
\item \textbf{Allumage} : Par \cle{bougie} (étincelle, essence) ou par \cle{compression} (auto-inflammation, diesel).
\item \textbf{Refroidissement} : \cle{À air} (ailettes sur cylindre/culasse, ventilateur souvent) ou \cle{À eau} (circuit liquide : pompe, radiateur, calorstat, ventilateur).
\item \textbf{Architecture} : \cle{Monocylindre} (1 cylindre), \cle{Bicylindre}, \cle{4 cylindres en ligne}, \cle{V6}, \cle{V8}... (nombre et disposition des cylindres).
\item \textbf{Proportionnalité (mélange 2 temps)} : Pour un pourcentage $p$ \%, $V_{\text{huile}} = \frac{p}{100} \times V_{\text{essence}}$ (mêmes unités).
\item \textbf{Équation de combustion simplifiée (essence $\approx$ octane)} : 
\[ \ce{2 C8H18 + 25 O2 -> 16 CO2 + 18 H2O + energie} \]
\end{itemize}
\end{aretenir}

\begin{attention}[Pièges fréquents et consignes de sécurité]
\begin{itemize}
\item \textbf{Confusion 2 temps / 4 temps} : Le 2 temps n'a \textbf{pas de soupapes} (luminaires) et \textbf{pas de carter d'huile séparé} (huile dans l'essence). Ne pas écrire « bougie » pour un diesel !
\item \textbf{Refroidissement} : « Refroidi par air » $\neq$ « pas de refroidissement ». Les ailettes sont indispensables. Ne \textbf{jamais} ouvrir un vase d'expansion ou bouchon de radiateur \textbf{à chaud} (risque de géyser d'eau bouillante \SI{>100}{°C} sous pression).
\item \textbf{Calcul mélange} : 4 \% = 0,04 = \SI{40}{cL}/\SI{10}{L} = \SI{4}{cL}/\SI{1}{L}. Vérifie l'unité demandée (L ou cL).
\item \textbf{Sécurité carburant} : Essence = très inflammable (vapeurs lourdes). Plein \textbf{moteur arrêté et froid}. Pas de téléphone, pas de cigarette. Jerrican homologué (rouge essence, jaune gasoil).
\item \textbf{Unités} : Cylindrée en \si{cm^3} ou \si{L} (\SI{125}{cm^3} = \SI{0,125}{L}). Consommation en \si{L/100 km}.
\end{itemize}
\end{attention}

\courssub{Démarche et corrigé commenté}

\begin{exoresolu}[Diagnostic au garage : corrigé complet]
Identifier les organes, classer les trois moteurs, calculer le mélange du hors-bord, rédiger les conseils de sécurité et trancher le débat du moulin à maïs.
\tcblower
\textbf{Question 1 -- Identification des organes (schéma en coupe).}
\begin{itemize}
\item \cle{Cylindre} : tube fixe dans lequel coulisse le piston ; paroi refroidie (ailettes ou chemise d'eau).
\item \cle{Piston} : pièce mobile qui assure l'étanchéité, comprime le mélange et reçoit la poussée des gaz ; transmet l'effort à la bielle.
\item \cle{Bielle} : organe de liaison piston -- vilebrequin ; transforme le mouvement alternatif (va-et-vient) en rotation.
\item \cle{Vilebrequin} : arbre tournant qui reçoit les efforts des bielles et délivre le couple moteur vers la boîte de vitesses / roue.
\item \cle{Chambre de combustion} : volume variable entre tête de piston et culasse ; siège de la combustion.
\item \cle{Bougie} : organe d'allumage (moteur essence) ; produit l'étincelle entre ses électrodes sous haute tension (\SIrange{15}{30}{kV}).
\item \cle{Soupape d'admission} : commande l'entrée du mélange frais (air+essence) ; \cle{soupape d'échappement} : commande la sortie des gaz brûlés.
\end{itemize}

\textbf{Question 2 -- Analyse de chaque machine (justifiée par le texte).}
\begin{itemize}
\item \textbf{Moto-taxi (Bouba)} : Carburant = \textbf{Essence} (texte : « moteur à essence ») ; Cycle = \textbf{4 temps} (texte) ; Allumage = \textbf{Par bougie} (texte : « allumage par bougie ») ; Refroidissement = \textbf{À air} (texte : « ailettes sur le cylindre » = signature du refroidissement air).
\item \textbf{Camion (Mama Ngo)} : Carburant = \textbf{Gasoil} (texte : « moteur diesel ») ; Cycle = \textbf{4 temps} (texte) ; Allumage = \textbf{Par compression} (texte : « allumage par compression (sans bougie) », définition du diesel) ; Refroidissement = \textbf{À eau} (texte : « radiateur » = signature du refroidissement eau).
\item \textbf{Hors-bord (Pa Pierre)} : Carburant = \textbf{Mélange essence-huile} (texte : « mélange essence-huile 4 \% ») ; Cycle = \textbf{2 temps} (texte) ; Allumage = \textbf{Par bougie} (texte : « allumage par bougie ») ; Refroidissement = \textbf{À eau} (texte : « refroidissement par l'eau du fleuve » = refroidissement direct par eau de mer/rivière, sans radiateur).
\end{itemize}

\textbf{Question 3 -- Tableau de classement complété.}
\begin{center}
\begin{tabularx}{\linewidth}{|l|X|X|X|}
\hline
\rowcolor{popblueL} \textbf{Critère} & \textbf{Moto-taxi} & \textbf{Camion} & \textbf{Hors-bord} \\
\hline
Carburant & Essence & Gasoil & Mélange essence-huile (4 \%) \\
\hline
Cycle (2 ou 4 temps) & 4 temps & 4 temps & 2 temps \\
\hline
Mode d'allumage & Par bougie & Par compression & Par bougie \\
\hline
Mode de refroidissement & À air (ailettes) & À eau (radiateur) & À eau (fleuve, refroidissement direct) \\
\hline
Nombre de cylindres & 1 (monocylindre) & 6 & 1 (généralement monocylindre pour cette gamme) \\
\hline
\end{tabularx}
\end{center}

\textbf{Question 4 -- Calcul du mélange essence-huile (proportionnalité).}
Le mélange contient 4 \% d'huile en volume.
\begin{align*}
V_{\text{huile}} &= \frac{4}{100} \times V_{\text{essence}} \\
&= 0{,}04 \times \SI{25}{L} \\
&= \SI{1}{L}
\end{align*}
\textbf{Vérification} : \SI{40}{cL} pour \SI{10}{L} $\rightarrow$ \SI{4}{cL} par litre. Pour \SI{25}{L} : $25 \times \SI{4}{cL} = \SI{100}{cL} = \SI{1}{L}$.
\textbf{Réponse} : Pa Pierre doit ajouter \textbf{\SI{1}{L} d'huile} à ses \SI{25}{L} d'essence.

\textbf{Question 5 -- Conseils de sécurité (courts, impératifs, adaptés).}
\begin{itemize}
\item \textbf{Bouba (moto essence 4 temps, refroidi air)} : « Ne fais jamais le plein moteur chaud ni en marche ; nettoie régulièrement les ailettes de refroidissement et contrôle le niveau d'huile moteur tous les 500 km. »
\item \textbf{Mama Ngo (camion diesel, refroidi eau)} : « Vérifie le niveau de liquide de refroidissement moteur froid avant chaque départ ; n'ouvre \textbf{jamais} le bouchon du radiateur ou du vase d'expansion à chaud (brûlures graves). »
\item \textbf{Pa Pierre (hors-bord 2 temps, mélange)} : « Prépare le mélange essence-huile \textbf{à l'air libre}, loin de toute flamme ou étincelle ; transporte le carburant dans un jerrican \textbf{homologué} (rouge), bien fermé, arrimé dans la pirogue. »
\end{itemize}

\textbf{Question 6 -- Débat : Quel moteur pour le moulin à maïs du village (6 h/jour) ?}
\begin{center}
\begin{tabularx}{\linewidth}{|l|X|X|}
\hline
\rowcolor{popblueL} \textbf{Solution} & \textbf{Avantages (min 2)} & \textbf{Inconvénients (min 2)} \\
\hline
\textbf{Moteur essence 4 temps} & Démarrage facile (lanceur/électrique) ; entretien simple (vidange, bougie, filtre) ; peu de fumée/odeur ; prix d'achat modéré. & Essence \textbf{chère} et parfois \textbf{rare} en zone rurale (ruptures SONARA) ; consommation élevée sur 6 h/jour (coût horaire fort) ; durée de vie moindre sous charge continue.
\\ \hline
\textbf{Moteur diesel} & Gasoil \textbf{moins cher} au litre ; \textbf{rendement meilleur} (conso \SI{200}{g/kWh} vs \SI{300}{g/kWh} essence) $\rightarrow$ autonomie \textbf{2$\times$ sup.} ; \textbf{couple élevé} à bas régime idéal pour moulin ; \textbf{robustesse} (forte compression, pièces massives) $\rightarrow$ durée de vie \textbf{très longue} (20 000 h+). & Prix d'achat \textbf{élevé} (injection, bloc renforcé) ; démarrage plus dur (batterie, préchauffage) ; entretien technique (filtre gasoil, purge, injecteurs) ; bruit et vibrations forts ; gasoil peut geler par grand froid (rare au Cameroun).
\\ \hline
\textbf{Moteur 2 temps} & Très léger, compact ; prix d'achat bas ; mécanique ultra-simple (peu de pannes). & \textbf{Obligation mélange huile-essence} (coût, dosage, pollution) ; \textbf{usure rapide} (lubrification imparfaite) ; \textbf{très polluant} (huile brûlée, gaz non brûlés) ; \textbf{inadapté} à 6 h/jour continu (surchauffe, casse) ; couple faible à bas régime.
\\ \hline
\end{tabularx}
\end{center}

\textbf{Conclusion argumentée} : Pour un moulin à maïs devant tourner \textbf{6 heures par jour} dans un village \textbf{non électrifié}, le \cle{moteur diesel} est le choix technique et économique le plus pertinent. 
\begin{itemize}
\item Son \textbf{coût horaire de carburant} est le plus bas (gasoil moins cher + meilleur rendement).
\item Sa \textbf{robustesse} garantit une durée de vie de plusieurs années en usage intensif, amortissant le surcoût d'achat.
\item Son \textbf{couple élevé} à bas régime évite le calage sous charge (maïs dur, humidité).
\end{itemize}
Le moteur essence 4 temps reste une solution de dépannage (disponibilité immédiate, budget initial très serré). Le moteur 2 temps est \textbf{formellement déconseillé} pour cet usage continu (fiabilité, pollution, coût huile).
\end{exoresolu}

\courssub{Bilan de la leçon}

Cette leçon t'a permis de mettre en évidence que :
\begin{itemize}
\item un moteur à combustion interne se reconnaît par ses \cle{organes} (cylindre, piston, bielle, vilebrequin, soupapes, bougie/injecteur) et par ses \cle{critères de classement} : carburant, cycle, allumage, refroidissement, architecture ;
\item le \cle{carburant} impose le mode d'allumage : essence $\rightarrow$ bougie (inflammation commandée) ; gasoil $\rightarrow$ compression (auto-inflammation) ;
\item le \cle{cycle 4 temps} (admission, compression, combustion-détente, échappement) est le standard pour les véhicules routiers (moto, voiture, camion) par son rendement et sa propreté ; le \cle{cycle 2 temps} persiste pour la légèreté (tronçonneuse, débroussailleuse, petit hors-bord) ;
\item le choix d'un moteur n'est jamais gratuit : il dépend de l'\cle{usage} (durée, charge, environnement), du \cle{coût global} (achat + carburant + entretien sur la durée) et de la \cle{sécurité} ;
\item les règles de sécurité (plein moteur arrêté/froid, radiateur jamais ouvert à chaud, mélange préparé loin des flammes, jerrican homologué) font partie intégrante du métier de technicien responsable ;
\item la technologie des moteurs, apprise en classe, sert directement à résoudre des problèmes concrets de la vie quotidienne au Cameroun : équiper un village en moulin à maïs, entretenir sa moto-taxi, gérer son camion de transport, pêcher sur le fleuve Noun.
\end{itemize}

\courssub{Exercices d'entraînement (BEPC type)}

\begin{exercice}[QCM -- Classification rapide]
Pour chacun des énoncés, coche la bonne réponse.
\begin{enumerate}
\item Un moteur monocylindre 125 cm$^3$, refroidi par ailettes, démarreur au kick, carburant essence : c'est un moteur \ldots
\begin{itemize}
\item[A.] Diesel 4 temps \item[B.] Essence 2 temps \item[C.] Essence 4 temps \item[D.] Diesel 2 temps
\end{itemize}
\item Le taux de compression d'un moteur diesel est typiquement \ldots
\begin{itemize}
\item[A.] 9 : 1 \item[B.] 11 : 1 \item[C.] 18 : 1 \item[D.] 5 : 1
\end{itemize}
\item L'organe qui transforme le mouvement alternatif du piston en rotation est \ldots
\begin{itemize}
\item[A.] La bielle \item[B.] Le vilebrequin \item[C.] L'arbre à cames \item[D.] Le volant moteur
\end{itemize}
\end{enumerate}
\end{exercice}

\begin{exercice}[Calcul de mélange et consommation]
Un débroussailleuse 2 temps fonctionne avec un mélange à 3 \% d'huile.
\begin{enumerate}
\item Quelle volume d'huile pour \SI{5}{L} d'essence ? (en cL et en L)
\item Si le réservoir fait \SI{0,75}{L} de mélange, quelle quantité d'essence pure et d'huile contient-il ?
\item Le moteur consomme \SI{0,6}{L} de mélange par heure. Combien d'heures de travail avec \SI{10}{L} d'essence pure (en préparant le mélange) ?
\end{enumerate}
\end{exercice}

\begin{exercice}[Analyse de cycle -- Schéma à compléter]
On donne le diagramme pression-volume (diagramme de Clapeyron) simplifié d'un cycle 4 temps essence.
\begin{center}
\begin{tikzpicture}[scale=0.7,>=Stealth]
\draw[->] (0,0) -- (7,0) node[right] {$V$};
\draw[->] (0,0) -- (0,5) node[above] {$P$};
% Admission (approx horizontal bas)
\draw[thick,popblue] (1,0.5) -- (6,0.5) node[midway,above] {1};
% Compression (courbe montante)
\draw[thick,poporange] (6,0.5) .. controls (5,1) and (3,2) .. (2,3) node[midway,right] {2};
% Combustion (isochore verticale)
\draw[thick,popdark] (2,3) -- (2,4.5) node[midway,right] {3};
% Détente (courbe descendante)
\draw[thick,popgreen] (2,4.5) .. controls (3,3.5) and (5,1.5) .. (6,0.7) node[midway,right] {4};
% Echappement (isochore descendante + horizontal)
\draw[thick,poppurple] (6,0.7) -- (6,0.5) node[midway,right] {5};
\draw[thick,poppurple] (6,0.5) -- (1,0.5);
% Point mort haut/bas
\node[below] at (1,0) {PMH};
\node[below] at (6,0) {PMB};
\end{tikzpicture}
\end{center}
\legende{Diagramme $P(V)$ qualitatif d'un cycle 4 temps essence.}
\begin{enumerate}
\item Nomme les 5 transformations numérotées 1 à 5.
\item Sur quelles transformations les soupapes sont-elles ouvertes ? Lesquelles ?
\item Pourquoi la transformation 3 est-elle verticale (isochore) ?
\item Quelle transformation correspond au \emph{temps moteur} ? Justifie.
\end{enumerate}
\end{exercice}

\begin{corrige}[Exercices d'entraînement]
\textbf{Exercice 1 (QCM)} : 1. \textbf{C} (Essence 4 temps : ailettes = air, essence = bougie, kick = 4 temps généralement). 2. \textbf{C} (Diesel : 14 à 22, typique 17-18). 3. \textbf{B} (Vilebrequin : bielle transmet, vilebrequin transforme par excentrique).
\textbf{Exercice 2 (Mélange)} : 1. $V_h = 0,03 \times 5 = \SI{0,15}{L} = \SI{15}{cL}$. 2. Mélange total 0,75 L. $V_h = 0,03 \times 0,75 = \SI{0,0225}{L} = \SI{2,25}{cL}$ ; $V_{es} = 0,75 - 0,0225 = \SI{0,7275}{L}$. 3. 10 L essence $\rightarrow$ mélange total $V_{mix} = V_{es} / 0,97 = 10 / 0,97 \approx \SI{10,31}{L}$. Durée $= 10,31 / 0,6 \approx \SI{17,2}{h}$.
\textbf{Exercice 3 (Diagramme)} : 1. 1=Admission (isobare), 2=Compression (adiabatique), 3=Combustion (isochore), 4=Détente (adiabatique), 5=Échappement (isochore + isobare). 2. Soupapes ouvertes : Admission (1) $\rightarrow$ soupape admission ; Échappement (5) $\rightarrow$ soupape échappement. 3. Combustion quasi instantanée $\rightarrow$ volume constant (piston au PMH) $\rightarrow$ isochore. 4. Transformation 4 (Détente) : seule où la pression pousse le piston vers le bas (travail moteur $>0$).
\end{corrige}

\newpage

\coursec{Méthodes et exercices résolus}

\coursec{Généralités sur les moteurs à combustion interne}

\courssub{Fiches méthodes et exercices résolus}

\begin{methode}[Identifier un moteur à combustion interne]
Pour reconnaître et définir un moteur à combustion interne dans une situation concrète :
\begin{enumerate}
\item Repérer la \cle{source d'énergie} : un carburant (essence, gasoil, gaz) qui brûle.
\item Vérifier le lieu de la combustion : si le combustible brûle \textbf{à l'intérieur} de la machine (dans le cylindre), c'est un moteur à \cle{combustion interne} ; si la chaleur vient de l'extérieur (chaudière), c'est une machine à combustion externe.
\item Identifier la transformation d'énergie : énergie \textbf{chimique} $\rightarrow$ énergie \textbf{thermique} $\rightarrow$ énergie \textbf{mécanique} (mouvement).
\item Citer des exemples de la vie courante : moteur de moto-taxi, groupe électrogène, moteur de voiture, tondeuse à gazon.
\end{enumerate}
\end{methode}

\begin{exoresolu}[Reconnaître un moteur à combustion interne]
Au quartier, trois machines fonctionnent : une moto-taxi, un groupe électrogène et une machine à vapeur de scierie dont la chaudière est chauffée au bois. Dire, en justifiant, lesquelles sont des moteurs à combustion interne.
\tcblower
\textbf{Raisonnement :} un moteur à combustion interne est une machine dans laquelle le combustible brûle \textbf{à l'intérieur} du moteur (dans le cylindre), transformant l'énergie chimique en énergie mécanique.
\begin{itemize}
\item \textbf{Moto-taxi :} l'essence brûle dans le cylindre du moteur. C'est un moteur à combustion interne.
\item \textbf{Groupe électrogène :} le gasoil (ou l'essence) brûle dans le cylindre du moteur qui entraîne l'alternateur. C'est un moteur à combustion interne.
\item \textbf{Machine à vapeur :} le bois brûle \textbf{à l'extérieur}, dans une chaudière séparée ; la vapeur produite actionne ensuite la machine. Ce n'est \textbf{pas} un moteur à combustion interne, mais une machine à combustion externe.
\end{itemize}
\textbf{Conclusion :} la moto-taxi et le groupe électrogène sont des moteurs à combustion interne ; la machine à vapeur est une machine à combustion externe.
\end{exoresolu}

\begin{experience}[Observer la constitution d'un moteur 4 temps (modèle réduit ou schéma éclaté)]
\textbf{Matériel :} modèle réduit transparent de moteur 4 temps (ou schéma éclaté projeté), fiche d'observation.
\textbf{Protocole :}
\begin{enumerate}
\item Identifier le \cle{cylindre} (corps fixe), le \cle{piston} (pièce mobile coulissant dans le cylindre), la \cle{bielle} (tige reliant piston et vilebrequin), le \cle{vilebrequin} (arbre à manivelle), les \cle{soupapes} (admission et échappement) et la \cle{bougie} (allumage).
\item Actionner manuellement le vilebrequin et observer le mouvement alternatif du piston (montée/descente) et l'ouverture/fermeture des soupapes.
\item Repérer le \cle{point mort haut} (PMH : piston tout en haut) et le \cle{point mort bas} (PMB : piston tout en bas).
\end{enumerate}
\textbf{Observations :} le piston effectue un aller-retour (2 courses) pendant 1 tour de vilebrequin. Les soupapes s'ouvrent et se ferment une fois chacun par cycle complet (2 tours).
\textbf{Interprétation :} l'ensemble \textbf{bielle -- vilebrequin} transforme le mouvement alternatif rectiligne du piston en mouvement de rotation continu utilisable par les roues ou l'alternateur.
\legende{Modèle réduit de moteur 4 temps : 1. Cylindre, 2. Piston, 3. Bielle, 4. Vilebrequin, 5. Soupape admission, 6. Soupape échappement, 7. Bougie.}
\end{experience}

\begin{methode}[Décrire la constitution d'un moteur à combustion interne]
Pour citer et situer les organes essentiels d'un moteur :
\begin{enumerate}
\item Apprendre la liste des organes : \cle{cylindre}, \cle{piston}, \cle{bielle}, \cle{vilebrequin}, \cle{soupapes} (admission et échappement), \cle{bougie} (moteur à essence), \cle{injecteur} ou carburateur.
\item Retenir le rôle de chacun : le piston coulisse dans le cylindre ; la bielle relie le piston au vilebrequin ; le vilebrequin transforme le mouvement alternatif en mouvement de \textbf{rotation} ; les soupapes commandent l'entrée des gaz frais et la sortie des gaz brûlés ; la bougie produit l'étincelle d'allumage.
\item Sur un schéma, flécher chaque organe sans inverser bielle et vilebrequin.
\end{enumerate}
\end{methode}

\begin{exoresolu}[Identifier les organes sur un schéma]
Sur le schéma d'un moteur à essence, quatre organes sont repérés par les lettres A (pièce qui coulisse dans le cylindre), B (tige reliant A à l'arbre), C (arbre qui tourne), D (pièce produisant l'étincelle). Nommer chaque organe et donner son rôle.
\tcblower
\textbf{Raisonnement :} on utilise la description fonctionnelle de chaque organe.
\begin{itemize}
\item \textbf{A = piston :} il coulisse dans le cylindre et reçoit la poussée des gaz brûlés.
\item \textbf{B = bielle :} elle relie le piston au vilebrequin et transmet l'effort.
\item \textbf{C = vilebrequin :} il transforme le mouvement alternatif (aller-retour) du piston en mouvement de rotation.
\item \textbf{D = bougie :} elle produit l'étincelle qui enflamme le mélange air-essence.
\end{itemize}
\textbf{Conclusion :} A est le piston, B la bielle, C le vilebrequin et D la bougie. L'ensemble bielle-vilebrequin est le système qui convertit la translation du piston en rotation utilisable par les roues.
\end{exoresolu}

\begin{methode}[Décrire le cycle à 4 temps]
Pour restituer les 4 temps d'un moteur à essence dans l'ordre :
\begin{enumerate}
\item Retenir le moyen mnémotechnique : \textbf{A}dmission – \textbf{C}ompression – \textbf{D}étente (explosion) – \textbf{É}chappement.
\item Préciser pour chaque temps : la position des soupapes (ouvertes ou fermées) et le mouvement du piston (descente ou montée).
\item Identifier le seul temps \textbf{moteur} : la détente (temps 3), qui fournit le travail ; les trois autres temps sont résistants.
\item Savoir que 4 temps = 2 tours de vilebrequin = 4 courses du piston.
\end{enumerate}
\end{methode}

\begin{exoresolu}[Analyser le cycle à 4 temps]
Un élève affirme : « Pendant la compression, la soupape d'admission est ouverte et le piston descend. » 1) Corriger cette affirmation. 2) Citer le temps moteur et justifier. 3) Combien de tours de vilebrequin effectue-t-on pendant un cycle complet ?
\tcblower
\textbf{1) Correction :} pendant la \cle{compression}, les \textbf{deux soupapes sont fermées} et le piston \textbf{monte}, comprimant le mélange air-essence admis au temps précédent. L'affirmation est donc doublement fausse.
\textbf{2) Temps moteur :} c'est la \cle{détente} (3\textsuperscript{e} temps) : la bougie enflamme le mélange comprimé, les gaz brûlés se détendent brutalement et \textbf{poussent le piston vers le bas} : c'est le seul temps qui fournit du travail mécanique. Les trois autres temps consomment de l'énergie.
\textbf{3) Nombre de tours :} un cycle à 4 temps comporte 4 courses du piston, et chaque course correspond à un demi-tour de vilebrequin :
\[ 4 \text{ courses} \times 0{,}5 \text{ tour/course} = 2 \text{ tours de vilebrequin.} \]
\textbf{Conclusion :} le cycle complet s'effectue en \textbf{2 tours de vilebrequin}, avec un seul temps moteur (la détente).
\end{exoresolu}

\begin{savaistu}[Le cycle 4 temps au Cameroun]
La quasi-totalité des moto-taxis (\emph{« bend skin »}) et des voitures particulières au Cameroun utilisent des moteurs 4 temps à essence. Ils sont préférés aux 2 temps pour leur rendement meilleur, leur consommation d'huile nulle (carter d'huile séparé) et leurs émissions de fumées moins visibles. Les groupes électrogènes domestiques (marques Elemax, Honda, etc.) sont aussi majoritairement des 4 temps essence ou diesel.
\end{savaistu}

\begin{methode}[Distinguer moteur 2 temps, moteur 4 temps et moteur diesel]
Pour classer un moteur à partir d'indices concrets :
\begin{enumerate}
\item \textbf{Moteur 2 temps} : le cycle se fait en \textbf{1 tour} de vilebrequin (2 courses) ; mélange essence-huile ; moteur léger (tronçonneuse, anciennes motos, débroussailleuses) ; plus polluant (fumées bleues).
\item \textbf{Moteur 4 temps à essence} : cycle en 2 tours ; allumage par \textbf{bougie} (étincelle) ; voitures et motos courantes ; lubrification par carter d'huile.
\item \textbf{Moteur diesel} : allumage par \textbf{compression} (pas de bougie, mais des injecteurs) ; carburant = gasoil ; taux de compression élevé ; couple élevé ; moteurs robustes des camions, bus et gros groupes électrogènes.
\item Réflexe : « bougie $\rightarrow$ essence » ; « injection + forte compression $\rightarrow$ diesel ».
\end{enumerate}
\end{methode}

\begin{exoresolu}[Classer des moteurs]
Trois engins : (a) une tronçonneuse dont le réservoir contient un mélange essence-huile, (b) une voiture à essence équipée de bougies d'allumage, (c) un camion fonctionnant au gasoil sans bougies. Classer chaque moteur et justifier.
\tcblower
\textbf{Raisonnement :} on utilise les critères carburant, mode d'allumage et nombre de temps.
\begin{itemize}
\item \textbf{(a) Tronçonneuse :} mélange essence-huile et moteur léger $\Rightarrow$ moteur à essence à \textbf{2 temps}. L'huile mélangée à l'essence assure la lubrification, car il n'y a pas de carter d'huile séparé.
\item \textbf{(b) Voiture à essence :} présence de bougies $\Rightarrow$ allumage commandé par étincelle $\Rightarrow$ moteur à essence à \textbf{4 temps}.
\item \textbf{(c) Camion au gasoil :} pas de bougies, allumage spontané du gasoil par forte compression de l'air $\Rightarrow$ moteur \textbf{diesel} (à injection).
\end{itemize}
\textbf{Conclusion :} (a) moteur 2 temps à essence ; (b) moteur 4 temps à essence ; (c) moteur diesel.
\end{exoresolu}

\begin{methode}[Classer les moteurs à combustion interne]
Pour présenter un classement complet et rigoureux :
\begin{enumerate}
\item Choisir un \textbf{critère de classement} et l'annoncer explicitement.
\item Critères usuels :
\begin{itemize}
\item selon le \textbf{nombre de temps} : moteurs 2 temps / 4 temps ;
\item selon le \textbf{carburant} : essence / diesel (gasoil) / gaz (GPL/GNV) ;
\item selon le \textbf{mode d'allumage} : allumage commandé (bougie) / allumage par compression ;
\item selon la \textbf{disposition des cylindres} : en ligne, en V, à plat (boxer).
\end{itemize}
\item Ne jamais mélanger deux critères dans une même case du tableau.
\end{enumerate}
\end{methode}

\begin{exoresolu}[Construire un tableau de classement]
Un moteur de moto-taxi est à essence, 4 temps, un cylindre, allumage par bougie. Un moteur de bus est au gasoil, 4 temps, 6 cylindres en ligne, allumage par compression. Présenter ces informations dans un tableau de classement selon trois critères.
\tcblower
\textbf{Raisonnement :} on choisit trois critères : carburant, nombre de temps, mode d'allumage (avec le nombre de cylindres en complément).
\begin{center}
\begin{tabularx}{\linewidth}{|l|X|X|}\hline
\rowcolor{popblueL} \textbf{Critère} & \textbf{Moto-taxi} & \textbf{Bus} \\ \hline
Carburant & Essence & Gasoil (diesel) \\ \hline
Nombre de temps & 4 temps & 4 temps \\ \hline
Mode d'allumage & Commandé (bougie) & Par compression \\ \hline
Cylindres & 1 cylindre & 6 cylindres en ligne \\ \hline
\end{tabularx}
\end{center}
\textbf{Conclusion :} les deux moteurs sont à 4 temps, mais ils diffèrent par le carburant, le mode d'allumage et le nombre de cylindres. Chaque ligne du tableau correspond à un seul critère : le classement est cohérent.
\end{exoresolu}

\begin{methode}[Appliquer les consignes de sécurité et d'entretien]
Pour répondre à une question de sécurité liée aux moteurs :
\begin{enumerate}
\item \textbf{Carburant} : stocker l'essence loin des flammes, dans un bidon homologué ; ne jamais faire le plein moteur chaud et en marche.
\item \textbf{Gaz d'échappement} : ne jamais faire tourner un groupe électrogène dans une pièce fermée (risque d'intoxication au monoxyde de carbone \ce{CO}, gaz invisible, inodore et mortel).
\item \textbf{Entretien} : vidange régulière de l'huile (tous les 5 000 à 10 000 km ou selon notice), nettoyage du filtre à air, vérification/remplacement des bougies (essence) ou du filtre à gasoil (diesel).
\item \textbf{Environnement} : les gaz d'échappement (\ce{CO2}, \ce{CO}, \ce{NOx}, particules fines) polluent l'air ; un moteur bien réglé consomme moins et pollue moins.
\end{enumerate}
\end{methode}

\begin{exoresolu}[Sécurité autour d'un groupe électrogène]
Pendant une coupure d'ENEO, une famille installe son groupe électrogène dans la cuisine fermée pour éviter le vol. Le lendemain, deux personnes se plaignent de maux de tête et de nausées. 1) Expliquer la cause probable. 2) Donner deux consignes de sécurité et une mesure d'entretien.
\tcblower
\textbf{1) Cause :} la combustion du carburant dans le moteur produit des gaz d'échappement contenant du \cle{monoxyde de carbone} \ce{CO}, gaz incolore et inodore. Dans une pièce fermée, le \ce{CO} s'accumule et provoque une intoxication (maux de tête, nausées, pouvant aller jusqu'à la mort).
\textbf{2) Consignes de sécurité :}
\begin{itemize}
\item Installer le groupe électrogène \textbf{à l'extérieur}, dans un endroit bien ventilé, loin des fenêtres et portes.
\item Ne jamais faire le plein d'essence moteur en marche ou moteur chaud, et éloigner le carburant de toute flamme.
\end{itemize}
\textbf{Mesure d'entretien :} effectuer la vidange de l'huile et nettoyer le filtre à air régulièrement ; un moteur bien entretenu brûle mieux le carburant et rejette moins de gaz toxiques.
\textbf{Conclusion :} les symptômes correspondent à une intoxication au \ce{CO} ; le groupe doit fonctionner en plein air et être entretenu régulièrement.
\end{exoresolu}

\begin{attention}[Les 5 erreurs qui coûtent des points au BEPC]
\begin{enumerate}
\item \textbf{Confondre combustion interne et externe} : dire qu'une machine à vapeur est un moteur à combustion interne. Rédige toujours : « le combustible brûle \emph{à l'intérieur} du cylindre ».
\item \textbf{Inverser l'ordre des 4 temps} ou oublier l'état des soupapes : admission (soupape d'admission ouverte, piston descend), compression (tout fermé, piston monte), détente (tout fermé, piston descend), échappement (soupape d'échappement ouverte, piston monte).
\item \textbf{Oublier le temps moteur} : seule la \textbf{détente} fournit du travail ; ne jamais écrire que la compression « fait avancer » le véhicule.
\item \textbf{Mettre une bougie dans un moteur diesel} : le diesel fonctionne par \textbf{allumage par compression} ; il possède des \textbf{injecteurs}, pas de bougies d'allumage (seules des bougies de préchauffage peuvent exister pour le démarrage à froid).
\item \textbf{Erreur de calcul sur le cycle} : 4 temps = 4 courses = \textbf{2 tours} de vilebrequin (et non 4 tours). Vérifie : 1 course = 1 demi-tour.
\end{enumerate}
\end{attention}

\newpage

\coursec{Exercices}

\courssub{Je vérifie mes connaissances}

\begin{exercice}[QCM : reconnaître un moteur à combustion interne]
Pour chaque question, recopie la (ou les) bonne(s) réponse(s).
\begin{enumerate}
\item Un moteur à combustion interne est une machine qui transforme :
\begin{enumerate}
\item l'énergie électrique en énergie mécanique ;
\item l'énergie chimique d'un carburant en énergie mécanique ;
\item l'énergie mécanique en énergie thermique.
\end{enumerate}
\item Dans un moteur à combustion interne, la combustion du carburant a lieu :
\begin{enumerate}
\item dans une chaudière extérieure au moteur ;
\item à l'intérieur même du moteur (dans le cylindre) ;
\item dans le réservoir de carburant.
\end{enumerate}
\item Parmi les machines suivantes, laquelle est un moteur à combustion \emph{externe} ?
\begin{enumerate}
\item le moteur d'une moto-taxi ;
\item la machine à vapeur d'une ancienne locomotive ;
\item le moteur diesel d'un camion.
\end{enumerate}
\item Le mélange air-essence est enflammé dans un moteur à essence par :
\begin{enumerate}
\item un injecteur ;
\item une bougie d'allumage ;
\item le vilebrequin.
\end{enumerate}
\end{enumerate}
\end{exercice}

\begin{exercice}[Vrai ou faux, à justifier]
Recopie chaque affirmation, indique si elle est \textbf{vraie} ou \textbf{fausse}, puis justifie ta réponse en une ou deux phrases.
\begin{enumerate}
\item « Dans un moteur à combustion interne, le carburant brûle à l'extérieur du moteur. »
\item « Le piston se déplace entre le point mort haut (PMH) et le point mort bas (PMB). »
\item « Un moteur diesel possède des bougies d'allumage pour enflammer le gazole. »
\item « Le vilebrequin transforme le mouvement de translation du piston en mouvement de rotation. »
\item « Un moteur à 2 temps effectue un cycle complet en quatre courses du piston. »
\end{enumerate}
\end{exercice}

\begin{exercice}[Définitions à compléter]
Recopie et complète les phrases suivantes avec les mots ou expressions de la liste : \emph{carburant ; cylindrée ; combustion interne ; point mort haut ; soupape ; bielle}.
\begin{enumerate}
\item On appelle moteur à \ldots\ldots\ldots{} un moteur dans lequel le carburant brûle à l'intérieur du cylindre.
\item Le \ldots\ldots\ldots{} est la position la plus haute atteinte par le piston dans le cylindre.
\item La \ldots\ldots\ldots{} relie le piston au vilebrequin.
\item La \ldots\ldots\ldots{} est le volume balayé par le piston entre le PMH et le PMB, multiplié par le nombre de cylindres.
\item L'essence et le gazole sont des \ldots\ldots\ldots{} dérivés du pétrole.
\end{enumerate}
\end{exercice}

\begin{exercice}[Calcul direct : cylindrée unitaire]
Le moteur d'une moto possède un seul cylindre. Le volume balayé par le piston entre le point mort haut et le point mort bas est $V = \SI{125}{cm^3}$.
\begin{enumerate}
\item Rappelle la définition de la cylindrée d'un moteur.
\item Calcule la cylindrée totale de ce moteur.
\item Un moteur de voiture possède 4 cylindres identiques de volume unitaire $\SI{400}{cm^3}$ chacun. Calcule sa cylindrée totale, en \si{cm^3} puis en litres.
\end{enumerate}
\end{exercice}

\courssub{Je m'entraîne}

\begin{exercice}[Schéma à légender : coupe d'un cylindre]
On donne ci-dessous la coupe schématique d'un cylindre de moteur à essence à 4 temps.

\begin{popfigure}
\begin{tikzpicture}[scale=1.0]
% Cylindre
\draw[thick, popdark] (0,0) rectangle (3,5.2);
\draw[thick, popdark] (-0.4,5.2) -- (3.4,5.2);
% Piston
\draw[fill=popblueL, thick, popdark] (0.15,1.6) rectangle (2.85,2.6);
\node[popdark] at (1.5,2.1) {\small 1};
% Segment repère piston
\draw[fill=popdark] (1.5,1.6) circle (0.09);
% Bielle
\draw[very thick, poporange] (1.5,1.6) -- (2.1,0.2);
\node[popdark] at (2.55,0.9) {\small 2};
% Vilebrequin
\draw[fill=popgoldL, thick, popdark] (2.1,-0.35) circle (0.55);
\draw[fill=popdark] (2.1,-0.35) circle (0.07);
\node[popdark] at (3.15,-0.35) {\small 3};
% Soupapes
\draw[very thick, popgreen] (0.8,5.2) -- (0.8,6.1);
\draw[very thick, popgreen] (0.5,5.2) -- (1.1,5.2);
\node[popdark] at (0.8,6.45) {\small 4};
\draw[very thick, popgreen] (2.2,5.2) -- (2.2,6.1);
\draw[very thick, popgreen] (1.9,5.2) -- (2.5,5.2);
\node[popdark] at (2.2,6.45) {\small 5};
% Bougie
\draw[very thick, poppink] (1.5,5.2) -- (1.5,6.0);
\draw[fill=poppink] (1.5,6.0) circle (0.12);
\node[popdark] at (1.5,6.45) {\small 6};
% PMH et PMB
\draw[dashed, popmuted] (3.0,4.4) -- (4.3,4.4) node[right, popdark]{\small 7};
\draw[dashed, popmuted] (3.0,1.6) -- (4.3,1.6) node[right, popdark]{\small 8};
\end{tikzpicture}
\legende{Coupe schématique d'un cylindre de moteur à essence. Les numéros 1 à 8 désignent des organes ou des positions à identifier.}
\end{popfigure}

Recopie et complète le tableau suivant en identifiant chaque numéro parmi : \emph{piston ; bielle ; vilebrequin ; soupape d'admission ; soupape d'échappement ; bougie ; point mort haut (PMH) ; point mort bas (PMB)}.

\begin{center}
\begin{tabularx}{\linewidth}{|c|X|c|X|}
\hline
\rowcolor{popblueL} Numéro & Nom de l'organe ou de la position & Numéro & Nom de l'organe ou de la position \\
\hline
1 & & 5 & \\
\hline
2 & & 6 & \\
\hline
3 & & 7 & \\
\hline
4 & & 8 & \\
\hline
\end{tabularx}
\end{center}
\end{exercice}

\begin{exercice}[Remettre les 4 temps dans l'ordre]
On a photographié un cylindre de moteur à essence à quatre instants différents du cycle. Les descriptions, données dans le désordre, sont les suivantes :
\begin{itemize}
\item \textbf{Situation A :} le piston descend du PMH vers le PMB ; la soupape d'admission est ouverte, le mélange air-essence entre dans le cylindre.
\item \textbf{Situation B :} le piston monte ; la soupape d'échappement est ouverte ; les gaz brûlés sont chassés vers l'extérieur.
\item \textbf{Situation C :} les deux soupapes sont fermées ; une étincelle jaillit à la bougie ; le piston est violemment repoussé vers le PMB.
\item \textbf{Situation D :} les deux soupapes sont fermées ; le piston monte du PMB vers le PMH et comprime le mélange.
\end{itemize}
\begin{enumerate}
\item Nomme le temps moteur correspondant à chaque situation A, B, C et D.
\item Remets les quatre situations dans l'ordre chronologique du cycle à 4 temps.
\item Justifie ton classement de la situation C en précisant pourquoi ce temps est appelé « temps moteur ».
\end{enumerate}
\end{exercice}

\begin{exercice}[Fiche technique d'un véhicule : classer le moteur]
Voici la fiche technique simplifiée du véhicule d'un commerçant de Douala :

\begin{center}
\begin{tabularx}{\linewidth}{|l|X|}
\hline
\rowcolor{popblueL} Caractéristique & Valeur \\
\hline
Carburant utilisé & Gazole (gasoil) \\
\hline
Allumage du mélange & Par compression (pas de bougie) \\
\hline
Nombre de temps du cycle & 4 temps \\
\hline
Refroidissement & À eau (radiateur) \\
\hline
Nombre de cylindres & 4 \\
\hline
\end{tabularx}
\end{center}

\begin{enumerate}
\item Rappelle les quatre critères de classement des moteurs à combustion interne étudiés en cours.
\item À partir de la fiche technique, classe ce moteur selon chacun des quatre critères.
\item Ce moteur possède-t-il des bougies d'allumage ou des injecteurs ? Justifie.
\end{enumerate}
\end{exercice}

\begin{exercice}[Sécurité : le groupe électrogène dans la maison]
Pendant une coupure d'ENEOS, M. Etoundi installe son groupe électrogène à essence dans le couloir fermé de la maison pour le protéger de la pluie. Au bout d'une heure, les occupants ressentent des maux de tête et des nausées.
\begin{enumerate}
\item Nomme le gaz toxique, invisible et inodore, dégagé par les gaz d'échappement du moteur.
\item Explique, en t'aidant de l'équation de combustion incomplète, pourquoi ce gaz se forme dans une pièce mal aérée.
\item Explique pourquoi il est dangereux de faire fonctionner un groupe électrogène dans une pièce fermée.
\item Propose deux mesures de prévention concrètes que M. Etoundi aurait dû prendre.
\end{enumerate}
\end{exercice}

\courssub{Je me prépare au Bac}

\begin{exercice}[Type BEPC : la moto-taxi de Bertoua]
Un « benskin » (moto-taxi) de Bertoua est équipé d'un moteur à essence monocylindre (un seul cylindre) fonctionnant suivant le cycle à 4 temps. Le volume balayé par le piston entre le PMH et le PMB est $V = \SI{150}{cm^3}$. Le moteur effectue \num{3000} cycles par minute en régime normal.
\begin{enumerate}
\item Définis un moteur à combustion interne et explique en quoi le moteur de cette moto en est un, contrairement à une machine à vapeur.
\item Calcule la cylindrée de ce moteur.
\item Nomme dans l'ordre les quatre temps du cycle de ce moteur, en précisant pour chacun l'état (ouvert ou fermé) des soupapes et le sens de déplacement du piston.
\item Le conducteur dit : « À chaque tour du vilebrequin, mon moteur produit une explosion. » A-t-il raison ? Justifie en précisant le nombre de tours du vilebrequin par cycle complet.
\item Calcule le nombre de tours du vilebrequin effectués en une minute, sachant qu'un cycle complet correspond à 2 tours de vilebrequin.
\item Cite deux gestes d'entretien que le conducteur doit effectuer régulièrement pour préserver son moteur.
\end{enumerate}
\end{exercice}

\begin{exercice}[Type BEPC : essence ou diesel à la SONARA]
À Limbé, la SONARA transforme le pétrole brut en plusieurs carburants, dont l'essence et le gazole. Deux camions quittent le dépôt : le camion A est équipé d'un moteur à essence, le camion B d'un moteur diesel, tous deux à 4 temps.
\begin{enumerate}
\item Recopie et complète le tableau comparatif suivant :

\begin{center}
\begin{tabularx}{\linewidth}{|l|X|X|}
\hline
\rowcolor{popblueL} Critère & Moteur à essence (camion A) & Moteur diesel (camion B) \\
\hline
Carburant & & \\
\hline
Mode d'allumage du mélange & & \\
\hline
Organe spécifique d'allumage ou d'injection & & \\
\hline
Un avantage & & \\
\hline
\end{tabularx}
\end{center}

\item Pour le camion B, explique pourquoi le mélange air-gazole s'enflamme sans étincelle : que subit l'air aspiré pendant le temps de compression ?
\item Écris l'équation-bilan de la combustion complète de l'octane, principal constituant de l'essence : \ce{2C8H18 + 25O2 -> ...}. Complète les produits formés.
\item Lorsque l'arrivée d'air est insuffisante (filtre à air encrassé), la combustion devient incomplète. Nomme les deux produits dangereux ou polluants alors formés et précise un danger pour la santé.
\item Les deux camions roulent la nuit, moteur chaud. Cite deux règles de sécurité à respecter lors du remplissage des réservoirs à la station-service.
\end{enumerate}
\end{exercice}

\begin{exercice}[Type BEPC : synthèse — le groupe électrogène du quartier]
Pour pallier les délestages, un cybercafé de Yaoundé utilise un groupe électrogène dont le moteur est un moteur à essence à 4 temps, refroidi à air, de cylindrée totale $\SI{420}{cm^3}$ répartie sur 2 cylindres identiques.
\begin{enumerate}
\item Calcule la cylindrée unitaire (volume d'un cylindre) de ce moteur.
\item Schématise, par quatre petits dessins légendés, les quatre temps du cycle d'un cylindre, en indiquant à chaque fois la position approximative du piston et l'état des soupapes.
\item Classe ce moteur selon les quatre critères suivants : nature du carburant, mode d'allumage, nombre de temps, mode de refroidissement.
\item Le propriétaire fait fonctionner le groupe dans la salle des serveurs, portes et fenêtres fermées, « pour éviter le vol ». À l'aide de tes connaissances sur les gaz d'échappement, rédige un court texte (5 à 8 lignes) pour le convaincre de changer d'installation : tu nommeras le gaz dangereux, expliqueras son mode d'action sur l'organisme et proposeras deux mesures de prévention.
\item Les gaz d'échappement des moteurs contribuent à la pollution de l'air des grandes villes camerounaises. Cite deux autres mesures, à l'échelle d'un quartier ou d'une ville, permettant de réduire cette pollution.
\end{enumerate}
\end{exercice}

\newpage

\coursec{Corrigés des exercices}

\begin{corrige}[Exercice 1 — QCM : reconnaître un moteur à combustion interne]
\begin{enumerate}
\item \textbf{Bonne réponse : b.} Un moteur à combustion interne transforme l'\cle{énergie chimique} contenue dans le carburant (essence, gazole) en \cle{énergie mécanique} (mouvement de rotation du vilebrequin). La réponse a décrit un moteur électrique ; la réponse c est l'inverse du rôle d'un moteur.
\item \textbf{Bonne réponse : b.} La combustion a lieu \emph{à l'intérieur même du moteur}, dans le cylindre (plus précisément dans la chambre de combustion). C'est précisément ce qui justifie le nom « combustion \emph{interne} ».
\item \textbf{Bonne réponse : b.} Dans la machine à vapeur, le combustible (charbon, bois) brûle dans une chaudière \emph{extérieure} au moteur : c'est un moteur à combustion \emph{externe}. Les moteurs de moto et de camion sont à combustion interne.
\item \textbf{Bonne réponse : b.} Dans un moteur à essence, le mélange air-essence est enflammé par l'étincelle de la \cle{bougie d'allumage}. L'injecteur ne fait que pulvériser le carburant ; le vilebrequin transforme le mouvement.
\end{enumerate}
\end{corrige}

\begin{corrige}[Exercice 2 — Vrai ou faux, à justifier]
\begin{enumerate}
\item \textbf{Faux.} Par définition, dans un moteur à combustion \emph{interne}, le carburant brûle à l'\emph{intérieur} du moteur, dans le cylindre. C'est dans un moteur à combustion \emph{externe} (machine à vapeur) que le combustible brûle à l'extérieur.
\item \textbf{Vrai.} Le piston effectue un mouvement de va-et-vient (translation) entre deux positions extrêmes : le \cle{point mort haut} (PMH), position la plus haute, et le \cle{point mort bas} (PMB), position la plus basse.
\item \textbf{Faux.} Dans un moteur diesel, il n'y a \emph{pas} de bougie d'allumage : l'air fortement comprimé s'échauffe suffisamment pour que le gazole injecté s'enflamme \emph{spontanément} (allumage par compression). Le moteur diesel possède des \emph{injecteurs}.
\item \textbf{Vrai.} Le système \cle{bielle}-\cle{vilebrequin} transforme le mouvement de translation (va-et-vient) du piston en mouvement de rotation du vilebrequin, lequel entraîne les roues du véhicule.
\item \textbf{Faux.} Un moteur à \cle{2 temps} effectue un cycle complet en \emph{deux} courses du piston (soit un seul tour de vilebrequin). C'est le moteur à 4 temps qui nécessite quatre courses (deux tours de vilebrequin).
\end{enumerate}
\end{corrige}

\begin{corrige}[Exercice 3 — Définitions à compléter]
\begin{enumerate}
\item On appelle moteur à \textbf{combustion interne} un moteur dans lequel le carburant brûle à l'intérieur du cylindre.
\item Le \textbf{point mort haut} est la position la plus haute atteinte par le piston dans le cylindre.
\item La \textbf{bielle} relie le piston au vilebrequin.
\item La \textbf{cylindrée} est le volume balayé par le piston entre le PMH et le PMB, multiplié par le nombre de cylindres.
\item L'essence et le gazole sont des \textbf{carburants} dérivés du pétrole.
\end{enumerate}
\textbf{Remarque :} le mot \emph{soupape} n'était pas utilisé ici ; il désigne l'organe qui ouvre ou ferme les conduits d'admission et d'échappement.
\end{corrige}

\begin{corrige}[Exercice 4 — Calcul direct : cylindrée unitaire]
\begin{enumerate}
\item La \cle{cylindrée} d'un moteur est le volume total balayé par l'ensemble des pistons entre le point mort haut et le point mort bas. Elle est égale au volume unitaire $V_u$ (volume balayé dans un cylindre) multiplié par le nombre $n$ de cylindres :
\[ C = n \times V_u \]
\item Le moteur de la moto possède un seul cylindre ($n = 1$) de volume unitaire $V_u = \SI{125}{cm^3}$ :
\[ C = 1 \times 125 = \SI{125}{cm^3} \]
La cylindrée totale de ce moteur est \SI{125}{cm^3} (on dit couramment « une 125 »).
\item Pour le moteur de voiture : $n = 4$ cylindres, $V_u = \SI{400}{cm^3}$ :
\[ C = 4 \times 400 = \SI{1600}{cm^3} \]
Conversion en litres : $1~\text{L} = \SI{1000}{cm^3}$, donc :
\[ C = \dfrac{1600}{1000} = \SI{1,6}{L} \]
La cylindrée totale est \SI{1600}{cm^3}, soit \SI{1,6}{L} (on parle d'un moteur « 1,6 L »).
\end{enumerate}
\textbf{Point de vigilance :} ne pas confondre le volume unitaire (un seul cylindre) et la cylindrée totale (tous les cylindres). Vérifier la conversion : $\SI{1}{L} = \SI{1}{dm^3} = \SI{1000}{cm^3}$.
\end{corrige}

\begin{corrige}[Exercice 5 — Schéma à légender : coupe d'un cylindre]
\begin{center}
\begin{tabularx}{\linewidth}{|c|X|c|X|}
\hline
\rowcolor{popblueL} Numéro & Nom de l'organe ou de la position & Numéro & Nom de l'organe ou de la position \\
\hline
1 & Piston & 5 & Soupape d'échappement \\
\hline
2 & Bielle & 6 & Bougie (d'allumage) \\
\hline
3 & Vilebrequin & 7 & Point mort haut (PMH) \\
\hline
4 & Soupape d'admission & 8 & Point mort bas (PMB) \\
\hline
\end{tabularx}
\end{center}
\textbf{Justifications :}
\begin{itemize}
\item Le repère 1 désigne la pièce coulissante dans le cylindre : le piston ; 2 est la tige articulée qui le relie au maneton : la bielle ; 3 est l'arbre à manivelle : le vilebrequin.
\item 4 et 5 sont les deux soupapes situées à la partie supérieure (culasse) : l'une pour l'admission du mélange frais, l'autre pour l'échappement des gaz brûlés. 6, placée entre les deux soupapes au sommet de la chambre de combustion, est la bougie.
\item 7 est le niveau le plus haut atteint par le piston (PMH) ; 8 le niveau le plus bas (PMB). Le volume compris entre ces deux niveaux, balayé par le piston, est la cylindrée unitaire.
\end{itemize}
\end{corrige}

\begin{corrige}[Exercice 6 — Remettre les 4 temps dans l'ordre]
\begin{enumerate}
\item Identification de chaque situation :
\begin{itemize}
\item \textbf{Situation A :} \cle{admission} (1\textsuperscript{er} temps) — soupape d'admission ouverte, piston descendant.
\item \textbf{Situation B :} \cle{échappement} (4\textsuperscript{e} temps) — soupape d'échappement ouverte, piston montant.
\item \textbf{Situation C :} \cle{combustion-détente} (3\textsuperscript{e} temps, temps moteur) — soupapes fermées, étincelle, piston repoussé vers le bas.
\item \textbf{Situation D :} \cle{compression} (2\textsuperscript{e} temps) — soupapes fermées, piston montant.
\end{itemize}
\item Ordre chronologique du cycle : \textbf{A} (admission) $\rightarrow$ \textbf{D} (compression) $\rightarrow$ \textbf{C} (combustion-détente) $\rightarrow$ \textbf{B} (échappement).
\item La situation C est appelée « \cle{temps moteur} » parce que c'est le \emph{seul} temps qui \emph{produit} de l'énergie mécanique : l'étincelle de la bougie enflamme le mélange comprimé, les gaz brûlés se dilatent violemment et \emph{poussent} le piston vers le PMB, ce qui fait tourner le vilebrequin. Les trois autres temps \emph{consomment} de l'énergie (fournie par l'inertie du volant moteur) : ils ne font que préparer ce temps moteur.
\end{enumerate}
\end{corrige}

\begin{corrige}[Exercice 7 — Fiche technique d'un véhicule : classer le moteur]
\begin{enumerate}
\item Les quatre critères de classement étudiés sont :
\begin{itemize}
\item la \cle{nature du carburant} (essence ou gazole) ;
\item le \cle{mode d'allumage} du mélange (allumage commandé par bougie ou allumage par compression) ;
\item le \cle{nombre de temps} du cycle (2 temps ou 4 temps) ;
\item le \cle{mode de refroidissement} (à air ou à eau/liquide).
\end{itemize}
(On peut aussi citer le nombre et la disposition des cylindres.)
\item Classement du moteur du commerçant :
\begin{itemize}
\item nature du carburant : moteur à \textbf{gazole} ;
\item mode d'allumage : allumage \textbf{par compression} (auto-inflammation) ;
\item nombre de temps : moteur à \textbf{4 temps} ;
\item refroidissement : \textbf{à eau} (par liquide de refroidissement et radiateur).
\end{itemize}
C'est donc un \textbf{moteur diesel à 4 temps, refroidi à eau}.
\item Ce moteur possède des \textbf{injecteurs} et \emph{pas} de bougies d'allumage. Justification : l'allumage se fait par compression : l'air aspiré est fortement comprimé, sa température s'élève au-dessus de la température d'auto-inflammation du gazole ; le gazole est alors pulvérisé finement par l'injecteur et s'enflamme spontanément au contact de l'air chaud. La bougie d'allumage est donc inutile.
\end{enumerate}
\end{corrige}

\begin{corrige}[Exercice 8 — Sécurité : le groupe électrogène dans la maison]
\begin{enumerate}
\item Le gaz toxique, invisible et inodore, est le \cle{monoxyde de carbone}, de formule \ce{CO}.
\item Dans une pièce mal aérée, l'apport de dioxygène \ce{O2} est insuffisant : la combustion de l'essence devient \emph{incomplète}. Au lieu de ne produire que du dioxyde de carbone et de l'eau, elle produit du monoxyde de carbone (et du carbone noir). Par exemple, pour l'octane :
\[ \ce{2C8H18 + 17O2 -> 16CO + 18H2O} \]
Le manque de \ce{O2} empêche l'oxydation complète du carbone en \ce{CO2}.
\item Le monoxyde de carbone est dangereux parce qu'il est \emph{indétectable} (sans couleur, sans odeur) et qu'inhalé, il se fixe sur l'hémoglobine du sang à la place du dioxygène : le sang ne transporte plus correctement le dioxygène vers les organes. Il en résulte des maux de tête, des nausées, puis des vertiges, une perte de conscience, et la mort par asphyxie en cas d'exposition prolongée — d'autant plus vite que la pièce est fermée et que le gaz s'y accumule.
\item Mesures de prévention :
\begin{itemize}
\item installer le groupe électrogène \emph{à l'extérieur} de la maison, dans un endroit dégagé et aéré, à distance des portes et fenêtres (on peut le protéger de la pluie par un auvent ouvert) ;
\item ne jamais le faire fonctionner dans une pièce fermée, et installer si possible un \emph{détecteur de monoxyde de carbone} dans l'habitation.
\end{itemize}
\end{enumerate}
\end{corrige}

\begin{corrige}[Exercice 9 — Type BEPC : la moto-taxi de Bertoua]
\begin{enumerate}
\item Un \cle{moteur à combustion interne} est une machine thermique qui transforme l'énergie chimique d'un carburant en énergie mécanique, la combustion du carburant ayant lieu \emph{à l'intérieur même du moteur}, dans le cylindre. Le moteur de la moto en est un : l'essence brûle directement dans son cylindre et la détente des gaz pousse le piston. Au contraire, dans une machine à vapeur, le combustible brûle dans une chaudière \emph{extérieure} qui chauffe de l'eau : c'est un moteur à combustion \emph{externe}.
\item Le moteur est monocylindre ($n = 1$) et le volume unitaire est $V_u = \SI{150}{cm^3}$ :
\[ C = n \times V_u = 1 \times 150 = \SI{150}{cm^3} \]
La cylindrée du moteur est \SI{150}{cm^3}.
\item Les quatre temps du cycle, dans l'ordre :
\begin{itemize}
\item \textbf{1\textsuperscript{er} temps — Admission :} soupape d'admission \emph{ouverte}, soupape d'échappement \emph{fermée} ; le piston \emph{descend} du PMH vers le PMB ; le mélange air-essence entre.
\item \textbf{2\textsuperscript{e} temps — Compression :} les deux soupapes sont \emph{fermées} ; le piston \emph{monte} du PMB vers le PMH et comprime le mélange.
\item \textbf{3\textsuperscript{e} temps — Combustion-détente (temps moteur) :} les deux soupapes sont \emph{fermées} ; l'étincelle de la bougie enflamme le mélange ; le piston est \emph{repoussé} du PMH vers le PMB.
\item \textbf{4\textsuperscript{e} temps — Échappement :} soupape d'échappement \emph{ouverte}, soupape d'admission \emph{fermée} ; le piston \emph{monte} du PMB vers le PMH et chasse les gaz brûlés.
\end{itemize}
\item Le conducteur a \textbf{tort}. Un cycle complet à 4 temps comporte 4 courses du piston, ce qui correspond à \textbf{2 tours de vilebrequin}. Il n'y a qu'\emph{une seule} explosion (temps moteur) par cycle, donc une explosion tous les \emph{deux} tours de vilebrequin, et non à chaque tour.
\item En une minute, le moteur effectue \num{3000} cycles. Or 1 cycle $=$ 2 tours de vilebrequin :
\[ N = 3000 \times 2 = \num{6000} \text{ tours} \]
Le vilebrequin effectue \num{6000} tours par minute (régime de \num{6000} tr/min).
\item Deux gestes d'entretien réguliers : vérifier et changer l'\emph{huile moteur} (vidange) ; nettoyer ou remplacer le \emph{filtre à air} ; on peut aussi citer le remplacement de la bougie usée ou le contrôle du niveau de carburant et l'utilisation d'une essence propre.
\end{enumerate}
\textbf{Barème indicatif (sur 10) :} question 1 : 2 pts ; question 2 : 1 pt ; question 3 : 3 pts (0,75 pt par temps : nom + soupapes + sens du piston) ; question 4 : 1,5 pt ; question 5 : 1,5 pt ; question 6 : 1 pt.

\textbf{Points de vigilance :} exiger la mention « combustion \emph{à l'intérieur} du cylindre » dans la définition ; pour la question 4, la justification doit citer explicitement « 2 tours de vilebrequin par cycle » ; ne pas confondre nombre de cycles par minute et régime moteur en tr/min.
\end{corrige}

\begin{corrige}[Exercice 10 — Type BEPC : essence ou diesel à la SONARA]
\begin{enumerate}
\item Tableau comparatif complété :
\begin{center}
\begin{tabularx}{\linewidth}{|l|X|X|}
\hline
\rowcolor{popblueL} Critère & Moteur à essence (camion A) & Moteur diesel (camion B) \\
\hline
Carburant & Essence & Gazole (gasoil) \\
\hline
Mode d'allumage du mélange & Allumage commandé (étincelle) & Allumage par compression (auto-inflammation) \\
\hline
Organe spécifique d'allumage ou d'injection & Bougie d'allumage & Injecteur \\
\hline
Un avantage & Fonctionnement plus souple et moins bruyant & Consommation plus faible, couple plus important (adapté aux charges lourdes) \\
\hline
\end{tabularx}
\end{center}
\item Pendant le temps de compression, l'air aspiré seul (sans carburant) est fortement comprimé par le piston qui monte, les soupapes étant fermées. Or comprimer un gaz l'\emph{échauffe} : la température de l'air dépasse alors la température d'auto-inflammation du gazole. Lorsque l'injecteur pulvérise le gazole finement dans cet air très chaud, le mélange s'enflamme \emph{spontanément}, sans étincelle.
\item Combustion complète de l'octane :
\[ \ce{2C8H18 + 25O2 -> 16CO2 + 18H2O} \]
Les produits formés sont le \cle{dioxyde de carbone} \ce{CO2} et l'\cle{eau} \ce{H2O}. Vérification de l'équilibre : C : $2\times 8 = 16$ ; H : $2\times 18 = 36 = 18\times 2$ ; O : $25\times 2 = 50 = 16\times 2 + 18$.
\item En combustion incomplète, il se forme du \cle{monoxyde de carbone} \ce{CO} et du \cle{carbone} \ce{C} (suies, fumées noires). Danger pour la santé : le monoxyde de carbone est un gaz toxique, invisible et inodore, qui se fixe sur l'hémoglobine du sang à la place du dioxygène et provoque l'asphyxie (maux de tête, nausées, pouvant aller jusqu'à la mort).
\item Deux règles de sécurité à la station-service :
\begin{itemize}
\item \emph{arrêter le moteur} pendant le remplissage et ne pas fumer ni approcher de flamme ou d'étincelle (les vapeurs d'essence sont inflammables) ;
\item ne pas remplir le réservoir à ras bord (éviter les débordements sur un moteur chaud) et ne pas utiliser de téléphone portable à proximité du pistolet.
\end{itemize}
\end{enumerate}
\textbf{Barème indicatif (sur 10) :} question 1 : 3 pts ; question 2 : 2 pts ; question 3 : 2 pts ; question 4 : 1,5 pt ; question 5 : 1,5 pt.

\textbf{Points de vigilance :} dans la question 2, préciser que dans un diesel \emph{l'air seul} est comprimé (le gazole n'arrive qu'à la fin, par l'injecteur) — c'est la différence essentielle avec le moteur à essence où le \emph{mélange} est comprimé ; l'équation-bilan doit être équilibrée (compter les atomes) ; citer \ce{CO} \emph{et} \ce{C} pour la combustion incomplète.
\end{corrige}

\begin{corrige}[Exercice 11 — Type BEPC : synthèse — le groupe électrogène du quartier]
\begin{enumerate}
\item La cylindrée totale est $C = \SI{420}{cm^3}$ répartie sur $n = 2$ cylindres identiques. La cylindrée unitaire est :
\[ V_u = \dfrac{C}{n} = \dfrac{420}{2} = \SI{210}{cm^3} \]
Chaque cylindre a un volume unitaire de \SI{210}{cm^3}.
\item Schéma des quatre temps (à reproduire sur la copie avec quatre petits cylindres) :

\begin{popfigure}
\begin{tikzpicture}[scale=0.62]
\begin{scope}[xshift=0cm]
\draw[thick, popdark] (0,0) rectangle (2,3.6);
\draw[fill=popblueL, thick, popdark] (0.1,0.9) rectangle (1.9,1.6);
\draw[->, very thick, popblue] (1,2.6) -- (1,1.9);
\draw[very thick, popgreen] (0.5,3.6) -- (0.5,4.3); \draw[very thick, popgreen] (0.3,3.75) -- (0.7,3.75);
\draw[very thick, popmuted] (1.5,3.6) -- (1.5,4.3); \draw[very thick, popmuted] (1.3,3.6) -- (1.7,3.6);
\draw[->, thick, poporange] (0.5,4.6) -- (0.5,4.0);
\node[popdark] at (1,-0.5) {\small 1. Admission};
\node[popdark] at (1,-1.1) {\small piston descend};
\end{scope}
\begin{scope}[xshift=3.6cm]
\draw[thick, popdark] (0,0) rectangle (2,3.6);
\draw[fill=popblueL, thick, popdark] (0.1,1.6) rectangle (1.9,2.3);
\draw[->, very thick, popblue] (1,0.9) -- (1,1.5);
\draw[very thick, popmuted] (0.5,3.6) -- (0.5,4.3); \draw[very thick, popmuted] (0.3,3.6) -- (0.7,3.6);
\draw[very thick, popmuted] (1.5,3.6) -- (1.5,4.3); \draw[very thick, popmuted] (1.3,3.6) -- (1.7,3.6);
\node[popdark] at (1,-0.5) {\small 2. Compression};
\node[popdark] at (1,-1.1) {\small piston monte};
\end{scope}
\begin{scope}[xshift=7.2cm]
\draw[thick, popdark] (0,0) rectangle (2,3.6);
\draw[fill=popblueL, thick, popdark] (0.1,2.6) rectangle (1.9,3.3);
\draw[->, very thick, popblue] (1,2.4) -- (1,1.6);
\draw[very thick, popmuted] (0.5,3.6) -- (0.5,4.3); \draw[very thick, popmuted] (0.3,3.6) -- (0.7,3.6);
\draw[very thick, popmuted] (1.5,3.6) -- (1.5,4.3); \draw[very thick, popmuted] (1.3,3.6) -- (1.7,3.6);
\draw[very thick, poppink] (1,3.6) -- (1,4.2);
\node[poppink] at (1,3.55) {\tiny étincelle};
\node[popdark] at (1,-0.5) {\small 3. Détente};
\node[popdark] at (1,-1.1) {\small piston repoussé};
\end{scope}
\begin{scope}[xshift=10.8cm]
\draw[thick, popdark] (0,0) rectangle (2,3.6);
\draw[fill=popblueL, thick, popdark] (0.1,1.6) rectangle (1.9,2.3);
\draw[->, very thick, popblue] (1,0.9) -- (1,1.5);
\draw[very thick, popmuted] (0.5,3.6) -- (0.5,4.3); \draw[very thick, popmuted] (0.3,3.6) -- (0.7,3.6);
\draw[very thick, popgreen] (1.5,3.6) -- (1.5,4.3); \draw[very thick, popgreen] (1.3,3.75) -- (1.7,3.75);
\draw[->, thick, poporange] (1.5,4.0) -- (1.5,4.7);
\node[popdark] at (1,-0.5) {\small 4. Échappement};
\node[popdark] at (1,-1.1) {\small piston monte};
\end{scope}
\end{tikzpicture}
\legende{Les quatre temps du cycle : 1. admission (soupape d'admission ouverte, piston descend) ; 2. compression (soupapes fermées, piston monte) ; 3. combustion-détente (soupapes fermées, étincelle, piston repoussé) ; 4. échappement (soupape d'échappement ouverte, piston monte).}
\end{popfigure}

\item Classement du moteur selon les quatre critères :
\begin{itemize}
\item nature du carburant : moteur à \textbf{essence} ;
\item mode d'allumage : \textbf{allumage commandé} par bougie (étincelle) ;
\item nombre de temps : moteur à \textbf{4 temps} ;
\item mode de refroidissement : refroidissement \textbf{à air}.
\end{itemize}
\item Texte de sensibilisation (exemple de réponse attendue) :

« Monsieur, faire fonctionner votre groupe électrogène dans une pièce fermée met votre vie en danger. Les gaz d'échappement du moteur contiennent du monoxyde de carbone (\ce{CO}), un gaz toxique invisible et sans odeur, formé par la combustion incomplète de l'essence lorsque l'air manque. Inhalé, ce gaz se fixe sur l'hémoglobine du sang à la place du dioxygène : les organes ne sont plus oxygénés, ce qui provoque maux de tête, nausées, perte de conscience, puis la mort par asphyxie. Dans une salle fermée, le gaz s'accumule rapidement. Je vous conseille donc d'installer le groupe à l'extérieur, dans un endroit dégagé et ventilé, loin des portes et fenêtres, et d'assurer une bonne aération permanente des locaux (ou d'installer un détecteur de monoxyde de carbone). »
\item Deux mesures à l'échelle d'un quartier ou d'une ville :
\begin{itemize}
\item développer les transports en commun et le covoiturage pour réduire le nombre de véhicules en circulation ;
\item faire contrôler régulièrement l'état des moteurs (visites techniques, interdiction des véhicules trop polluants) ; on peut aussi citer la création d'espaces verts, l'encadrement de l'importation des véhicules âgés ou la sensibilisation à l'entretien des moteurs.
\end{itemize}
\end{enumerate}
\textbf{Barème indicatif (sur 10) :} question 1 : 1 pt ; question 2 : 3 pts (schémas légendés avec position du piston et état des soupapes) ; question 3 : 2 pts (0,5 pt par critère) ; question 4 : 3 pts (gaz nommé 0,5 pt ; mode d'action 1 pt ; deux mesures 1,5 pt) ; question 5 : 1 pt.

\textbf{Points de vigilance :} pour la question 2, chaque dessin doit montrer la \emph{position du piston} ET l'\emph{état des soupapes} (ouvertes/fermées) : c'est ce qui est évalué ; dans le texte de la question 4, le jury attend explicitement le nom du gaz (\ce{CO}), son mode d'action (fixation sur l'hémoglobine) et deux mesures concrètes ; pour la question 1, bien diviser la cylindrée totale par le nombre de cylindres, et non l'inverse.
\end{corrige}

\newpage

\coursec{Fiche bilan}

\begin{popfigure}
\begin{center}\tcbox[colback=popdark,colframe=popdark,coltext=white,arc=9pt,boxsep=4pt,left=6pt,right=6pt]{\bfseries\small Moteurs à combustion interne}\end{center}
\vspace{-2pt}
\begin{tcbraster}[raster columns=3,raster equal height=rows,raster column skip=6pt,raster row skip=6pt,enhanced,blankest]
\begin{tcolorbox}[enhanced,colback=white,colframe=popblue,boxrule=0.9pt,arc=3pt,title={Définition},fonttitle=\bfseries\scriptsize,coltitle=white,colbacktitle=popblue,left=4pt,right=4pt,top=2pt,bottom=2pt,boxsep=1pt,toptitle=1.5pt,bottomtitle=1.5pt,halign=left]
\begin{itemize}[nosep,leftmargin=*,itemsep=1pt,font=\scriptsize]
\item Énergie chimique $\to$ mécanique
\item Combustion dans le cylindre
\end{itemize}
\end{tcolorbox}
\begin{tcolorbox}[enhanced,colback=white,colframe=poporange,boxrule=0.9pt,arc=3pt,title={Constitution},fonttitle=\bfseries\scriptsize,coltitle=white,colbacktitle=poporange,left=4pt,right=4pt,top=2pt,bottom=2pt,boxsep=1pt,toptitle=1.5pt,bottomtitle=1.5pt,halign=left]
\begin{itemize}[nosep,leftmargin=*,itemsep=1pt,font=\scriptsize]
\item Cylindre, piston, bielle
\item Vilebrequin, soupapes, bougie
\end{itemize}
\end{tcolorbox}
\begin{tcolorbox}[enhanced,colback=white,colframe=popgreen,boxrule=0.9pt,arc=3pt,title={Cycle 4 temps},fonttitle=\bfseries\scriptsize,coltitle=white,colbacktitle=popgreen,left=4pt,right=4pt,top=2pt,bottom=2pt,boxsep=1pt,toptitle=1.5pt,bottomtitle=1.5pt,halign=left]
\begin{itemize}[nosep,leftmargin=*,itemsep=1pt,font=\scriptsize]
\item Admission, compression
\item Explosion, échappement
\end{itemize}
\end{tcolorbox}
\begin{tcolorbox}[enhanced,colback=white,colframe=poppurple,boxrule=0.9pt,arc=3pt,title={2 temps et Diesel},fonttitle=\bfseries\scriptsize,coltitle=white,colbacktitle=poppurple,left=4pt,right=4pt,top=2pt,bottom=2pt,boxsep=1pt,toptitle=1.5pt,bottomtitle=1.5pt,halign=left]
\begin{itemize}[nosep,leftmargin=*,itemsep=1pt,font=\scriptsize]
\item 2 temps : 1 tour
\item Diesel : auto-allumage
\end{itemize}
\end{tcolorbox}
\begin{tcolorbox}[enhanced,colback=white,colframe=poppink,boxrule=0.9pt,arc=3pt,title={Classement},fonttitle=\bfseries\scriptsize,coltitle=white,colbacktitle=poppink,left=4pt,right=4pt,top=2pt,bottom=2pt,boxsep=1pt,toptitle=1.5pt,bottomtitle=1.5pt,halign=left]
\begin{itemize}[nosep,leftmargin=*,itemsep=1pt,font=\scriptsize]
\item Selon le carburant
\item Selon le cycle
\end{itemize}
\end{tcolorbox}
\begin{tcolorbox}[enhanced,colback=white,colframe=popgold,boxrule=0.9pt,arc=3pt,title={Sécurité et environnement},fonttitle=\bfseries\scriptsize,coltitle=white,colbacktitle=popgold,left=4pt,right=4pt,top=2pt,bottom=2pt,boxsep=1pt,toptitle=1.5pt,bottomtitle=1.5pt,halign=left]
\begin{itemize}[nosep,leftmargin=*,itemsep=1pt,font=\scriptsize]
\item Entretien, pollution
\end{itemize}
\end{tcolorbox}
\end{tcbraster}

\legende{Carte mentale de la leçon : les six idées clés sur les moteurs à combustion interne.}
\end{popfigure}

\begin{bilan}
\textbf{Définitions clés}
\begin{itemize}
\item Un \cle{moteur à combustion interne} est une machine qui transforme l'\cle{énergie chimique} d'un carburant en \cle{énergie mécanique} : la combustion du mélange air-carburant se fait \emph{à l'intérieur} du moteur (dans le cylindre).
\item Le \cle{cylindre} est le tube où coulisse le \cle{piston} ; la \cle{bielle} relie le piston au \cle{vilebrequin}, qui transforme le mouvement de translation du piston en mouvement de rotation.
\item La \cle{cylindrée} est le volume balayé par le piston entre le point mort haut (PMH) et le point mort bas (PMB) ; elle s'exprime en \si{cm^3} (ex. moto 125 : \SI{125}{cm^3}).
\item Un \cle{temps moteur} correspond à une course du piston (montée ou descente).
\end{itemize}
\textbf{Le cycle à 4 temps (moteur à essence)} — à connaître par cœur
\begin{center}
\begin{tabularx}{\linewidth}{|c|l|X|}\hline
\rowcolor{popblueL} \textbf{Temps} & \textbf{Nom} & \textbf{Ce qui se passe} \\ \hline
1 & Admission & Le piston descend, la soupape d'admission s'ouvre : le mélange air-essence entre. \\ \hline
2 & Compression & Soupapes fermées, le piston remonte et comprime le mélange. \\ \hline
3 & Explosion-détente & La bougie produit l'étincelle : la combustion pousse le piston (seul temps moteur). \\ \hline
4 & Échappement & La soupape d'échappement s'ouvre, le piston remonte et chasse les gaz brûlés. \\ \hline
\end{tabularx}
\end{center}
\textbf{Autres moteurs et classement}
\begin{itemize}
\item \cle{Moteur à 2 temps} : le cycle se fait en 2 courses (1 tour de vilebrequin) ; mélange essence-huile ; léger et simple (motos, tondeuses), mais pollue davantage.
\item \cle{Moteur Diesel} : pas de bougie ; l'air seul est fortement comprimé puis le gasoil injecté \cle{s'enflamme spontanément} (auto-allumage) ; robuste et économique (camions, groupes électrogènes).
\item Classement : selon le \cle{carburant} (essence, Diesel), selon le \cle{cycle} (2 ou 4 temps), selon l'\cle{allumage} (commandé par bougie ou par compression).
\end{itemize}
\textbf{Méthode express} : pour identifier un temps moteur sur un schéma, regarder (1) le sens de déplacement du piston, (2) l'état des soupapes (ouvertes/fermées), (3) l'étincelle de la bougie.
\end{bilan}

\begin{aretenir}[Checklist avant le BEPC]
\begin{itemize}
\item[$\square$] Je sais définir un moteur à combustion interne et citer la transformation d'énergie réalisée.
\item[$\square$] Je sais nommer les organes principaux : cylindre, piston, bielle, vilebrequin, soupapes, bougie.
\item[$\square$] Je sais citer dans l'ordre les 4 temps du moteur à essence.
\item[$\square$] Je sais décrire ce qui se passe à chaque temps (piston, soupapes, étincelle).
\item[$\square$] Je sais que seul le temps d'explosion-détente est moteur.
\item[$\square$] Je sais que le cycle 4 temps correspond à 2 tours de vilebrequin.
\item[$\square$] Je sais expliquer la différence entre moteur à 2 temps et moteur à 4 temps.
\item[$\square$] Je sais expliquer l'allumage d'un moteur Diesel (auto-allumage, sans bougie).
\item[$\square$] Je sais classer les moteurs selon le carburant, le cycle et le mode d'allumage.
\item[$\square$] Je sais citer deux règles de sécurité (pas d'essence près d'une flamme, moteur chaud) et un geste d'entretien (vidange, filtre à air).
\item[$\square$] Je sais citer un impact environnemental des moteurs (gaz polluants, \ce{CO2}) et une solution (entretien, catalyseur).
\end{itemize}
\end{aretenir}

\findecours

\end{document}
